% ============================================================
% Supplement to the concise conference-facing main manuscript (v12)
% Detailed methods/audit archive plus frozen post-closure descriptive diagnostic
% Development audits were executed; no primary or reserve seed was accessed.
% 29 August 2026
%
% BUILD MODES
%   Anonymous review build (DEFAULT):
%       pdflatex IC_TORS26_SHAP_MCDM_conference_supplement_2026-08-29_v5.tex
%   Identified camera-ready build:
%       pdflatex "\def\CAMERAREADY{1}% ============================================================
% Supplement to the concise conference-facing main manuscript (v12)
% Detailed methods/audit archive plus frozen post-closure descriptive diagnostic
% Development audits were executed; no primary or reserve seed was accessed.
% 29 August 2026
%
% BUILD MODES
%   Anonymous review build (DEFAULT):
%       pdflatex IC_TORS26_SHAP_MCDM_conference_supplement_2026-08-29_v5.tex
%   Identified camera-ready build:
%       pdflatex "\def\CAMERAREADY{1}% ============================================================
% Supplement to the concise conference-facing main manuscript (v12)
% Detailed methods/audit archive plus frozen post-closure descriptive diagnostic
% Development audits were executed; no primary or reserve seed was accessed.
% 29 August 2026
%
% BUILD MODES
%   Anonymous review build (DEFAULT):
%       pdflatex IC_TORS26_SHAP_MCDM_conference_supplement_2026-08-29_v5.tex
%   Identified camera-ready build:
%       pdflatex "\def\CAMERAREADY{1}% ============================================================
% Supplement to the concise conference-facing main manuscript (v12)
% Detailed methods/audit archive plus frozen post-closure descriptive diagnostic
% Development audits were executed; no primary or reserve seed was accessed.
% 29 August 2026
%
% BUILD MODES
%   Anonymous review build (DEFAULT):
%       pdflatex IC_TORS26_SHAP_MCDM_conference_supplement_2026-08-29_v5.tex
%   Identified camera-ready build:
%       pdflatex "\def\CAMERAREADY{1}\input{IC_TORS26_SHAP_MCDM_conference_supplement_2026-08-29_v5}"
% ============================================================
% ChkTeX: intentional en dashes, literal digest ellipses, label placement,
% acronym punctuation and LNCS name/bibliography conventions.
% chktex-file 8
% chktex-file 11
% chktex-file 13
% chktex-file 24
% chktex-file 36
% chktex-file 38

\documentclass[runningheads,orivec]{llncs}

\usepackage[T1]{fontenc}
\usepackage[utf8]{inputenc}
\usepackage{amsmath,amssymb,amsfonts}
\usepackage{booktabs}
\usepackage{graphicx}
\usepackage{tabularx}
\usepackage{hyperref}
\hypersetup{hidelinks}

\newcommand{\E}{\mathbb{E}}
\newcommand{\median}{\operatorname{median}}
\newcommand{\IQR}{\operatorname{IQR}}
\newcommand{\SD}{\operatorname{SD}}
\newcommand{\clip}{\operatorname{clip}}
\newcommand{\XGB}{XGBoost}

% ---- build-mode switch: anonymous by default -----------------
\providecommand{\CAMERAREADY}{0}
\newif\ifanon
\anontrue
\ifnum\CAMERAREADY=1 \anonfalse \fi
% --------------------------------------------------------------

\begin{document}
\emergencystretch=1.5em

\title{Supplementary Methods and Audit Archive for\texorpdfstring{\\}{ }
``From Predictive Importance to Decision Value: Evaluating XAI-Based Weights for
ITS Prioritisation''}
\titlerunning{Supplementary Benchmark-Adequacy Record}

\ifanon
  \author{Anonymous Author(s)}
  \authorrunning{Anonymous Author(s)}
  \institute{Anonymous Institution(s)}
\else
  \author{Mohamed Yassine Neifar\inst{1,2} \and
  Ahmed Frikha\inst{2} \and
  Nadir Farhi\inst{1}}
  \authorrunning{Neifar et al.}
  \institute{Universit\'e Gustave Eiffel, COSYS--GRETTIA,
  Champs-sur-Marne, F-77447 Marne-la-Vall\'ee Cedex 2, France
  \and
  Institut Sup\'erieur de Gestion Industrielle de Sfax (ISGIS),
  OLID, University of Sfax, Sfax 3038, Tunisia}
\fi

\maketitle

\section{Status and Interpretation Boundary}

This supplement reports executed development audits and preserves the detailed primary and
robustness analyses that had been specified before the prospectively frozen Pilot-B
benchmark-adequacy gate was executed. The latter analyses were conditional on all five Pilot-B
hard-stop flags being false. The modal-winner effective-tie stop fired, so the primary factorial
was not authorized. Primary seeds 11001--11030 and reserve seeds 30001--30005 were never
accessed; all sections explicitly labelled ``Archived'' remain unexecuted protocols rather than
results.

The closed v1 state is anchored by an immutable version-control reference.
\ifanon
Exact repository identifiers and complete digests are withheld from this anonymous reviewer
file; they are available confidentially to editors and will accompany the public
reproducibility archive after review.
\else
Exact repository identifiers and complete digests accompany the public reproducibility archive.
\fi
This document is therefore a methods, audit-results and archival-protocol supplement. It is not
authorization to execute the protected design.
After permanent Pilot-B closure, one separately frozen development-only descriptive diagnostic
reconstructed the context-level MOORA selection maps of the seven observed archived fixed
vectors. It used no primary/reserve world, generated no target outcome, fitted or queried no
model, re-estimated no weight and could not alter the gate. Its protocol, implementation,
independent audit and result were frozen as a distinct post-closure layer.
Methodological citations are provided in the main manuscript and are not duplicated here.

\subsection*{Question Structure and Evidential Status}

The main manuscript states two prospective research questions and one separate post-closure
descriptive question. RQ1 asks whether structural diagnostics can reveal that apparent
contextual variation disappears under the normalization used by the downstream MCDM operator.
RQ2 asks whether, after that structural defect is corrected, structured global-weight pipelines
provide sufficient incremental decision value over random-weight and constant-decision
references to justify protected primary evaluation. Both are addressed by the prospectively
structured benchmark-development record and the frozen Pilot-B result reported here.

DQ1 asks what the archived context-level selection maps reveal about the observed decision
behaviour of the seven frozen fixed-weight vectors. DQ1 is post-closure and descriptive. It is
non-confirmatory, it carries no pass/fail role, and it cannot modify, reinterpret or reopen the
Pilot-B gate. Section~\ref{supp:sec:fixedselection} reports it under that label.

Throughout this supplement, \emph{each replication seed defines one benchmark world}, and the
world is the statistical unit. Where a numbered identifier such as 21001 appears it names a
world; the phrase ``replication seed'' is retained only when the frozen artefact field or
protocol clause is itself named that way. \emph{Validity} denotes conceptual scientific
properties of the benchmark; \emph{adequacy} denotes the operational prospectively frozen gate
used to decide whether protected primary execution was warranted.

\subsection*{Reader Map}

The conference-facing main manuscript follows a problem--method--result--meaning structure and
deliberately omits most development chronology. Table~\ref{supp:tab:readermap} directs readers
to the corresponding detailed evidence. This separation changes presentation only: the frozen
scientific record, Pilot-B decision and post-closure interpretation boundary are unchanged.

\begin{table}[t]
\centering
\caption{Navigation from the concise main argument to the detailed reproducibility record,
with the evidential status of each block. Status codes are: EDA, executed development audit;
PBA, Pilot-B confirmatory adequacy result; PCD, post-closure descriptive diagnostic; ANE,
archived and not executed.}
\label{supp:tab:readermap}
\small
\begin{tabularx}{\linewidth}{@{}lXc@{}}
\toprule
\textbf{Main-paper element} & \textbf{Detailed supplementary evidence} & \textbf{Status}\\
\midrule
Controlled ITS laboratory & Criterion definitions, capability pathways, copula contexts and
opportunity functions in Sect.~2. & EDA\\
Oracle and global attribution & Closed-form functional, joint-background moments, \XGB,
TreeSHAP and B100 checks in Sect.~3. & EDA\\
Structural correction & Calibration, positive control, numerical null and candidate screening
in Sect.~4. & EDA\\
D29 decision geometry & World-resolved oracle orderings, winners, modal shares and regret in
Sect.~5. & EDA\\
Pilot-B evidence (RQ2) & Per-world method contrasts, gate decision and objective-weight vectors
in Sects.~5 and 7. & PBA\\
Observed selection maps (DQ1) & All audited modal identities, shares, unique-choice counts and
pairwise agreement summaries in Sect.~5. & PCD\\
Decision operators and gate & Exact weighting, tie, regret and hard-stop rules in Sects.~6--7.
& EDA\\
Unexecuted primary plan & Clearly labelled archival protocols in Sects.~8--12. & ANE\\
Integrity and closure & Content digests, software tests and permanent-stop record in
Sects.~13--14. & EDA\\
\bottomrule
\end{tabularx}
\end{table}

\section{Detailed Benchmark Construction}

\subsection{Criteria and Capability Pathways}

The six alternatives in the main manuscript were evaluated on the ten technology-agnostic
criteria in Table~\ref{supp:tab:criteria}. The original lifecycle-burden criterion was a cost;
all other criteria were benefits.

\begin{table}[t]
\centering
\caption{Controlled evaluation criteria.}
\label{supp:tab:criteria}
\small
\begin{tabularx}{\linewidth}{@{}c X c@{}}
\toprule
Code & Criterion & Original direction\\
\midrule
$C_1$ & Mobility-efficiency improvement & Benefit\\
$C_2$ & Travel-time reliability improvement & Benefit\\
$C_3$ & Person-throughput improvement & Benefit\\
$C_4$ & Safety-surrogate improvement & Benefit\\
$C_5$ & Non-recurring-congestion resilience & Benefit\\
$C_6$ & Environmental-efficiency improvement & Benefit\\
$C_7$ & User and accessibility benefit & Benefit\\
$C_8$ & Infrastructure and data readiness & Benefit\\
$C_9$ & Interoperability and scalability & Benefit\\
$C_{10}$ & Annualised lifecycle deployment burden & Cost\\
\bottomrule
\end{tabularx}
\end{table}

The pre-specified capability mask in Table~\ref{supp:tab:mask} defines synthetic functional
pathways, not empirical impact magnitudes. Direct (D), indirect (I) and zero (0) pathways set
broad response-amplitude bands.

\begin{table}[t]
\centering
\caption{Pre-specified structural capability mask.}
\label{supp:tab:mask}
\scriptsize
\begin{tabular}{@{}lccccccc ccc@{}}
\toprule
 & $C_1$ & $C_2$ & $C_3$ & $C_4$ & $C_5$ & $C_6$ & $C_7$ & Ready & Interop. & Burden\\
\midrule
ATSC      & D & D & D & I & I & I & I & M    & H & M\\
TSP       & I & D & D & I & 0 & I & D & M    & H & L--M\\
TIDM      & I & D & I & D & D & I & I & M    & H & M\\
RMTI--DRG & I & I & I & I & D & I & D & L--M & H & L--M\\
V2X--CS   & I & I & 0 & D & I & I & I & H    & H & H\\
DSM       & D & D & I & D & D & I & I & M--H & H & M--H\\
\bottomrule
\end{tabular}
\end{table}

\subsection{Gaussian-Copula Contexts and Opportunity Functions}

For demand, incident, person-movement, environmental and readiness factors, generate
$z_{0,s},z_{1,s},\ldots,z_{5,s}\overset{iid}{\sim}\mathcal N(0,1)$ and set
\begin{equation}
\widetilde h_{k,s}^{(\rho)}
=\sqrt{\rho}\,z_{0,s}+\sqrt{1-\rho}\,z_{k,s},
\qquad h_{k,s}^{(\rho)}=\Phi(\widetilde h_{k,s}^{(\rho)}).
\label{supp:eq:copula}
\end{equation}
The same base draws were reused across dependence levels, while estimation and external TEST
streams used disjoint namespaces. The first seven criterion opportunities were
\begin{align}
o_{s1}&=h_D, & o_{s2}&=0.60h_D+0.40h_I, & o_{s3}&=0.60h_D+0.40h_T,\nonumber\\
o_{s4}&=0.50h_D+0.50h_I, & o_{s5}&=h_I, & o_{s6}&=0.60h_D+0.40h_E,\nonumber\\
o_{s7}&=0.40h_D+0.60h_T.&&&
\label{supp:eq:opportunities}
\end{align}

Each replication seed $r$ defined one benchmark world $W_r$: technology parameters, a master
context stream and all common-random-number streams. The conditional method target had been
\begin{equation}
\Psi_m(x)=\E_W\!\left[\E_{H,\varepsilon}
\{L_m(x;W,H,\varepsilon)\mid W\}\right].
\label{supp:eq:estimand}
\end{equation}
The world, rather than an alternative--context row, was the intended statistical unit. The
primary worlds needed to estimate this quantity were never accessed after the gate fired.

\subsection{Historical and Frozen Response Layers}

For $C_1$--$C_7$, capability amplitudes were drawn once per world:
\begin{equation}
\theta_{aj}\sim
\begin{cases}
0,&M_{aj}=0,\\
\mathcal U(0.05,0.20),&M_{aj}=I,\\
\mathcal U(0.20,0.40),&M_{aj}=D.
\end{cases}
\label{supp:eq:effectbands}
\end{equation}
The historical v2.1 response was
\begin{equation}
x_{asj}=\clip\!\left(\theta_{aj}o_{sj}+\sigma_x\xi_{asj},0,1\right).
\label{supp:eq:v21}
\end{equation}
The deployment classes were
\begin{equation}
L\sim U(0.25,0.40),\qquad
M\sim U(0.45,0.60),\qquad
H\sim U(0.65,0.80).
\label{supp:eq:deploymentclasses}
\end{equation}
Compound classes were equal-probability mixtures. With readiness
requirement $r_a$, interoperability $\iota_a$ and burden $\kappa_a$,
\begin{align}
x_{as8}&=\clip\!\left(1-\max(0,r_a-h_{R,s})+\eta_{as8},0,1\right),\nonumber\\
x_{as9}&=\clip\!\left(\iota_a+0.10h_{R,s}+\eta_{as9},0,1\right),\nonumber\\
x_{as,10}&=\clip\!\left(\kappa_a+0.20(1-h_{R,s})r_a+\eta_{as,10},0,1\right).
\label{supp:eq:c8c10}
\end{align}
Benefit orientation used $g_{asj}=x_{asj}$ for $j\le9$ and
$g_{as,10}=1-x_{as,10}$. At $\sigma_x=0$, the first seven criteria reduced to
$g_{asj}=\theta_{aj}o_{sj}$, which explains the exact normalization cancellation in the main
paper.

The frozen D29 kernel was
\begin{equation}
g_{asj}=\theta_{aj}o_{sj}+0.05v_{aj}o_{sj}(1-o_{sj}),\qquad j\le7.
\label{supp:eq:d29}
\end{equation}
For $m_j$ active alternatives, its ordered contrast grid was
$v^{grid}_{\ell j}=-1+2(\ell-1)/(m_j-1)$. After sorting active IDs, one deterministic NumPy
permutation used
\begin{equation}
\texttt{SeedSequence([r, 2010, j])}.
\label{supp:eq:d29seedsequence}
\end{equation}
No outcome or
winner information entered this assignment. Active $\theta_{aj}\ge0.05$ ensured
$\partial g/\partial o=\theta_{aj}+0.05v_{aj}(1-2o)\ge0$. Structural-zero pathways remained
zero, and $C_8$--$C_{10}$ retained Eq.~\eqref{supp:eq:c8c10}.

For $o_{sj}>0$, within-context vector normalization over active alternatives gives the exact
geometry
\begin{equation}
r_{asj}=
\frac{\theta_{aj}+0.05v_{aj}(1-o_{sj})}
{\sqrt{\sum_{b\in\mathcal A_j}
[\theta_{bj}+0.05v_{bj}(1-o_{sj})]^2}}.
\label{supp:eq:d29normalized}
\end{equation}
Thus D29 removes the v2.1 cancellation, but each criterion-specific normalized profile remains
on a one-parameter curve indexed by $o_{sj}$. Because opportunities differ across criteria,
Eq.~\eqref{supp:eq:d29normalized} alone does not prove that a fixed-weight MOORA winner is
context invariant. At Pilot-B closure, fixed-weight invariance was therefore an untested
hypothesis rather than an explanation of the stop. The subsequently frozen descriptive
selection-map diagnostic reported below resolves this question for the seven observed vectors
only; it does not establish a causal effect of Eq.~\eqref{supp:eq:d29normalized} or characterize
unobserved weight vectors.

\section{Oracle, Learning and Background Details}

\subsection{Oracle Utility and Closed-Form Attribution}

The monotone transform, coefficient vector and interaction graph were
\begin{align}
q_j(g)&=\frac{\log(1+2g)}{\log3},\nonumber\\
\boldsymbol\alpha&=(0.11,0.04,0.05,0.07,0.10,0.14,0.12,0.16,0.08,0.13),\nonumber\\
\mathcal E&=\{(1,2),(3,7),(4,5),(6,10),(8,9)\},\qquad \beta_{jk}=0.20.
\label{supp:eq:oraclecomponents}
\end{align}
For $z_j=q_j(g_j)$, the empirical FIT-background moments were
$\mu_j=\E_{bg}(z_j)$ and $\mu_{jk}=\E_{bg}(z_jz_k)$; dependence was retained by not replacing
$\mu_{jk}$ with $\mu_j\mu_k$. The main-effect contribution was
$\phi_j^{main}=\alpha_j(z_j-\mu_j)$. For pair $\beta_{jk}z_jz_k$, feature $j$ received
\begin{equation}
\phi_j^{(jk)}=\frac{\beta_{jk}}{2}
\{z_j\mu_k-\mu_{jk}+z_jz_k-\mu_jz_k\},
\label{supp:eq:pairshap}
\end{equation}
with the symmetric expression for $k$. The closed-form attribution was
\begin{equation}
\phi_j^\star(\mathbf g)=\frac{1}{1+\lambda}
\left[\alpha_j(z_j-\mu_j)+\lambda\!\sum_{k:(j,k)\in\mathcal E}\phi_j^{(jk)}\right],
\label{supp:eq:closedshap}
\end{equation}
with symmetric handling when $j$ was second in a pair. Exhaustive coalition enumeration was a
unit-test reference. Global weights normalized mean absolute WEIGHT-row attributions; total
importance at most $10^{-12}$ made the vector undefined, with no Equal substitution.

\subsection{XGBoost and TreeSHAP Freeze}

XGBoost was selected once by context-grouped cross-validation on development FIT rows at the
reference condition. Candidate 45 attained mean RMSE $0.017757058913365983$ with
\begin{equation}
\begin{split}
&n_{\mathrm{estimators}}=600,\quad d_{\max}=2,\quad \eta=0.05,\\
&\text{subsample}=0.8,\quad \text{colsample\_bytree}=0.8,\quad
\text{min\_child\_weight}=5,\\
&\text{reg\_alpha}=0,\quad \text{reg\_lambda}=1.
\end{split}
\label{supp:eq:xgbparams}
\end{equation}
The objective was \texttt{reg:squarederror}, with one computational thread and estimator
random-state namespace 82002. The selected model was then fixed. Interventional TreeSHAP used 100 deterministic FIT
identities selected with row-priority namespace 83001 and evaluated mean absolute attributions
on WEIGHT. External TEST did not enter calibration, background selection or weighting.

\subsection{Matched B100 Oracle Diagnostic}

The primary oracle used all eligible FIT rows for $\mu_j$ and $\mu_{jk}$. The secondary B100
diagnostic changed only these moments to the exact 100 frozen TreeSHAP identities; WEIGHT rows,
the oracle functional and every other setting were unchanged. It compared full-FIT and B100
weight vectors through mean absolute difference, total variation, maximum absolute difference,
Spearman correlation, Top-3/Top-5 overlap and maximum marginal/joint moment changes. Undefined
vectors remained undefined. The only authorized pre-primary execution was seed 21001 at
$N=250$, $\rho=0.4$ and $\lambda=0.5$; no $Y$, SHAP package call or MCDM operation was involved.

\section{Detailed Structural Calibration and Candidate Screening}

\subsection{Structural and Reachability Metrics}

Let $G_j\in\mathbb R^{S\times A_j}$ denote the active-alternative response matrix for criterion
$j$ at zero measurement noise. If its singular values are
$\sigma_{1j}\ge\sigma_{2j}\ge\cdots$ and
$E_j=\sum_\ell\sigma_{\ell j}^2$, the rank-one residual was
\begin{equation}
NS_j=1-\frac{\sigma_{1j}^2}{E_j},
\qquad E_j>10^{-24}.
\label{supp:eq:NS}
\end{equation}
No epsilon was inserted into this ratio; structural-zero matrices were reported as undefined.
For each jointly positive active alternative pair,
\begin{equation}
LRV_{ab,j}=\SD_s\!\left[\log(g_{asj}/g_{bsj})\right],
\qquad
LRV50_j=\median_{a<b}LRV_{ab,j},
\label{supp:eq:LRV}
\end{equation}
using population SD. For normalization family $m$,
\begin{equation}
NSV_{j,m}=\left\{S^{-1}\sum_s
\left\|\mathbf r_{sj}^{(m)}-\bar{\mathbf r}_j^{(m)}\right\|_2^2\right\}^{1/2}.
\label{supp:eq:NSV}
\end{equation}
Vector normalization was the primary structural gate; sum, max and min--max variants were
robustness diagnostics.

Reachability reflected one idiosyncratic Gaussian-copula shock while holding the shared and all
other idiosyncratic shocks fixed. With opportunity functions recomputed after reflection,
\begin{equation}
RE_{kj}=\median_{s,a\in\mathcal A_j}|g_{asj}^{(-k)}-g_{asj}|,
\qquad
SRE_{kj}=\frac{RE_{kj}}{\IQR(g_j)+10^{-12}}.
\label{supp:eq:SRE}
\end{equation}

\subsection{Scientific Positive Control}

For active $C_1$--$C_7$ pathways, the dedicated positive control was
\begin{equation}
g_{asj}^{PC,\ell}=\left[\theta_{aj}+0.05c_a^{(\ell)}(2o_{sj}-1)\right]o_{sj},
\label{supp:eq:positivecontrol}
\end{equation}
with contrast base $(-1,-0.6,-0.2,0.2,0.6,1)$ and all six cyclic rotations. Structural-zero
pathways remained zero and $C_8$--$C_{10}$ were unchanged. Rotation medians were formed within
development worlds 21001--21005 and world medians were then used as scientific references.

\begin{table}[t]
\centering
\caption{Frozen D2.7 criterion-level scientific references.}
\label{supp:tab:d27layera}
\small
\begin{tabular}{@{}ccc@{}}
\toprule
Criterion & $T^{sci}_{LRV,j}$ & $T^{sci}_{NSV,j}$\\
\midrule
$C_1$ & 0.185606 & 0.101685\\
$C_2$ & 0.107107 & 0.069436\\
$C_3$ & 0.163029 & 0.064900\\
$C_4$ & 0.115796 & 0.069423\\
$C_5$ & 0.127964 & 0.073681\\
$C_6$ & 0.194899 & 0.129935\\
$C_7$ & 0.144080 & 0.073092\\
\bottomrule
\end{tabular}
\end{table}

\begin{table}[t]
\centering
\caption{Frozen D2.7 latent-pathway reachability references.}
\label{supp:tab:d27sre}
\small
\begin{tabular}{@{}lc@{}}
\toprule
Pathway & $T^{sci}_{SRE,kj}$\\
\midrule
$h_D\rightarrow C_1$ & 0.556245\\
$h_D\rightarrow C_2$ & 0.334437\\
$h_D\rightarrow C_3$ & 0.284878\\
$h_D\rightarrow C_4$ & 0.292844\\
$h_D\rightarrow C_6$ & 0.431612\\
$h_D\rightarrow C_7$ & 0.252576\\
$h_E\rightarrow C_6$ & 0.296869\\
$h_I\rightarrow C_2$ & 0.212043\\
$h_I\rightarrow C_4$ & 0.294905\\
$h_I\rightarrow C_5$ & 0.455271\\
$h_T\rightarrow C_3$ & 0.186995\\
$h_T\rightarrow C_7$ & 0.373380\\
\bottomrule
\end{tabular}
\end{table}

The paired positive-control uplift in mean normalized constant-modal regret was
$(0.001485,0.020022,-0.013836,0.015534,0.006161)$ for worlds 21001--21005. Because the activation
rule required positive uplift in all five worlds, no Layer-B scientific regret reference was
defined and no post-hoc value was introduced.

\subsection{Numerical Null and Candidate Outcomes}

Numerical thresholds were calibrated from 2000 exact rank-one replicates for each of five- and
six-alternative active sets:
\begin{equation}
T^{num}_{NS}=10^{-12},\qquad
T^{num}_{LRV}=10^{-10},\qquad
T^{num}_{NSV}=3.2506084454\times10^{-8}.
\label{supp:eq:nullthresholds}
\end{equation}
All 35 historical v2.1 world--criterion rows fell below every threshold. For active pathways
\(j\leq7\), the D4-F1 curvature family was
\(g_{asj}^{D4F1}(\chi)=\theta_{aj}o_{sj}[1+\chi v_{aj}(1-o_{sj})]\), with structural-zero
pathways fixed at zero. It cleared some structural references at higher curvature but lost reachability; no
$\chi\in\{0.1,\ldots,1.0\}$ passed all gates, so the family was rejected.

The later exact D2.9 audit evaluated all $6^5=7776$ prespecified positive-control profile
assignments. None met the old simultaneous binary conjunction. This established non-viability
for that calibrated design, not mathematical impossibility. Under the corrected prospective
rule, 54 candidates were eligible. D29 was retained solely for protocol continuity and minimal
intervention, not scientific preference, downstream performance, SHAP pattern or ITS identity.

The count of 54 is composed rather than directly observed in a single artefact, and is
decomposed here to prevent confusion. The nonselective retrospective re-adjudication examined 54
historical parameter points and returned 44 \texttt{VERIFIED\_\allowbreak ELIGIBLE},
0 \texttt{VERIFIED\_\allowbreak INELIGIBLE} and 10
\texttt{INSUFFICIENT\_\allowbreak FROZEN\_\allowbreak EVIDENCE}. The 10
unresolved points were all v2.2 points whose frozen artefact schema did not record C8--C10
invariance. A separate supplemental resolution then established that invariance from the
contemporaneous frozen source contract, and all 10 satisfied the corrected hard constraints.
The frozen retention artefact therefore records an eligible universe of
$44+10=54$ specifications. Eligibility is explicitly non-unique: the retention of D29 is a
minimal-intervention and provenance-stability decision, not a superiority, uniqueness or
statistical-selection claim, and no ranking or re-optimization over the 54 points was
performed.

\section{Detailed Development Audit Results}

\subsection{D29 Layer-B Decision Geometry}

The one-shot characterization used development worlds 21001--21005, the complete $N=1000$
FIT+WEIGHT pool, $\rho=0.4$, $\sigma_x=0$ and $\lambda=0.5$. It generated no $Y$ and executed
no learned model, SHAP package or MCDM method. Regret used the exact context-specific oracle
range and was undefined only if that range was at most $10^{-12}$; no context was undefined.

\begin{table}[t]
\centering
\caption{Complete frozen D29 Layer-B development record.}
\label{supp:tab:d29layerb}
\scriptsize
\resizebox{\linewidth}{!}{%
\begin{tabular}{@{}r r r l r r r r@{}}
\toprule
World & Orderings & Winners & Modal & Modal share & Mean regret & Median regret & $P_{95}$ regret\\
\midrule
21001 & 23 & 2 & $A_6$ & 0.983 & 0.00042018 & 0 & 0\\
21002 & 78 & 5 & $A_1$ & 0.519 & 0.05992337 & 0 & 0.29155475\\
21003 & 24 & 4 & $A_2$ & 0.745 & 0.04523642 & 0 & 0.28747533\\
21004 & 17 & 3 & $A_2$ & 0.827 & 0.01734351 & 0 & 0.12641643\\
21005 & 38 & 4 & $A_2$ & 0.809 & 0.01606539 & 0 & 0.12654230\\
\bottomrule
\end{tabular}}
\end{table}

Across worlds, the medians were 24 orderings, four winners, modal share $0.809$, mean regret
$0.01734351$, median regret zero and $P_{95}$ regret $0.12654230$. At the directly comparable
world 21001, the historical v2.1 record had 19 orderings, two winners, modal share $0.962$ and
mean always-modal regret $0.0011106$ under its frozen $D_s+10^{-12}$ denominator. D29 had 23
orderings, two winners, modal share $0.983$ and mean regret $0.00042018$ under exact $D_s$.
The comparison therefore supports neither universal degeneration nor universal improvement.

\subsection{XGBoost and TreeSHAP Development Checks}

The selected XGBoost hyperparameters are reported in Eq.~\eqref{supp:eq:xgbparams}; candidate
number and grouped-CV RMSE are given above. At world 21001 and the reference condition, total mean-absolute
TreeSHAP importance was $0.07006131011183041$ and maximum local-accuracy error was
$6.45\times10^{-7}$. The development-only weight vector for $(C_1,\ldots,C_{10})$ was
\begin{equation}
\begin{split}
&(0.124860,\ 0.113194,\ 0.069426,\ 0.025945,\ 0.155169,\\
&\qquad 0.078339,\ 0.071621,\ 0.265059,\ 0.020729,\ 0.075658).
\end{split}
\label{supp:eq:devshapweights}
\end{equation}
These observations did not tune comparator definitions, MCDM rules or winner semantics.

\subsection{B100-Matched Background Result}

Both oracle vectors were defined, and the 100 identities exactly matched the frozen TreeSHAP
membership. Total mean-absolute importance was $0.06961866590$ under full-FIT moments and
$0.06992189186$ under B100 moments.

\begin{table}[t]
\centering
\caption{Development-only oracle-attribution weights under full-FIT and matched B100 moments.}
\label{supp:tab:b100weights}
\small
\begin{tabular}{@{}lrrr@{}}
\toprule
Criterion & Full-FIT & B100 & Absolute difference\\
\midrule
$C_1$  & 0.09663438 & 0.09712037 & 0.00048600\\
$C_2$  & 0.04849942 & 0.04854063 & 0.00004122\\
$C_3$  & 0.05910838 & 0.05913637 & 0.00002799\\
$C_4$  & 0.06665455 & 0.06559847 & 0.00105608\\
$C_5$  & 0.14880967 & 0.14667297 & 0.00213669\\
$C_6$  & 0.09368358 & 0.09388953 & 0.00020596\\
$C_7$  & 0.07615062 & 0.07594850 & 0.00020212\\
$C_8$  & 0.24766693 & 0.25232065 & 0.00465372\\
$C_9$  & 0.04819959 & 0.04935317 & 0.00115358\\
$C_{10}$ & 0.11459290 & 0.11141932 & 0.00317357\\
\bottomrule
\end{tabular}
\end{table}

The weight-vector discrepancies were
\begin{equation}
\operatorname{MAE}_w=0.0013136928,\quad
TV_w=0.0065684640,\quad
\max_j|\Delta w_j|=0.0046537223.
\label{supp:eq:b100distances}
\end{equation}
Spearman correlation was $0.98787879$, and Top-3 and Top-5 overlaps were both $1.0$.
Maximum absolute marginal- and joint-moment changes were $0.0098907938$ and $0.0041811926$,
respectively. Equal fallback was not used. These results show small
finite-background sensitivity at this one development reference, not uniform robustness.

\subsection{World-Resolved Pilot-B Oracle and SHAP Comparison}

Table~\ref{supp:tab:pilotboracleshap} retains the full comparison moved from the main paper.
Each entry is an advantage over the world-specific MajorityWinner reference; positive values
favour the structured method.

\begin{table}[t]
\centering
\caption{Per-world closed-form oracle-attribution and SHAP advantages over MajorityWinner.}
\label{supp:tab:pilotboracleshap}
\small
\begin{tabular}{@{}crrrr@{}}
\toprule
World & Oracle $\Delta$Top-1 & SHAP $\Delta$Top-1 & Oracle $\Delta$regret & SHAP $\Delta$regret\\
\midrule
21001 & $ 0.000$ & $ 0.000$ & $ 0.000000$ & $ 0.000000$\\
21002 & $ 0.035$ & $ 0.030$ & $ 0.001538$ & $ 0.003064$\\
21003 & $-0.465$ & $ 0.020$ & $-0.143522$ & $ 0.006717$\\
21004 & $-0.010$ & $ 0.000$ & $-0.002119$ & $ 0.000000$\\
21005 & $-0.380$ & $-0.035$ & $-0.090328$ & $-0.002986$\\
\bottomrule
\end{tabular}
\end{table}

The closed-form method exactly implements the frozen global interventional attribution
functional; it is not an exact representation of context-specific oracle utility. Thus its
failure to escape MajorityWinner is evidence about the frozen global-compression--MOORA
pipeline, not a numerical error in the attribution formula.

\subsection{Illustrative Frozen Objective-Weight Vectors}

Table~\ref{supp:tab:objectivevectors} directly transcribes the CRITIC and Entropy vectors stored
for development world 21001; no new transformation of TEST decisions was performed. The display
cannot alter the frozen gate or establish a performance mechanism. Both vectors were defined
without Equal fallback, and neither has its largest component on C8 or C10. Thus concentration
on those two criteria is not supported in this world. The table is deliberately not combined
with five-world median performance to infer cross-world geometry.

\begin{table}[t]
\centering
\caption{Frozen objective-weight vectors for development world 21001.}
\label{supp:tab:objectivevectors}
\scriptsize
\setlength{\tabcolsep}{2.8pt}
\begin{tabular}{@{}lrrrrrrrrrr@{}}
\toprule
Method & C1 & C2 & C3 & C4 & C5 & C6 & C7 & C8 & C9 & C10\\
\midrule
CRITIC & .079 & .075 & .101 & .093 & .150 & .099 & .097 & .111 & .084 & .111\\
Entropy & .131 & .114 & .175 & .099 & .236 & .064 & .084 & .027 & .032 & .038\\
\bottomrule
\end{tabular}
\end{table}

The aggregate coincidence in world 21003 did not, by itself, identify a common decision map.
Equal, CRITIC and Entropy each recorded Top-1 accuracy $0.190$ and mean normalized regret
$0.316758$, but their mean Kendall $\tau_b$ values were respectively $0.5493$, $0.5427$ and
$0.5113$. Top-1 accuracy and normalized regret depend on the selected alternative, whereas
Kendall $\tau_b$ depends on the complete ranking. The $\tau_b$ differences therefore establish
different lower-order ranking behaviour on average but are logically uninformative about
argmax equality. That ambiguity motivated the separately frozen diagnostic below.

\subsection{Frozen Post-Closure Fixed-Weight Selection-Map Diagnostic (DQ1)}
\label{supp:sec:fixedselection}

\noindent\textbf{Evidential status: post-closure, descriptive, non-confirmatory.} This
subsection answers DQ1 only. It cannot modify, reinterpret, rescue or reopen the Pilot-B gate,
it defines no pass/fail criterion, and none of its content is confirmatory evidence for RQ1 or
RQ2.

The diagnostic was specified only after the Pilot-B result had been independently audited,
frozen and permanently closed. Its protocol was committed before any context-level
fixed-weight result was reconstructed. It used worlds 21001--21005 at the same reference
condition and the 200 frozen external TEST identities per world. The six structured vectors
were read from the immutable Pilot-B result without re-estimation; Equal was exactly
$(0.1,\ldots,0.1)$. D29 TEST decision matrices and oracle utilities were reconstructed only to
validate the archived aggregates. No target $Y$, model fit/prediction, SHAP, PI, RandomWeights,
bootstrap or protected seed was executed.

For each world, method and context the result stores the selected alternative; all six MOORA
scores and ranks; oracle utilities, ranks and selected alternative; Top-1 match; normalized
regret contribution; and agreement with the frozen MajorityWinner identity. Exactly
$5\times7\times200=7000$ records were independently audited. The audit reconstructed every
selection, anchor-based rank and regret directly from the stored full-precision vectors,
matched all 35 archived method--world aggregates within $10^{-12}$, and rebuilt all 105
pairwise selection-agreement summaries without executing the generator.

Table~\ref{supp:tab:fixedselectionmaps} reports each method's modal selected alternative,
modal share and number of distinct selected alternatives. Entries use the format
``alternative/\allowbreak share/\allowbreak $k$'', where $k=1$ denotes
context-invariant selection in that world.
MajorityWinner (MW) is a single frozen alternative by construction.

\begin{table}[t]
\centering
\caption{Audited selection maps for the seven observed fixed vectors. Each method entry is
modal alternative/modal share/number of distinct selected alternatives over 200 TEST contexts.}
\label{supp:tab:fixedselectionmaps}
\scriptsize
\setlength{\tabcolsep}{2.1pt}
\begin{tabular}{@{}rclllllll@{}}
\toprule
World & MW & Oracle & SHAP & PI & Ridge+ & CRITIC & Entropy & Equal\\
\midrule
21001 & A6 & A6/1.000/1 & A6/1.000/1 & A6/.980/2 & A6/1.000/1 & A6/1.000/1 & A6/1.000/1 & A6/1.000/1\\
21002 & A1 & A1/.700/3 & A1/.950/2 & A1/.630/2 & A1/.875/3 & A3/.800/3 & A1/1.000/1 & A1/1.000/1\\
21003 & A2 & A6/.720/4 & A2/.690/3 & A2/.530/3 & A6/.495/3 & A6/1.000/1 & A6/1.000/1 & A6/1.000/1\\
21004 & A2 & A2/.990/2 & A2/1.000/1 & A2/.775/4 & A2/.990/2 & A3/.750/3 & A3/.985/3 & A2/1.000/1\\
21005 & A2 & A6/.350/4 & A2/.790/3 & A3/.675/5 & A6/.375/4 & A6/.915/3 & A6/.800/3 & A2/1.000/1\\
\bottomrule
\end{tabular}
\end{table}

Equal was context invariant in all five worlds and exactly matched MajorityWinner across all
200 contexts in worlds 21001, 21002, 21004 and 21005. In world 21003, Equal instead selected A6
throughout while MajorityWinner was A2. The numbers of constant-selection worlds were
Oracle 1, SHAP 2, PI 0, Ridge+ 1, CRITIC 2, Entropy 3 and Equal 5; exact all-context
MajorityWinner agreement occurred in 1, 2, 0, 1, 1, 2 and 4 worlds, respectively.

The world-21003 ambiguity is now directly resolved: Equal, CRITIC and Entropy all selected A6
in every TEST context. Their different mean $\tau_b$ values therefore reflect differences below
the argmax, not different selected alternatives. More broadly, the heterogeneous constancy
counts show that fixed-vector MOORA was neither universally context invariant nor uniformly
context responsive over the seven observed vectors. The result supports coarse
observed-vector decision regions as a descriptive hypothesis, but cannot establish the causal
mechanism of the Pilot-B stop, characterize unobserved regions of the weight simplex, rank
methods or generalize beyond these five development worlds.

The new artefact makes paired context resampling technically possible, but the prospectively
frozen diagnostic explicitly excluded bootstrap and defined no inferential or pass/fail role.
No resampling was executed. Any future paired sensitivity analysis requires its own protocol,
frozen before inspecting bootstrap output, and cannot revise the closed Pilot-B decision.

\section{Detailed Weighting and Decision Operators}

The primary information partitions were fixed: Oracle, SHAP, PI, CRITIC and Entropy used WEIGHT;
Ridge+ used FIT; no method used external TEST for fitting or weight estimation. Oracle and SHAP
normalized mean absolute attributions. SHAP was declared undefined if total importance was at
most $10^{-12}$ and was never replaced by Equal weights.

CRITIC applied criterion-wise min--max scaling, population SD and Pearson conflict over
nonconstant columns. Entropy applied the same min--max scaling followed by Shannon
diversification with $0\log0=0$. Ridge+ standardized predictors and target with fold-local
population moments, solved non-negative L2-augmented least squares without an intercept, and
used the frozen grid
\begin{equation}
\tau\in\{10^{-4},10^{-3},10^{-2},10^{-1},1,10,100\}.
\label{supp:eq:ridgegrid}
\end{equation}
Five-fold context-grouped original-scale RMSE selected $\tau$, with the smallest value as tie
break.

Permutation importance kept the fitted XGBoost model fixed. For each criterion, it copied the
complete six-alternative block from source context $\pi(d)$ to destination context $d$ under 20
unique derangements. Their frozen generator was
\begin{equation}
\texttt{SeedSequence([replication, 84001, N])}.
\label{supp:eq:piseedsequence}
\end{equation}
The same
derangements were reused across criteria and factor levels for a given seed and $N$. Importance
was permuted minus baseline WEIGHT MSE; negative means were truncated before normalization and
repeat dispersion used sample SD. CRITIC, Entropy, Ridge+ and PI used Equal weights only under
their frozen zero-information rules, with fallback recorded.

Equal was $(0.1,\ldots,0.1)$. RandomWeights used 200
$\operatorname{Dirichlet}(1,\ldots,1)$ vectors per world from namespace 81001 and reused them
across conditions. MajorityWinner estimated the oracle modal winner on WEIGHT and applied that
constant choice to TEST. MajorityWinner was intentionally oracle-informed: it was the
context-free oracle-modal alternative estimated on WEIGHT and evaluated on TEST, and it was
optimal only for empirical oracle Top-1 frequency among constant choices. It is not claimed to
minimize regret or to constitute an operational transportation policy. It was neither a
deployable competitor nor a universally optimal choice, and it was not stakeholder-derived or
normative. The regret arm of the frozen gate uses this same
prospectively frozen Top-1-derived constant reference rather than a separately optimized
regret-minimizing constant policy; a future benchmark could prospectively define metric-matched
constant references. The frozen gate itself is unchanged by this observation. Direct-XGBoost was a parallel
predictive-decision reference rather than a weighting
method. MOORA used within-context vector normalization and a weighted sum; TOPSIS used the same
normalized benefit matrix as a conditional aggregation robustness operator.

\subsection{Frozen Tie Semantics and Archived Decision Margin}

Within each context, alternatives were sorted by descending raw score and ascending
\texttt{alternative\_id}. The highest unassigned score anchored a tie group; every remaining
alternative within absolute distance $10^{-12}$ of that anchor joined the group and received
the average of the occupied one-based rank positions. Ties were not chained through adjacent
near-equalities. Deterministic Top-1 used
\begin{equation}
\mathcal T_s=\{a:S_{\max,s}-S_{as}\le10^{-12}\}
\label{supp:eq:top1tieset}
\end{equation}
and selected the smallest alternative ID in $\mathcal T_s$. Relative tolerance was zero. These
semantics applied to oracle, MOORA, TOPSIS and Direct-XGBoost scores.

The conditional primary plan also defined the guarded oracle decision margin
\begin{equation}
M_s=\frac{U_{(1),s}^\star-U_{(2),s}^\star}
{\max_aU_{as}^\star-\min_aU_{as}^\star+10^{-12}}.
\label{supp:eq:margin}
\end{equation}
Margin-conditioned analyses were not executed because Pilot B stopped the primary study.

\section{Frozen Pilot-B Gate and Complete Decision}

Pilot B reused development worlds 21001--21005 at the frozen reference
\begin{equation}
(N,c,\rho,\lambda)=(250,0.30,0.4,0.5).
\label{supp:eq:pilotbreference}
\end{equation}
Each world used 200 external TEST contexts and 200 common
RandomWeights draws per world. It was a protocol-level adequacy screen, not independent
population validation. A structured method was eligible only if it was defined without Equal
fallback in all five worlds and Kendall's $\tau_b$ was defined in at least 190 TEST contexts per
world; at least 190 RandomWeights draws also had to be eligible. Cross-world advantages used
medians and required strict positivity in at least four worlds.

The five independently sufficient hard stops were: incomplete coverage; no eligible method
clearing median $\Delta\tau_b=0.05$ or $\Delta$mean-regret $=0.01$ versus RandomWeights; no
eligible method clearing median $\Delta$Top-1 $=0.02$ or $\Delta$mean-regret $=0.01$ versus
MajorityWinner; pooled oracle modal share at least $0.90$; or the frozen random-weight
insensitivity criterion, including its $0.90$ invariant fraction and joint width limits
$0.05$ ($\tau_b$), $0.02$ (Top-1) and $0.01$ (mean regret) in at least four worlds. Passage
required every stop to be false.

\begin{table}[t]
\centering
\caption{Complete frozen Pilot-B gate decision.}
\label{supp:tab:pilotbgates}
\small
\begin{tabularx}{\linewidth}{@{}lXc@{}}
\toprule
Gate & Audited result & Stop\\
\midrule
Coverage & Six structured methods eligible without fallback & No\\
Random equivalence & Four of six methods separated from RandomWeights & No\\
Modal-winner effective tie & Zero of six methods escaped MajorityWinner & \textbf{Yes}\\
Oracle-winner dominance & Pooled modal share $0.526<0.90$ & No\\
Weight insensitivity & Invariant fraction $0.004<0.90$; width-qualified worlds $0<4$ & No\\
\midrule
Overall OR rule & At least one frozen stop fired & \textbf{Yes}\\
\bottomrule
\end{tabularx}
\end{table}

The best strictly positive consistency count was two of five worlds for either Top-1 or mean
regret, and the best cross-world medians were zero. The independently reconstructed decision
therefore barred the protected factorial. This is a negative result about the adequacy of the
frozen comparison at one condition, not evidence that SHAP, weights or MCDM fail universally.

\section{Archived Stage-Wise Fidelity Ladder}

The conditional primary analysis contained nine representations:
\begin{enumerate}
\item \texttt{OracleUtility}, the full interacting oracle;
\item \texttt{OracleMainEffectReference}, with the interaction removed;
\item \texttt{OracleGlobalAttributionNonlinearQ};
\item \texttt{SHAPGlobalAttributionNonlinearQ};
\item \texttt{OracleGlobalAttributionLinearG};
\item \texttt{SHAPGlobalAttributionLinearG};
\item \texttt{OracleWeightMOORA};
\item \texttt{SHAPWeightMOORA};
\item \texttt{DirectXGBoost}, a parallel reference rather than a sequential stage.
\end{enumerate}
Eight paired contrasts isolated interaction removal, oracle global compression, SHAP estimation
on the nonlinear scale, nonlinear-to-linear transformation under each weight vector, MOORA
transformation under each vector and post-MOORA Oracle-versus-SHAP attribution. Kendall's
$\tau_b$ and normalized oracle regret were specified for every representation. The ladder was
diagnostic, not an additive error decomposition, and was not executed because Pilot B stopped
the primary study.

\section{Conditional Primary Factorial}

The four planned primary factors were
\begin{align}
N&\in\{25,50,100,250,1000\}, &
c&\in\{0.10,0.30,0.60\},\label{supp:eq:factors1}\\
\rho&\in\{0,0.4,0.8\}, &
\lambda&\in\{0,0.5,1\}.
\label{supp:eq:factors2}
\end{align}
The resulting $5\times3\times3\times3=135$ conditions were to be repeated over 30
prespecified primary benchmark-instance seeds, giving 4050 seed--condition replications. The
reference condition was
\begin{equation}
(N,c,\rho,\lambda)=(250,0.30,0.4,0.5).
\label{supp:eq:reference}
\end{equation}

\begin{table}[t]
\centering
\caption{Archived nested estimation partitions and fixed TEST size.}
\label{supp:tab:partitions}
\small
\begin{tabular}{@{}rrrr@{}}
\toprule
$N$ & FIT & WEIGHT & external TEST\\
\midrule
25 & 20 & 5 & 200\\
50 & 40 & 10 & 200\\
100 & 80 & 20 & 200\\
250 & 200 & 50 & 200\\
1000 & 800 & 200 & 200\\
\bottomrule
\end{tabular}
\end{table}

Technology-response parameters and common-random-number streams were to be sampled once per
seed and reused across factor levels. For each $(\rho,\lambda)$ cell, the oracle signal scale
$s_U$ was to be calculated on the complete 1000-context FIT+WEIGHT master pool so that changing
$N$ would not change the imposed noise-to-signal ratio. Realized signal-to-noise ratio was to
be logged for every cell. FIT, WEIGHT and external TEST partitions were to remain
context-disjoint, and no TEST row was to be used for model fitting or weight estimation.

\section{Archived Alpha-Dispersion Stress Analysis}

The conditional primary analysis retained the frozen heterogeneous coefficient vector. At the
reference condition only, three additional coefficient structures had been fixed: balanced,
dispersion-aligned and dispersion-anti-aligned. The mapping was obtained from an independent
200-world technology ensemble using
\begin{equation}
D_{j,w}^{attr}
=\frac{1}{M_w}\sum_i\left|q_j(g_{ijw})-\bar q_{j,w}\right|,
\qquad
\bar D_j^{attr}=\E_W[D_{j,W}^{attr}].
\label{supp:eq:qmad}
\end{equation}
The complete ordering agreed between 150 and 200 worlds:
\begin{equation}
C_8>C_{10}>C_5>C_2>C_3>C_4>C_1>C_7>C_6>C_9.
\label{supp:eq:d3order}
\end{equation}

\begin{table}[t]
\centering
\caption{Frozen 200-world mean attribution-scale dispersion.}
\label{supp:tab:d3dispersion}
\small
\begin{tabular}{@{}lc@{}}
\toprule
Criterion & Mean dispersion\\
\midrule
$C_8$ & 0.122873493\\
$C_{10}$ & 0.119403289\\
$C_5$ & 0.110114266\\
$C_2$ & 0.094820283\\
$C_3$ & 0.090482741\\
$C_4$ & 0.088452971\\
$C_1$ & 0.087135074\\
$C_7$ & 0.078706704\\
$C_6$ & 0.049591336\\
$C_9$ & 0.028266887\\
\bottomrule
\end{tabular}
\end{table}

Balanced coefficients were $\alpha_j=0.10$. The aligned and anti-aligned mappings were
\begin{align}
\boldsymbol\alpha^{aligned}:
&(C_8,C_{10},C_5,C_2,C_3,C_4,C_1,C_7,C_6,C_9)
\nonumber\\
&\mapsto(0.16,0.14,0.13,0.12,0.11,0.10,0.08,0.07,0.05,0.04),
\label{supp:eq:alphaaligned}\\
\boldsymbol\alpha^{anti}:
&(C_9,C_6,C_7,C_1,C_4,C_3,C_2,C_5,C_{10},C_8)
\nonumber\\
&\mapsto(0.16,0.14,0.13,0.12,0.11,0.10,0.08,0.07,0.05,0.04).
\label{supp:eq:alphaanti}
\end{align}
The frozen primary heterogeneous assignment had weak rank association with the dispersion
ordering ($\rho_S=0.1273$, $\tau_b=0.0222$), so it was not changed. Although the D3 mapping
itself was a completed development artefact, the planned primary stress comparisons were never
executed.

\section{Archived Bootstrap and Robustness Plan}

At the reference condition, $B=200$ context-level bootstrap replications were to resample FIT
and WEIGHT contexts separately, refit the model, recompute attributions and weights, and
evaluate decisions on the fixed original external TEST contexts. Confidence intervals were to
be reported for criterion weights, total-variation weight error, Kendall's $\tau_b$, Top-1
accuracy and normalized oracle regret. No cross-context probability that a particular ITS was
ranked first was to be interpreted as a global deployment probability.

A 50-oracle ensemble was to leave the frozen D29 technology-response kernel unchanged while
varying oracle coefficient vectors, nonlinearities and semantically admissible interaction
pairs. Relative outcome noise was to be recalibrated to each oracle variant's signal standard
deviation. This ensemble would not replace the response generator and therefore would not
reintroduce the historical v2.1 separability defect through oracle-only perturbations.

Two distribution-shift analyses were specified: holding out the upper 20\% of demand pressure
and holding out the lower 20\% of infrastructure readiness. Secondary robustness analyses were
to include criterion removal, $\pm10/20/30\%$ SHAP-weight perturbations,
leave-one-alternative-out analysis and MOORA--TOPSIS comparison.

\section{Archived Statistical Analysis Plan}

Factor effects were to be summarized with repeated-seed regression or mixed-effects models,
emphasizing effect sizes and uncertainty intervals rather than binary significance alone.
Method comparisons were paired because every method would share the same benchmark instance.
The 30 primary seeds were to be treated as benchmark-instance clusters; alternative--context
rows were not to be interpreted as independent superpopulation replications.

Method availability was to be reported explicitly. For each factorial cell, undefined
Oracle/SHAP weight rates and recorded Equal-fallback rates were to accompany conditional
decision-fidelity summaries. Pairwise contrasts were to use only benchmark instances on which
both methods were defined, with coverage and fallback rates reported alongside them. Undefined
Oracle or SHAP vectors were never to be silently replaced by Equal weights.

Planned outputs included predictive-performance and Direct-XGBoost references, weight-recovery
tables, decision-fidelity summaries, factorial effect summaries, alpha-alignment stress results,
bootstrap intervals, validity-boundary figures and the prespecified stage-wise fidelity ladder.
None of these primary outputs exists under v1.

\section{Reproducibility and Content-Integrity Record}
\label{supp:sec:integrity}

Table~\ref{supp:tab:digests} gives non-identifying truncated content digests. The selected B100
background exactly matched the frozen TreeSHAP identity set.

\begin{table}[t]
\centering
\caption{Truncated, non-identifying content digests of the frozen artefacts underlying the
reported numbers.}
\label{supp:tab:digests}
\small
\begin{tabular}{@{}ll@{}}
\toprule
Artifact & Truncated SHA-256\\
\midrule
Oracle B100 result & \texttt{7f9864bd...439c9fb0}\\
TreeSHAP background identities & \texttt{3184c675...1646eda1}\\
Pilot-B result & \texttt{05072857...9f787a88}\\
Fixed-weight selection-map result & \texttt{a1f293b6...d7e20f4e}\\
\bottomrule
\end{tabular}
\end{table}

The full anonymous provenance package records, in order, the
threshold freeze, evaluator/executor freeze, scientifically invariant direct-file loader repair,
audited Pilot-B result freeze, immutable closure reference, post-closure diagnostic protocol
freeze, implementation freeze and independently audited result freeze. The worktree was clean
after each closure point, and the final complete test run reported 836 passes with three
pre-existing external plotting warnings.
\ifanon
Exact repository identifiers are intentionally omitted here to protect anonymous review.
\else
Exact repository identifiers accompany the public reproducibility archive.
\fi
These facts document software and artefact integrity; they are not statistical evidence.

\section{Prospective Stop and Permanent Closure}

Pilot B was governed by five independently sufficient hard stops: coverage failure, failure to
separate RandomWeights, failure to escape MajorityWinner, pooled oracle-winner dominance and
weight insensitivity. Passage required every flag to be false. In the frozen result, only the
modal-winner effective-tie flag was true, but the overall OR rule therefore barred every
analysis in this supplement.

The immutable version-control, digest and test record appears in
Section~\ref{supp:sec:integrity}. Pilot B v1 is closed and must not be rerun. Any later study
requires a new specification and version reference, explicit disclosure of the v1 failure and
development-only design work; it cannot alter the v1 thresholds, result or closure reference.
The fixed-weight selection-map diagnostic in Section~\ref{supp:sec:fixedselection} is such a
separately frozen descriptive layer; it neither reopened Pilot B nor authorized any protected
analysis.

\end{document}
"
% ============================================================
% ChkTeX: intentional en dashes, literal digest ellipses, label placement,
% acronym punctuation and LNCS name/bibliography conventions.
% chktex-file 8
% chktex-file 11
% chktex-file 13
% chktex-file 24
% chktex-file 36
% chktex-file 38

\documentclass[runningheads,orivec]{llncs}

\usepackage[T1]{fontenc}
\usepackage[utf8]{inputenc}
\usepackage{amsmath,amssymb,amsfonts}
\usepackage{booktabs}
\usepackage{graphicx}
\usepackage{tabularx}
\usepackage{hyperref}
\hypersetup{hidelinks}

\newcommand{\E}{\mathbb{E}}
\newcommand{\median}{\operatorname{median}}
\newcommand{\IQR}{\operatorname{IQR}}
\newcommand{\SD}{\operatorname{SD}}
\newcommand{\clip}{\operatorname{clip}}
\newcommand{\XGB}{XGBoost}

% ---- build-mode switch: anonymous by default -----------------
\providecommand{\CAMERAREADY}{0}
\newif\ifanon
\anontrue
\ifnum\CAMERAREADY=1 \anonfalse \fi
% --------------------------------------------------------------

\begin{document}
\emergencystretch=1.5em

\title{Supplementary Methods and Audit Archive for\texorpdfstring{\\}{ }
``From Predictive Importance to Decision Value: Evaluating XAI-Based Weights for
ITS Prioritisation''}
\titlerunning{Supplementary Benchmark-Adequacy Record}

\ifanon
  \author{Anonymous Author(s)}
  \authorrunning{Anonymous Author(s)}
  \institute{Anonymous Institution(s)}
\else
  \author{Mohamed Yassine Neifar\inst{1,2} \and
  Ahmed Frikha\inst{2} \and
  Nadir Farhi\inst{1}}
  \authorrunning{Neifar et al.}
  \institute{Universit\'e Gustave Eiffel, COSYS--GRETTIA,
  Champs-sur-Marne, F-77447 Marne-la-Vall\'ee Cedex 2, France
  \and
  Institut Sup\'erieur de Gestion Industrielle de Sfax (ISGIS),
  OLID, University of Sfax, Sfax 3038, Tunisia}
\fi

\maketitle

\section{Status and Interpretation Boundary}

This supplement reports executed development audits and preserves the detailed primary and
robustness analyses that had been specified before the prospectively frozen Pilot-B
benchmark-adequacy gate was executed. The latter analyses were conditional on all five Pilot-B
hard-stop flags being false. The modal-winner effective-tie stop fired, so the primary factorial
was not authorized. Primary seeds 11001--11030 and reserve seeds 30001--30005 were never
accessed; all sections explicitly labelled ``Archived'' remain unexecuted protocols rather than
results.

The closed v1 state is anchored by an immutable version-control reference.
\ifanon
Exact repository identifiers and complete digests are withheld from this anonymous reviewer
file; they are available confidentially to editors and will accompany the public
reproducibility archive after review.
\else
Exact repository identifiers and complete digests accompany the public reproducibility archive.
\fi
This document is therefore a methods, audit-results and archival-protocol supplement. It is not
authorization to execute the protected design.
After permanent Pilot-B closure, one separately frozen development-only descriptive diagnostic
reconstructed the context-level MOORA selection maps of the seven observed archived fixed
vectors. It used no primary/reserve world, generated no target outcome, fitted or queried no
model, re-estimated no weight and could not alter the gate. Its protocol, implementation,
independent audit and result were frozen as a distinct post-closure layer.
Methodological citations are provided in the main manuscript and are not duplicated here.

\subsection*{Question Structure and Evidential Status}

The main manuscript states two prospective research questions and one separate post-closure
descriptive question. RQ1 asks whether structural diagnostics can reveal that apparent
contextual variation disappears under the normalization used by the downstream MCDM operator.
RQ2 asks whether, after that structural defect is corrected, structured global-weight pipelines
provide sufficient incremental decision value over random-weight and constant-decision
references to justify protected primary evaluation. Both are addressed by the prospectively
structured benchmark-development record and the frozen Pilot-B result reported here.

DQ1 asks what the archived context-level selection maps reveal about the observed decision
behaviour of the seven frozen fixed-weight vectors. DQ1 is post-closure and descriptive. It is
non-confirmatory, it carries no pass/fail role, and it cannot modify, reinterpret or reopen the
Pilot-B gate. Section~\ref{supp:sec:fixedselection} reports it under that label.

Throughout this supplement, \emph{each replication seed defines one benchmark world}, and the
world is the statistical unit. Where a numbered identifier such as 21001 appears it names a
world; the phrase ``replication seed'' is retained only when the frozen artefact field or
protocol clause is itself named that way. \emph{Validity} denotes conceptual scientific
properties of the benchmark; \emph{adequacy} denotes the operational prospectively frozen gate
used to decide whether protected primary execution was warranted.

\subsection*{Reader Map}

The conference-facing main manuscript follows a problem--method--result--meaning structure and
deliberately omits most development chronology. Table~\ref{supp:tab:readermap} directs readers
to the corresponding detailed evidence. This separation changes presentation only: the frozen
scientific record, Pilot-B decision and post-closure interpretation boundary are unchanged.

\begin{table}[t]
\centering
\caption{Navigation from the concise main argument to the detailed reproducibility record,
with the evidential status of each block. Status codes are: EDA, executed development audit;
PBA, Pilot-B confirmatory adequacy result; PCD, post-closure descriptive diagnostic; ANE,
archived and not executed.}
\label{supp:tab:readermap}
\small
\begin{tabularx}{\linewidth}{@{}lXc@{}}
\toprule
\textbf{Main-paper element} & \textbf{Detailed supplementary evidence} & \textbf{Status}\\
\midrule
Controlled ITS laboratory & Criterion definitions, capability pathways, copula contexts and
opportunity functions in Sect.~2. & EDA\\
Oracle and global attribution & Closed-form functional, joint-background moments, \XGB,
TreeSHAP and B100 checks in Sect.~3. & EDA\\
Structural correction & Calibration, positive control, numerical null and candidate screening
in Sect.~4. & EDA\\
D29 decision geometry & World-resolved oracle orderings, winners, modal shares and regret in
Sect.~5. & EDA\\
Pilot-B evidence (RQ2) & Per-world method contrasts, gate decision and objective-weight vectors
in Sects.~5 and 7. & PBA\\
Observed selection maps (DQ1) & All audited modal identities, shares, unique-choice counts and
pairwise agreement summaries in Sect.~5. & PCD\\
Decision operators and gate & Exact weighting, tie, regret and hard-stop rules in Sects.~6--7.
& EDA\\
Unexecuted primary plan & Clearly labelled archival protocols in Sects.~8--12. & ANE\\
Integrity and closure & Content digests, software tests and permanent-stop record in
Sects.~13--14. & EDA\\
\bottomrule
\end{tabularx}
\end{table}

\section{Detailed Benchmark Construction}

\subsection{Criteria and Capability Pathways}

The six alternatives in the main manuscript were evaluated on the ten technology-agnostic
criteria in Table~\ref{supp:tab:criteria}. The original lifecycle-burden criterion was a cost;
all other criteria were benefits.

\begin{table}[t]
\centering
\caption{Controlled evaluation criteria.}
\label{supp:tab:criteria}
\small
\begin{tabularx}{\linewidth}{@{}c X c@{}}
\toprule
Code & Criterion & Original direction\\
\midrule
$C_1$ & Mobility-efficiency improvement & Benefit\\
$C_2$ & Travel-time reliability improvement & Benefit\\
$C_3$ & Person-throughput improvement & Benefit\\
$C_4$ & Safety-surrogate improvement & Benefit\\
$C_5$ & Non-recurring-congestion resilience & Benefit\\
$C_6$ & Environmental-efficiency improvement & Benefit\\
$C_7$ & User and accessibility benefit & Benefit\\
$C_8$ & Infrastructure and data readiness & Benefit\\
$C_9$ & Interoperability and scalability & Benefit\\
$C_{10}$ & Annualised lifecycle deployment burden & Cost\\
\bottomrule
\end{tabularx}
\end{table}

The pre-specified capability mask in Table~\ref{supp:tab:mask} defines synthetic functional
pathways, not empirical impact magnitudes. Direct (D), indirect (I) and zero (0) pathways set
broad response-amplitude bands.

\begin{table}[t]
\centering
\caption{Pre-specified structural capability mask.}
\label{supp:tab:mask}
\scriptsize
\begin{tabular}{@{}lccccccc ccc@{}}
\toprule
 & $C_1$ & $C_2$ & $C_3$ & $C_4$ & $C_5$ & $C_6$ & $C_7$ & Ready & Interop. & Burden\\
\midrule
ATSC      & D & D & D & I & I & I & I & M    & H & M\\
TSP       & I & D & D & I & 0 & I & D & M    & H & L--M\\
TIDM      & I & D & I & D & D & I & I & M    & H & M\\
RMTI--DRG & I & I & I & I & D & I & D & L--M & H & L--M\\
V2X--CS   & I & I & 0 & D & I & I & I & H    & H & H\\
DSM       & D & D & I & D & D & I & I & M--H & H & M--H\\
\bottomrule
\end{tabular}
\end{table}

\subsection{Gaussian-Copula Contexts and Opportunity Functions}

For demand, incident, person-movement, environmental and readiness factors, generate
$z_{0,s},z_{1,s},\ldots,z_{5,s}\overset{iid}{\sim}\mathcal N(0,1)$ and set
\begin{equation}
\widetilde h_{k,s}^{(\rho)}
=\sqrt{\rho}\,z_{0,s}+\sqrt{1-\rho}\,z_{k,s},
\qquad h_{k,s}^{(\rho)}=\Phi(\widetilde h_{k,s}^{(\rho)}).
\label{supp:eq:copula}
\end{equation}
The same base draws were reused across dependence levels, while estimation and external TEST
streams used disjoint namespaces. The first seven criterion opportunities were
\begin{align}
o_{s1}&=h_D, & o_{s2}&=0.60h_D+0.40h_I, & o_{s3}&=0.60h_D+0.40h_T,\nonumber\\
o_{s4}&=0.50h_D+0.50h_I, & o_{s5}&=h_I, & o_{s6}&=0.60h_D+0.40h_E,\nonumber\\
o_{s7}&=0.40h_D+0.60h_T.&&&
\label{supp:eq:opportunities}
\end{align}

Each replication seed $r$ defined one benchmark world $W_r$: technology parameters, a master
context stream and all common-random-number streams. The conditional method target had been
\begin{equation}
\Psi_m(x)=\E_W\!\left[\E_{H,\varepsilon}
\{L_m(x;W,H,\varepsilon)\mid W\}\right].
\label{supp:eq:estimand}
\end{equation}
The world, rather than an alternative--context row, was the intended statistical unit. The
primary worlds needed to estimate this quantity were never accessed after the gate fired.

\subsection{Historical and Frozen Response Layers}

For $C_1$--$C_7$, capability amplitudes were drawn once per world:
\begin{equation}
\theta_{aj}\sim
\begin{cases}
0,&M_{aj}=0,\\
\mathcal U(0.05,0.20),&M_{aj}=I,\\
\mathcal U(0.20,0.40),&M_{aj}=D.
\end{cases}
\label{supp:eq:effectbands}
\end{equation}
The historical v2.1 response was
\begin{equation}
x_{asj}=\clip\!\left(\theta_{aj}o_{sj}+\sigma_x\xi_{asj},0,1\right).
\label{supp:eq:v21}
\end{equation}
The deployment classes were
\begin{equation}
L\sim U(0.25,0.40),\qquad
M\sim U(0.45,0.60),\qquad
H\sim U(0.65,0.80).
\label{supp:eq:deploymentclasses}
\end{equation}
Compound classes were equal-probability mixtures. With readiness
requirement $r_a$, interoperability $\iota_a$ and burden $\kappa_a$,
\begin{align}
x_{as8}&=\clip\!\left(1-\max(0,r_a-h_{R,s})+\eta_{as8},0,1\right),\nonumber\\
x_{as9}&=\clip\!\left(\iota_a+0.10h_{R,s}+\eta_{as9},0,1\right),\nonumber\\
x_{as,10}&=\clip\!\left(\kappa_a+0.20(1-h_{R,s})r_a+\eta_{as,10},0,1\right).
\label{supp:eq:c8c10}
\end{align}
Benefit orientation used $g_{asj}=x_{asj}$ for $j\le9$ and
$g_{as,10}=1-x_{as,10}$. At $\sigma_x=0$, the first seven criteria reduced to
$g_{asj}=\theta_{aj}o_{sj}$, which explains the exact normalization cancellation in the main
paper.

The frozen D29 kernel was
\begin{equation}
g_{asj}=\theta_{aj}o_{sj}+0.05v_{aj}o_{sj}(1-o_{sj}),\qquad j\le7.
\label{supp:eq:d29}
\end{equation}
For $m_j$ active alternatives, its ordered contrast grid was
$v^{grid}_{\ell j}=-1+2(\ell-1)/(m_j-1)$. After sorting active IDs, one deterministic NumPy
permutation used
\begin{equation}
\texttt{SeedSequence([r, 2010, j])}.
\label{supp:eq:d29seedsequence}
\end{equation}
No outcome or
winner information entered this assignment. Active $\theta_{aj}\ge0.05$ ensured
$\partial g/\partial o=\theta_{aj}+0.05v_{aj}(1-2o)\ge0$. Structural-zero pathways remained
zero, and $C_8$--$C_{10}$ retained Eq.~\eqref{supp:eq:c8c10}.

For $o_{sj}>0$, within-context vector normalization over active alternatives gives the exact
geometry
\begin{equation}
r_{asj}=
\frac{\theta_{aj}+0.05v_{aj}(1-o_{sj})}
{\sqrt{\sum_{b\in\mathcal A_j}
[\theta_{bj}+0.05v_{bj}(1-o_{sj})]^2}}.
\label{supp:eq:d29normalized}
\end{equation}
Thus D29 removes the v2.1 cancellation, but each criterion-specific normalized profile remains
on a one-parameter curve indexed by $o_{sj}$. Because opportunities differ across criteria,
Eq.~\eqref{supp:eq:d29normalized} alone does not prove that a fixed-weight MOORA winner is
context invariant. At Pilot-B closure, fixed-weight invariance was therefore an untested
hypothesis rather than an explanation of the stop. The subsequently frozen descriptive
selection-map diagnostic reported below resolves this question for the seven observed vectors
only; it does not establish a causal effect of Eq.~\eqref{supp:eq:d29normalized} or characterize
unobserved weight vectors.

\section{Oracle, Learning and Background Details}

\subsection{Oracle Utility and Closed-Form Attribution}

The monotone transform, coefficient vector and interaction graph were
\begin{align}
q_j(g)&=\frac{\log(1+2g)}{\log3},\nonumber\\
\boldsymbol\alpha&=(0.11,0.04,0.05,0.07,0.10,0.14,0.12,0.16,0.08,0.13),\nonumber\\
\mathcal E&=\{(1,2),(3,7),(4,5),(6,10),(8,9)\},\qquad \beta_{jk}=0.20.
\label{supp:eq:oraclecomponents}
\end{align}
For $z_j=q_j(g_j)$, the empirical FIT-background moments were
$\mu_j=\E_{bg}(z_j)$ and $\mu_{jk}=\E_{bg}(z_jz_k)$; dependence was retained by not replacing
$\mu_{jk}$ with $\mu_j\mu_k$. The main-effect contribution was
$\phi_j^{main}=\alpha_j(z_j-\mu_j)$. For pair $\beta_{jk}z_jz_k$, feature $j$ received
\begin{equation}
\phi_j^{(jk)}=\frac{\beta_{jk}}{2}
\{z_j\mu_k-\mu_{jk}+z_jz_k-\mu_jz_k\},
\label{supp:eq:pairshap}
\end{equation}
with the symmetric expression for $k$. The closed-form attribution was
\begin{equation}
\phi_j^\star(\mathbf g)=\frac{1}{1+\lambda}
\left[\alpha_j(z_j-\mu_j)+\lambda\!\sum_{k:(j,k)\in\mathcal E}\phi_j^{(jk)}\right],
\label{supp:eq:closedshap}
\end{equation}
with symmetric handling when $j$ was second in a pair. Exhaustive coalition enumeration was a
unit-test reference. Global weights normalized mean absolute WEIGHT-row attributions; total
importance at most $10^{-12}$ made the vector undefined, with no Equal substitution.

\subsection{XGBoost and TreeSHAP Freeze}

XGBoost was selected once by context-grouped cross-validation on development FIT rows at the
reference condition. Candidate 45 attained mean RMSE $0.017757058913365983$ with
\begin{equation}
\begin{split}
&n_{\mathrm{estimators}}=600,\quad d_{\max}=2,\quad \eta=0.05,\\
&\text{subsample}=0.8,\quad \text{colsample\_bytree}=0.8,\quad
\text{min\_child\_weight}=5,\\
&\text{reg\_alpha}=0,\quad \text{reg\_lambda}=1.
\end{split}
\label{supp:eq:xgbparams}
\end{equation}
The objective was \texttt{reg:squarederror}, with one computational thread and estimator
random-state namespace 82002. The selected model was then fixed. Interventional TreeSHAP used 100 deterministic FIT
identities selected with row-priority namespace 83001 and evaluated mean absolute attributions
on WEIGHT. External TEST did not enter calibration, background selection or weighting.

\subsection{Matched B100 Oracle Diagnostic}

The primary oracle used all eligible FIT rows for $\mu_j$ and $\mu_{jk}$. The secondary B100
diagnostic changed only these moments to the exact 100 frozen TreeSHAP identities; WEIGHT rows,
the oracle functional and every other setting were unchanged. It compared full-FIT and B100
weight vectors through mean absolute difference, total variation, maximum absolute difference,
Spearman correlation, Top-3/Top-5 overlap and maximum marginal/joint moment changes. Undefined
vectors remained undefined. The only authorized pre-primary execution was seed 21001 at
$N=250$, $\rho=0.4$ and $\lambda=0.5$; no $Y$, SHAP package call or MCDM operation was involved.

\section{Detailed Structural Calibration and Candidate Screening}

\subsection{Structural and Reachability Metrics}

Let $G_j\in\mathbb R^{S\times A_j}$ denote the active-alternative response matrix for criterion
$j$ at zero measurement noise. If its singular values are
$\sigma_{1j}\ge\sigma_{2j}\ge\cdots$ and
$E_j=\sum_\ell\sigma_{\ell j}^2$, the rank-one residual was
\begin{equation}
NS_j=1-\frac{\sigma_{1j}^2}{E_j},
\qquad E_j>10^{-24}.
\label{supp:eq:NS}
\end{equation}
No epsilon was inserted into this ratio; structural-zero matrices were reported as undefined.
For each jointly positive active alternative pair,
\begin{equation}
LRV_{ab,j}=\SD_s\!\left[\log(g_{asj}/g_{bsj})\right],
\qquad
LRV50_j=\median_{a<b}LRV_{ab,j},
\label{supp:eq:LRV}
\end{equation}
using population SD. For normalization family $m$,
\begin{equation}
NSV_{j,m}=\left\{S^{-1}\sum_s
\left\|\mathbf r_{sj}^{(m)}-\bar{\mathbf r}_j^{(m)}\right\|_2^2\right\}^{1/2}.
\label{supp:eq:NSV}
\end{equation}
Vector normalization was the primary structural gate; sum, max and min--max variants were
robustness diagnostics.

Reachability reflected one idiosyncratic Gaussian-copula shock while holding the shared and all
other idiosyncratic shocks fixed. With opportunity functions recomputed after reflection,
\begin{equation}
RE_{kj}=\median_{s,a\in\mathcal A_j}|g_{asj}^{(-k)}-g_{asj}|,
\qquad
SRE_{kj}=\frac{RE_{kj}}{\IQR(g_j)+10^{-12}}.
\label{supp:eq:SRE}
\end{equation}

\subsection{Scientific Positive Control}

For active $C_1$--$C_7$ pathways, the dedicated positive control was
\begin{equation}
g_{asj}^{PC,\ell}=\left[\theta_{aj}+0.05c_a^{(\ell)}(2o_{sj}-1)\right]o_{sj},
\label{supp:eq:positivecontrol}
\end{equation}
with contrast base $(-1,-0.6,-0.2,0.2,0.6,1)$ and all six cyclic rotations. Structural-zero
pathways remained zero and $C_8$--$C_{10}$ were unchanged. Rotation medians were formed within
development worlds 21001--21005 and world medians were then used as scientific references.

\begin{table}[t]
\centering
\caption{Frozen D2.7 criterion-level scientific references.}
\label{supp:tab:d27layera}
\small
\begin{tabular}{@{}ccc@{}}
\toprule
Criterion & $T^{sci}_{LRV,j}$ & $T^{sci}_{NSV,j}$\\
\midrule
$C_1$ & 0.185606 & 0.101685\\
$C_2$ & 0.107107 & 0.069436\\
$C_3$ & 0.163029 & 0.064900\\
$C_4$ & 0.115796 & 0.069423\\
$C_5$ & 0.127964 & 0.073681\\
$C_6$ & 0.194899 & 0.129935\\
$C_7$ & 0.144080 & 0.073092\\
\bottomrule
\end{tabular}
\end{table}

\begin{table}[t]
\centering
\caption{Frozen D2.7 latent-pathway reachability references.}
\label{supp:tab:d27sre}
\small
\begin{tabular}{@{}lc@{}}
\toprule
Pathway & $T^{sci}_{SRE,kj}$\\
\midrule
$h_D\rightarrow C_1$ & 0.556245\\
$h_D\rightarrow C_2$ & 0.334437\\
$h_D\rightarrow C_3$ & 0.284878\\
$h_D\rightarrow C_4$ & 0.292844\\
$h_D\rightarrow C_6$ & 0.431612\\
$h_D\rightarrow C_7$ & 0.252576\\
$h_E\rightarrow C_6$ & 0.296869\\
$h_I\rightarrow C_2$ & 0.212043\\
$h_I\rightarrow C_4$ & 0.294905\\
$h_I\rightarrow C_5$ & 0.455271\\
$h_T\rightarrow C_3$ & 0.186995\\
$h_T\rightarrow C_7$ & 0.373380\\
\bottomrule
\end{tabular}
\end{table}

The paired positive-control uplift in mean normalized constant-modal regret was
$(0.001485,0.020022,-0.013836,0.015534,0.006161)$ for worlds 21001--21005. Because the activation
rule required positive uplift in all five worlds, no Layer-B scientific regret reference was
defined and no post-hoc value was introduced.

\subsection{Numerical Null and Candidate Outcomes}

Numerical thresholds were calibrated from 2000 exact rank-one replicates for each of five- and
six-alternative active sets:
\begin{equation}
T^{num}_{NS}=10^{-12},\qquad
T^{num}_{LRV}=10^{-10},\qquad
T^{num}_{NSV}=3.2506084454\times10^{-8}.
\label{supp:eq:nullthresholds}
\end{equation}
All 35 historical v2.1 world--criterion rows fell below every threshold. For active pathways
\(j\leq7\), the D4-F1 curvature family was
\(g_{asj}^{D4F1}(\chi)=\theta_{aj}o_{sj}[1+\chi v_{aj}(1-o_{sj})]\), with structural-zero
pathways fixed at zero. It cleared some structural references at higher curvature but lost reachability; no
$\chi\in\{0.1,\ldots,1.0\}$ passed all gates, so the family was rejected.

The later exact D2.9 audit evaluated all $6^5=7776$ prespecified positive-control profile
assignments. None met the old simultaneous binary conjunction. This established non-viability
for that calibrated design, not mathematical impossibility. Under the corrected prospective
rule, 54 candidates were eligible. D29 was retained solely for protocol continuity and minimal
intervention, not scientific preference, downstream performance, SHAP pattern or ITS identity.

The count of 54 is composed rather than directly observed in a single artefact, and is
decomposed here to prevent confusion. The nonselective retrospective re-adjudication examined 54
historical parameter points and returned 44 \texttt{VERIFIED\_\allowbreak ELIGIBLE},
0 \texttt{VERIFIED\_\allowbreak INELIGIBLE} and 10
\texttt{INSUFFICIENT\_\allowbreak FROZEN\_\allowbreak EVIDENCE}. The 10
unresolved points were all v2.2 points whose frozen artefact schema did not record C8--C10
invariance. A separate supplemental resolution then established that invariance from the
contemporaneous frozen source contract, and all 10 satisfied the corrected hard constraints.
The frozen retention artefact therefore records an eligible universe of
$44+10=54$ specifications. Eligibility is explicitly non-unique: the retention of D29 is a
minimal-intervention and provenance-stability decision, not a superiority, uniqueness or
statistical-selection claim, and no ranking or re-optimization over the 54 points was
performed.

\section{Detailed Development Audit Results}

\subsection{D29 Layer-B Decision Geometry}

The one-shot characterization used development worlds 21001--21005, the complete $N=1000$
FIT+WEIGHT pool, $\rho=0.4$, $\sigma_x=0$ and $\lambda=0.5$. It generated no $Y$ and executed
no learned model, SHAP package or MCDM method. Regret used the exact context-specific oracle
range and was undefined only if that range was at most $10^{-12}$; no context was undefined.

\begin{table}[t]
\centering
\caption{Complete frozen D29 Layer-B development record.}
\label{supp:tab:d29layerb}
\scriptsize
\resizebox{\linewidth}{!}{%
\begin{tabular}{@{}r r r l r r r r@{}}
\toprule
World & Orderings & Winners & Modal & Modal share & Mean regret & Median regret & $P_{95}$ regret\\
\midrule
21001 & 23 & 2 & $A_6$ & 0.983 & 0.00042018 & 0 & 0\\
21002 & 78 & 5 & $A_1$ & 0.519 & 0.05992337 & 0 & 0.29155475\\
21003 & 24 & 4 & $A_2$ & 0.745 & 0.04523642 & 0 & 0.28747533\\
21004 & 17 & 3 & $A_2$ & 0.827 & 0.01734351 & 0 & 0.12641643\\
21005 & 38 & 4 & $A_2$ & 0.809 & 0.01606539 & 0 & 0.12654230\\
\bottomrule
\end{tabular}}
\end{table}

Across worlds, the medians were 24 orderings, four winners, modal share $0.809$, mean regret
$0.01734351$, median regret zero and $P_{95}$ regret $0.12654230$. At the directly comparable
world 21001, the historical v2.1 record had 19 orderings, two winners, modal share $0.962$ and
mean always-modal regret $0.0011106$ under its frozen $D_s+10^{-12}$ denominator. D29 had 23
orderings, two winners, modal share $0.983$ and mean regret $0.00042018$ under exact $D_s$.
The comparison therefore supports neither universal degeneration nor universal improvement.

\subsection{XGBoost and TreeSHAP Development Checks}

The selected XGBoost hyperparameters are reported in Eq.~\eqref{supp:eq:xgbparams}; candidate
number and grouped-CV RMSE are given above. At world 21001 and the reference condition, total mean-absolute
TreeSHAP importance was $0.07006131011183041$ and maximum local-accuracy error was
$6.45\times10^{-7}$. The development-only weight vector for $(C_1,\ldots,C_{10})$ was
\begin{equation}
\begin{split}
&(0.124860,\ 0.113194,\ 0.069426,\ 0.025945,\ 0.155169,\\
&\qquad 0.078339,\ 0.071621,\ 0.265059,\ 0.020729,\ 0.075658).
\end{split}
\label{supp:eq:devshapweights}
\end{equation}
These observations did not tune comparator definitions, MCDM rules or winner semantics.

\subsection{B100-Matched Background Result}

Both oracle vectors were defined, and the 100 identities exactly matched the frozen TreeSHAP
membership. Total mean-absolute importance was $0.06961866590$ under full-FIT moments and
$0.06992189186$ under B100 moments.

\begin{table}[t]
\centering
\caption{Development-only oracle-attribution weights under full-FIT and matched B100 moments.}
\label{supp:tab:b100weights}
\small
\begin{tabular}{@{}lrrr@{}}
\toprule
Criterion & Full-FIT & B100 & Absolute difference\\
\midrule
$C_1$  & 0.09663438 & 0.09712037 & 0.00048600\\
$C_2$  & 0.04849942 & 0.04854063 & 0.00004122\\
$C_3$  & 0.05910838 & 0.05913637 & 0.00002799\\
$C_4$  & 0.06665455 & 0.06559847 & 0.00105608\\
$C_5$  & 0.14880967 & 0.14667297 & 0.00213669\\
$C_6$  & 0.09368358 & 0.09388953 & 0.00020596\\
$C_7$  & 0.07615062 & 0.07594850 & 0.00020212\\
$C_8$  & 0.24766693 & 0.25232065 & 0.00465372\\
$C_9$  & 0.04819959 & 0.04935317 & 0.00115358\\
$C_{10}$ & 0.11459290 & 0.11141932 & 0.00317357\\
\bottomrule
\end{tabular}
\end{table}

The weight-vector discrepancies were
\begin{equation}
\operatorname{MAE}_w=0.0013136928,\quad
TV_w=0.0065684640,\quad
\max_j|\Delta w_j|=0.0046537223.
\label{supp:eq:b100distances}
\end{equation}
Spearman correlation was $0.98787879$, and Top-3 and Top-5 overlaps were both $1.0$.
Maximum absolute marginal- and joint-moment changes were $0.0098907938$ and $0.0041811926$,
respectively. Equal fallback was not used. These results show small
finite-background sensitivity at this one development reference, not uniform robustness.

\subsection{World-Resolved Pilot-B Oracle and SHAP Comparison}

Table~\ref{supp:tab:pilotboracleshap} retains the full comparison moved from the main paper.
Each entry is an advantage over the world-specific MajorityWinner reference; positive values
favour the structured method.

\begin{table}[t]
\centering
\caption{Per-world closed-form oracle-attribution and SHAP advantages over MajorityWinner.}
\label{supp:tab:pilotboracleshap}
\small
\begin{tabular}{@{}crrrr@{}}
\toprule
World & Oracle $\Delta$Top-1 & SHAP $\Delta$Top-1 & Oracle $\Delta$regret & SHAP $\Delta$regret\\
\midrule
21001 & $ 0.000$ & $ 0.000$ & $ 0.000000$ & $ 0.000000$\\
21002 & $ 0.035$ & $ 0.030$ & $ 0.001538$ & $ 0.003064$\\
21003 & $-0.465$ & $ 0.020$ & $-0.143522$ & $ 0.006717$\\
21004 & $-0.010$ & $ 0.000$ & $-0.002119$ & $ 0.000000$\\
21005 & $-0.380$ & $-0.035$ & $-0.090328$ & $-0.002986$\\
\bottomrule
\end{tabular}
\end{table}

The closed-form method exactly implements the frozen global interventional attribution
functional; it is not an exact representation of context-specific oracle utility. Thus its
failure to escape MajorityWinner is evidence about the frozen global-compression--MOORA
pipeline, not a numerical error in the attribution formula.

\subsection{Illustrative Frozen Objective-Weight Vectors}

Table~\ref{supp:tab:objectivevectors} directly transcribes the CRITIC and Entropy vectors stored
for development world 21001; no new transformation of TEST decisions was performed. The display
cannot alter the frozen gate or establish a performance mechanism. Both vectors were defined
without Equal fallback, and neither has its largest component on C8 or C10. Thus concentration
on those two criteria is not supported in this world. The table is deliberately not combined
with five-world median performance to infer cross-world geometry.

\begin{table}[t]
\centering
\caption{Frozen objective-weight vectors for development world 21001.}
\label{supp:tab:objectivevectors}
\scriptsize
\setlength{\tabcolsep}{2.8pt}
\begin{tabular}{@{}lrrrrrrrrrr@{}}
\toprule
Method & C1 & C2 & C3 & C4 & C5 & C6 & C7 & C8 & C9 & C10\\
\midrule
CRITIC & .079 & .075 & .101 & .093 & .150 & .099 & .097 & .111 & .084 & .111\\
Entropy & .131 & .114 & .175 & .099 & .236 & .064 & .084 & .027 & .032 & .038\\
\bottomrule
\end{tabular}
\end{table}

The aggregate coincidence in world 21003 did not, by itself, identify a common decision map.
Equal, CRITIC and Entropy each recorded Top-1 accuracy $0.190$ and mean normalized regret
$0.316758$, but their mean Kendall $\tau_b$ values were respectively $0.5493$, $0.5427$ and
$0.5113$. Top-1 accuracy and normalized regret depend on the selected alternative, whereas
Kendall $\tau_b$ depends on the complete ranking. The $\tau_b$ differences therefore establish
different lower-order ranking behaviour on average but are logically uninformative about
argmax equality. That ambiguity motivated the separately frozen diagnostic below.

\subsection{Frozen Post-Closure Fixed-Weight Selection-Map Diagnostic (DQ1)}
\label{supp:sec:fixedselection}

\noindent\textbf{Evidential status: post-closure, descriptive, non-confirmatory.} This
subsection answers DQ1 only. It cannot modify, reinterpret, rescue or reopen the Pilot-B gate,
it defines no pass/fail criterion, and none of its content is confirmatory evidence for RQ1 or
RQ2.

The diagnostic was specified only after the Pilot-B result had been independently audited,
frozen and permanently closed. Its protocol was committed before any context-level
fixed-weight result was reconstructed. It used worlds 21001--21005 at the same reference
condition and the 200 frozen external TEST identities per world. The six structured vectors
were read from the immutable Pilot-B result without re-estimation; Equal was exactly
$(0.1,\ldots,0.1)$. D29 TEST decision matrices and oracle utilities were reconstructed only to
validate the archived aggregates. No target $Y$, model fit/prediction, SHAP, PI, RandomWeights,
bootstrap or protected seed was executed.

For each world, method and context the result stores the selected alternative; all six MOORA
scores and ranks; oracle utilities, ranks and selected alternative; Top-1 match; normalized
regret contribution; and agreement with the frozen MajorityWinner identity. Exactly
$5\times7\times200=7000$ records were independently audited. The audit reconstructed every
selection, anchor-based rank and regret directly from the stored full-precision vectors,
matched all 35 archived method--world aggregates within $10^{-12}$, and rebuilt all 105
pairwise selection-agreement summaries without executing the generator.

Table~\ref{supp:tab:fixedselectionmaps} reports each method's modal selected alternative,
modal share and number of distinct selected alternatives. Entries use the format
``alternative/\allowbreak share/\allowbreak $k$'', where $k=1$ denotes
context-invariant selection in that world.
MajorityWinner (MW) is a single frozen alternative by construction.

\begin{table}[t]
\centering
\caption{Audited selection maps for the seven observed fixed vectors. Each method entry is
modal alternative/modal share/number of distinct selected alternatives over 200 TEST contexts.}
\label{supp:tab:fixedselectionmaps}
\scriptsize
\setlength{\tabcolsep}{2.1pt}
\begin{tabular}{@{}rclllllll@{}}
\toprule
World & MW & Oracle & SHAP & PI & Ridge+ & CRITIC & Entropy & Equal\\
\midrule
21001 & A6 & A6/1.000/1 & A6/1.000/1 & A6/.980/2 & A6/1.000/1 & A6/1.000/1 & A6/1.000/1 & A6/1.000/1\\
21002 & A1 & A1/.700/3 & A1/.950/2 & A1/.630/2 & A1/.875/3 & A3/.800/3 & A1/1.000/1 & A1/1.000/1\\
21003 & A2 & A6/.720/4 & A2/.690/3 & A2/.530/3 & A6/.495/3 & A6/1.000/1 & A6/1.000/1 & A6/1.000/1\\
21004 & A2 & A2/.990/2 & A2/1.000/1 & A2/.775/4 & A2/.990/2 & A3/.750/3 & A3/.985/3 & A2/1.000/1\\
21005 & A2 & A6/.350/4 & A2/.790/3 & A3/.675/5 & A6/.375/4 & A6/.915/3 & A6/.800/3 & A2/1.000/1\\
\bottomrule
\end{tabular}
\end{table}

Equal was context invariant in all five worlds and exactly matched MajorityWinner across all
200 contexts in worlds 21001, 21002, 21004 and 21005. In world 21003, Equal instead selected A6
throughout while MajorityWinner was A2. The numbers of constant-selection worlds were
Oracle 1, SHAP 2, PI 0, Ridge+ 1, CRITIC 2, Entropy 3 and Equal 5; exact all-context
MajorityWinner agreement occurred in 1, 2, 0, 1, 1, 2 and 4 worlds, respectively.

The world-21003 ambiguity is now directly resolved: Equal, CRITIC and Entropy all selected A6
in every TEST context. Their different mean $\tau_b$ values therefore reflect differences below
the argmax, not different selected alternatives. More broadly, the heterogeneous constancy
counts show that fixed-vector MOORA was neither universally context invariant nor uniformly
context responsive over the seven observed vectors. The result supports coarse
observed-vector decision regions as a descriptive hypothesis, but cannot establish the causal
mechanism of the Pilot-B stop, characterize unobserved regions of the weight simplex, rank
methods or generalize beyond these five development worlds.

The new artefact makes paired context resampling technically possible, but the prospectively
frozen diagnostic explicitly excluded bootstrap and defined no inferential or pass/fail role.
No resampling was executed. Any future paired sensitivity analysis requires its own protocol,
frozen before inspecting bootstrap output, and cannot revise the closed Pilot-B decision.

\section{Detailed Weighting and Decision Operators}

The primary information partitions were fixed: Oracle, SHAP, PI, CRITIC and Entropy used WEIGHT;
Ridge+ used FIT; no method used external TEST for fitting or weight estimation. Oracle and SHAP
normalized mean absolute attributions. SHAP was declared undefined if total importance was at
most $10^{-12}$ and was never replaced by Equal weights.

CRITIC applied criterion-wise min--max scaling, population SD and Pearson conflict over
nonconstant columns. Entropy applied the same min--max scaling followed by Shannon
diversification with $0\log0=0$. Ridge+ standardized predictors and target with fold-local
population moments, solved non-negative L2-augmented least squares without an intercept, and
used the frozen grid
\begin{equation}
\tau\in\{10^{-4},10^{-3},10^{-2},10^{-1},1,10,100\}.
\label{supp:eq:ridgegrid}
\end{equation}
Five-fold context-grouped original-scale RMSE selected $\tau$, with the smallest value as tie
break.

Permutation importance kept the fitted XGBoost model fixed. For each criterion, it copied the
complete six-alternative block from source context $\pi(d)$ to destination context $d$ under 20
unique derangements. Their frozen generator was
\begin{equation}
\texttt{SeedSequence([replication, 84001, N])}.
\label{supp:eq:piseedsequence}
\end{equation}
The same
derangements were reused across criteria and factor levels for a given seed and $N$. Importance
was permuted minus baseline WEIGHT MSE; negative means were truncated before normalization and
repeat dispersion used sample SD. CRITIC, Entropy, Ridge+ and PI used Equal weights only under
their frozen zero-information rules, with fallback recorded.

Equal was $(0.1,\ldots,0.1)$. RandomWeights used 200
$\operatorname{Dirichlet}(1,\ldots,1)$ vectors per world from namespace 81001 and reused them
across conditions. MajorityWinner estimated the oracle modal winner on WEIGHT and applied that
constant choice to TEST. MajorityWinner was intentionally oracle-informed: it was the
context-free oracle-modal alternative estimated on WEIGHT and evaluated on TEST, and it was
optimal only for empirical oracle Top-1 frequency among constant choices. It is not claimed to
minimize regret or to constitute an operational transportation policy. It was neither a
deployable competitor nor a universally optimal choice, and it was not stakeholder-derived or
normative. The regret arm of the frozen gate uses this same
prospectively frozen Top-1-derived constant reference rather than a separately optimized
regret-minimizing constant policy; a future benchmark could prospectively define metric-matched
constant references. The frozen gate itself is unchanged by this observation. Direct-XGBoost was a parallel
predictive-decision reference rather than a weighting
method. MOORA used within-context vector normalization and a weighted sum; TOPSIS used the same
normalized benefit matrix as a conditional aggregation robustness operator.

\subsection{Frozen Tie Semantics and Archived Decision Margin}

Within each context, alternatives were sorted by descending raw score and ascending
\texttt{alternative\_id}. The highest unassigned score anchored a tie group; every remaining
alternative within absolute distance $10^{-12}$ of that anchor joined the group and received
the average of the occupied one-based rank positions. Ties were not chained through adjacent
near-equalities. Deterministic Top-1 used
\begin{equation}
\mathcal T_s=\{a:S_{\max,s}-S_{as}\le10^{-12}\}
\label{supp:eq:top1tieset}
\end{equation}
and selected the smallest alternative ID in $\mathcal T_s$. Relative tolerance was zero. These
semantics applied to oracle, MOORA, TOPSIS and Direct-XGBoost scores.

The conditional primary plan also defined the guarded oracle decision margin
\begin{equation}
M_s=\frac{U_{(1),s}^\star-U_{(2),s}^\star}
{\max_aU_{as}^\star-\min_aU_{as}^\star+10^{-12}}.
\label{supp:eq:margin}
\end{equation}
Margin-conditioned analyses were not executed because Pilot B stopped the primary study.

\section{Frozen Pilot-B Gate and Complete Decision}

Pilot B reused development worlds 21001--21005 at the frozen reference
\begin{equation}
(N,c,\rho,\lambda)=(250,0.30,0.4,0.5).
\label{supp:eq:pilotbreference}
\end{equation}
Each world used 200 external TEST contexts and 200 common
RandomWeights draws per world. It was a protocol-level adequacy screen, not independent
population validation. A structured method was eligible only if it was defined without Equal
fallback in all five worlds and Kendall's $\tau_b$ was defined in at least 190 TEST contexts per
world; at least 190 RandomWeights draws also had to be eligible. Cross-world advantages used
medians and required strict positivity in at least four worlds.

The five independently sufficient hard stops were: incomplete coverage; no eligible method
clearing median $\Delta\tau_b=0.05$ or $\Delta$mean-regret $=0.01$ versus RandomWeights; no
eligible method clearing median $\Delta$Top-1 $=0.02$ or $\Delta$mean-regret $=0.01$ versus
MajorityWinner; pooled oracle modal share at least $0.90$; or the frozen random-weight
insensitivity criterion, including its $0.90$ invariant fraction and joint width limits
$0.05$ ($\tau_b$), $0.02$ (Top-1) and $0.01$ (mean regret) in at least four worlds. Passage
required every stop to be false.

\begin{table}[t]
\centering
\caption{Complete frozen Pilot-B gate decision.}
\label{supp:tab:pilotbgates}
\small
\begin{tabularx}{\linewidth}{@{}lXc@{}}
\toprule
Gate & Audited result & Stop\\
\midrule
Coverage & Six structured methods eligible without fallback & No\\
Random equivalence & Four of six methods separated from RandomWeights & No\\
Modal-winner effective tie & Zero of six methods escaped MajorityWinner & \textbf{Yes}\\
Oracle-winner dominance & Pooled modal share $0.526<0.90$ & No\\
Weight insensitivity & Invariant fraction $0.004<0.90$; width-qualified worlds $0<4$ & No\\
\midrule
Overall OR rule & At least one frozen stop fired & \textbf{Yes}\\
\bottomrule
\end{tabularx}
\end{table}

The best strictly positive consistency count was two of five worlds for either Top-1 or mean
regret, and the best cross-world medians were zero. The independently reconstructed decision
therefore barred the protected factorial. This is a negative result about the adequacy of the
frozen comparison at one condition, not evidence that SHAP, weights or MCDM fail universally.

\section{Archived Stage-Wise Fidelity Ladder}

The conditional primary analysis contained nine representations:
\begin{enumerate}
\item \texttt{OracleUtility}, the full interacting oracle;
\item \texttt{OracleMainEffectReference}, with the interaction removed;
\item \texttt{OracleGlobalAttributionNonlinearQ};
\item \texttt{SHAPGlobalAttributionNonlinearQ};
\item \texttt{OracleGlobalAttributionLinearG};
\item \texttt{SHAPGlobalAttributionLinearG};
\item \texttt{OracleWeightMOORA};
\item \texttt{SHAPWeightMOORA};
\item \texttt{DirectXGBoost}, a parallel reference rather than a sequential stage.
\end{enumerate}
Eight paired contrasts isolated interaction removal, oracle global compression, SHAP estimation
on the nonlinear scale, nonlinear-to-linear transformation under each weight vector, MOORA
transformation under each vector and post-MOORA Oracle-versus-SHAP attribution. Kendall's
$\tau_b$ and normalized oracle regret were specified for every representation. The ladder was
diagnostic, not an additive error decomposition, and was not executed because Pilot B stopped
the primary study.

\section{Conditional Primary Factorial}

The four planned primary factors were
\begin{align}
N&\in\{25,50,100,250,1000\}, &
c&\in\{0.10,0.30,0.60\},\label{supp:eq:factors1}\\
\rho&\in\{0,0.4,0.8\}, &
\lambda&\in\{0,0.5,1\}.
\label{supp:eq:factors2}
\end{align}
The resulting $5\times3\times3\times3=135$ conditions were to be repeated over 30
prespecified primary benchmark-instance seeds, giving 4050 seed--condition replications. The
reference condition was
\begin{equation}
(N,c,\rho,\lambda)=(250,0.30,0.4,0.5).
\label{supp:eq:reference}
\end{equation}

\begin{table}[t]
\centering
\caption{Archived nested estimation partitions and fixed TEST size.}
\label{supp:tab:partitions}
\small
\begin{tabular}{@{}rrrr@{}}
\toprule
$N$ & FIT & WEIGHT & external TEST\\
\midrule
25 & 20 & 5 & 200\\
50 & 40 & 10 & 200\\
100 & 80 & 20 & 200\\
250 & 200 & 50 & 200\\
1000 & 800 & 200 & 200\\
\bottomrule
\end{tabular}
\end{table}

Technology-response parameters and common-random-number streams were to be sampled once per
seed and reused across factor levels. For each $(\rho,\lambda)$ cell, the oracle signal scale
$s_U$ was to be calculated on the complete 1000-context FIT+WEIGHT master pool so that changing
$N$ would not change the imposed noise-to-signal ratio. Realized signal-to-noise ratio was to
be logged for every cell. FIT, WEIGHT and external TEST partitions were to remain
context-disjoint, and no TEST row was to be used for model fitting or weight estimation.

\section{Archived Alpha-Dispersion Stress Analysis}

The conditional primary analysis retained the frozen heterogeneous coefficient vector. At the
reference condition only, three additional coefficient structures had been fixed: balanced,
dispersion-aligned and dispersion-anti-aligned. The mapping was obtained from an independent
200-world technology ensemble using
\begin{equation}
D_{j,w}^{attr}
=\frac{1}{M_w}\sum_i\left|q_j(g_{ijw})-\bar q_{j,w}\right|,
\qquad
\bar D_j^{attr}=\E_W[D_{j,W}^{attr}].
\label{supp:eq:qmad}
\end{equation}
The complete ordering agreed between 150 and 200 worlds:
\begin{equation}
C_8>C_{10}>C_5>C_2>C_3>C_4>C_1>C_7>C_6>C_9.
\label{supp:eq:d3order}
\end{equation}

\begin{table}[t]
\centering
\caption{Frozen 200-world mean attribution-scale dispersion.}
\label{supp:tab:d3dispersion}
\small
\begin{tabular}{@{}lc@{}}
\toprule
Criterion & Mean dispersion\\
\midrule
$C_8$ & 0.122873493\\
$C_{10}$ & 0.119403289\\
$C_5$ & 0.110114266\\
$C_2$ & 0.094820283\\
$C_3$ & 0.090482741\\
$C_4$ & 0.088452971\\
$C_1$ & 0.087135074\\
$C_7$ & 0.078706704\\
$C_6$ & 0.049591336\\
$C_9$ & 0.028266887\\
\bottomrule
\end{tabular}
\end{table}

Balanced coefficients were $\alpha_j=0.10$. The aligned and anti-aligned mappings were
\begin{align}
\boldsymbol\alpha^{aligned}:
&(C_8,C_{10},C_5,C_2,C_3,C_4,C_1,C_7,C_6,C_9)
\nonumber\\
&\mapsto(0.16,0.14,0.13,0.12,0.11,0.10,0.08,0.07,0.05,0.04),
\label{supp:eq:alphaaligned}\\
\boldsymbol\alpha^{anti}:
&(C_9,C_6,C_7,C_1,C_4,C_3,C_2,C_5,C_{10},C_8)
\nonumber\\
&\mapsto(0.16,0.14,0.13,0.12,0.11,0.10,0.08,0.07,0.05,0.04).
\label{supp:eq:alphaanti}
\end{align}
The frozen primary heterogeneous assignment had weak rank association with the dispersion
ordering ($\rho_S=0.1273$, $\tau_b=0.0222$), so it was not changed. Although the D3 mapping
itself was a completed development artefact, the planned primary stress comparisons were never
executed.

\section{Archived Bootstrap and Robustness Plan}

At the reference condition, $B=200$ context-level bootstrap replications were to resample FIT
and WEIGHT contexts separately, refit the model, recompute attributions and weights, and
evaluate decisions on the fixed original external TEST contexts. Confidence intervals were to
be reported for criterion weights, total-variation weight error, Kendall's $\tau_b$, Top-1
accuracy and normalized oracle regret. No cross-context probability that a particular ITS was
ranked first was to be interpreted as a global deployment probability.

A 50-oracle ensemble was to leave the frozen D29 technology-response kernel unchanged while
varying oracle coefficient vectors, nonlinearities and semantically admissible interaction
pairs. Relative outcome noise was to be recalibrated to each oracle variant's signal standard
deviation. This ensemble would not replace the response generator and therefore would not
reintroduce the historical v2.1 separability defect through oracle-only perturbations.

Two distribution-shift analyses were specified: holding out the upper 20\% of demand pressure
and holding out the lower 20\% of infrastructure readiness. Secondary robustness analyses were
to include criterion removal, $\pm10/20/30\%$ SHAP-weight perturbations,
leave-one-alternative-out analysis and MOORA--TOPSIS comparison.

\section{Archived Statistical Analysis Plan}

Factor effects were to be summarized with repeated-seed regression or mixed-effects models,
emphasizing effect sizes and uncertainty intervals rather than binary significance alone.
Method comparisons were paired because every method would share the same benchmark instance.
The 30 primary seeds were to be treated as benchmark-instance clusters; alternative--context
rows were not to be interpreted as independent superpopulation replications.

Method availability was to be reported explicitly. For each factorial cell, undefined
Oracle/SHAP weight rates and recorded Equal-fallback rates were to accompany conditional
decision-fidelity summaries. Pairwise contrasts were to use only benchmark instances on which
both methods were defined, with coverage and fallback rates reported alongside them. Undefined
Oracle or SHAP vectors were never to be silently replaced by Equal weights.

Planned outputs included predictive-performance and Direct-XGBoost references, weight-recovery
tables, decision-fidelity summaries, factorial effect summaries, alpha-alignment stress results,
bootstrap intervals, validity-boundary figures and the prespecified stage-wise fidelity ladder.
None of these primary outputs exists under v1.

\section{Reproducibility and Content-Integrity Record}
\label{supp:sec:integrity}

Table~\ref{supp:tab:digests} gives non-identifying truncated content digests. The selected B100
background exactly matched the frozen TreeSHAP identity set.

\begin{table}[t]
\centering
\caption{Truncated, non-identifying content digests of the frozen artefacts underlying the
reported numbers.}
\label{supp:tab:digests}
\small
\begin{tabular}{@{}ll@{}}
\toprule
Artifact & Truncated SHA-256\\
\midrule
Oracle B100 result & \texttt{7f9864bd...439c9fb0}\\
TreeSHAP background identities & \texttt{3184c675...1646eda1}\\
Pilot-B result & \texttt{05072857...9f787a88}\\
Fixed-weight selection-map result & \texttt{a1f293b6...d7e20f4e}\\
\bottomrule
\end{tabular}
\end{table}

The full anonymous provenance package records, in order, the
threshold freeze, evaluator/executor freeze, scientifically invariant direct-file loader repair,
audited Pilot-B result freeze, immutable closure reference, post-closure diagnostic protocol
freeze, implementation freeze and independently audited result freeze. The worktree was clean
after each closure point, and the final complete test run reported 836 passes with three
pre-existing external plotting warnings.
\ifanon
Exact repository identifiers are intentionally omitted here to protect anonymous review.
\else
Exact repository identifiers accompany the public reproducibility archive.
\fi
These facts document software and artefact integrity; they are not statistical evidence.

\section{Prospective Stop and Permanent Closure}

Pilot B was governed by five independently sufficient hard stops: coverage failure, failure to
separate RandomWeights, failure to escape MajorityWinner, pooled oracle-winner dominance and
weight insensitivity. Passage required every flag to be false. In the frozen result, only the
modal-winner effective-tie flag was true, but the overall OR rule therefore barred every
analysis in this supplement.

The immutable version-control, digest and test record appears in
Section~\ref{supp:sec:integrity}. Pilot B v1 is closed and must not be rerun. Any later study
requires a new specification and version reference, explicit disclosure of the v1 failure and
development-only design work; it cannot alter the v1 thresholds, result or closure reference.
The fixed-weight selection-map diagnostic in Section~\ref{supp:sec:fixedselection} is such a
separately frozen descriptive layer; it neither reopened Pilot B nor authorized any protected
analysis.

\end{document}
"
% ============================================================
% ChkTeX: intentional en dashes, literal digest ellipses, label placement,
% acronym punctuation and LNCS name/bibliography conventions.
% chktex-file 8
% chktex-file 11
% chktex-file 13
% chktex-file 24
% chktex-file 36
% chktex-file 38

\documentclass[runningheads,orivec]{llncs}

\usepackage[T1]{fontenc}
\usepackage[utf8]{inputenc}
\usepackage{amsmath,amssymb,amsfonts}
\usepackage{booktabs}
\usepackage{graphicx}
\usepackage{tabularx}
\usepackage{hyperref}
\hypersetup{hidelinks}

\newcommand{\E}{\mathbb{E}}
\newcommand{\median}{\operatorname{median}}
\newcommand{\IQR}{\operatorname{IQR}}
\newcommand{\SD}{\operatorname{SD}}
\newcommand{\clip}{\operatorname{clip}}
\newcommand{\XGB}{XGBoost}

% ---- build-mode switch: anonymous by default -----------------
\providecommand{\CAMERAREADY}{0}
\newif\ifanon
\anontrue
\ifnum\CAMERAREADY=1 \anonfalse \fi
% --------------------------------------------------------------

\begin{document}
\emergencystretch=1.5em

\title{Supplementary Methods and Audit Archive for\texorpdfstring{\\}{ }
``From Predictive Importance to Decision Value: Evaluating XAI-Based Weights for
ITS Prioritisation''}
\titlerunning{Supplementary Benchmark-Adequacy Record}

\ifanon
  \author{Anonymous Author(s)}
  \authorrunning{Anonymous Author(s)}
  \institute{Anonymous Institution(s)}
\else
  \author{Mohamed Yassine Neifar\inst{1,2} \and
  Ahmed Frikha\inst{2} \and
  Nadir Farhi\inst{1}}
  \authorrunning{Neifar et al.}
  \institute{Universit\'e Gustave Eiffel, COSYS--GRETTIA,
  Champs-sur-Marne, F-77447 Marne-la-Vall\'ee Cedex 2, France
  \and
  Institut Sup\'erieur de Gestion Industrielle de Sfax (ISGIS),
  OLID, University of Sfax, Sfax 3038, Tunisia}
\fi

\maketitle

\section{Status and Interpretation Boundary}

This supplement reports executed development audits and preserves the detailed primary and
robustness analyses that had been specified before the prospectively frozen Pilot-B
benchmark-adequacy gate was executed. The latter analyses were conditional on all five Pilot-B
hard-stop flags being false. The modal-winner effective-tie stop fired, so the primary factorial
was not authorized. Primary seeds 11001--11030 and reserve seeds 30001--30005 were never
accessed; all sections explicitly labelled ``Archived'' remain unexecuted protocols rather than
results.

The closed v1 state is anchored by an immutable version-control reference.
\ifanon
Exact repository identifiers and complete digests are withheld from this anonymous reviewer
file; they are available confidentially to editors and will accompany the public
reproducibility archive after review.
\else
Exact repository identifiers and complete digests accompany the public reproducibility archive.
\fi
This document is therefore a methods, audit-results and archival-protocol supplement. It is not
authorization to execute the protected design.
After permanent Pilot-B closure, one separately frozen development-only descriptive diagnostic
reconstructed the context-level MOORA selection maps of the seven observed archived fixed
vectors. It used no primary/reserve world, generated no target outcome, fitted or queried no
model, re-estimated no weight and could not alter the gate. Its protocol, implementation,
independent audit and result were frozen as a distinct post-closure layer.
Methodological citations are provided in the main manuscript and are not duplicated here.

\subsection*{Question Structure and Evidential Status}

The main manuscript states two prospective research questions and one separate post-closure
descriptive question. RQ1 asks whether structural diagnostics can reveal that apparent
contextual variation disappears under the normalization used by the downstream MCDM operator.
RQ2 asks whether, after that structural defect is corrected, structured global-weight pipelines
provide sufficient incremental decision value over random-weight and constant-decision
references to justify protected primary evaluation. Both are addressed by the prospectively
structured benchmark-development record and the frozen Pilot-B result reported here.

DQ1 asks what the archived context-level selection maps reveal about the observed decision
behaviour of the seven frozen fixed-weight vectors. DQ1 is post-closure and descriptive. It is
non-confirmatory, it carries no pass/fail role, and it cannot modify, reinterpret or reopen the
Pilot-B gate. Section~\ref{supp:sec:fixedselection} reports it under that label.

Throughout this supplement, \emph{each replication seed defines one benchmark world}, and the
world is the statistical unit. Where a numbered identifier such as 21001 appears it names a
world; the phrase ``replication seed'' is retained only when the frozen artefact field or
protocol clause is itself named that way. \emph{Validity} denotes conceptual scientific
properties of the benchmark; \emph{adequacy} denotes the operational prospectively frozen gate
used to decide whether protected primary execution was warranted.

\subsection*{Reader Map}

The conference-facing main manuscript follows a problem--method--result--meaning structure and
deliberately omits most development chronology. Table~\ref{supp:tab:readermap} directs readers
to the corresponding detailed evidence. This separation changes presentation only: the frozen
scientific record, Pilot-B decision and post-closure interpretation boundary are unchanged.

\begin{table}[t]
\centering
\caption{Navigation from the concise main argument to the detailed reproducibility record,
with the evidential status of each block. Status codes are: EDA, executed development audit;
PBA, Pilot-B confirmatory adequacy result; PCD, post-closure descriptive diagnostic; ANE,
archived and not executed.}
\label{supp:tab:readermap}
\small
\begin{tabularx}{\linewidth}{@{}lXc@{}}
\toprule
\textbf{Main-paper element} & \textbf{Detailed supplementary evidence} & \textbf{Status}\\
\midrule
Controlled ITS laboratory & Criterion definitions, capability pathways, copula contexts and
opportunity functions in Sect.~2. & EDA\\
Oracle and global attribution & Closed-form functional, joint-background moments, \XGB,
TreeSHAP and B100 checks in Sect.~3. & EDA\\
Structural correction & Calibration, positive control, numerical null and candidate screening
in Sect.~4. & EDA\\
D29 decision geometry & World-resolved oracle orderings, winners, modal shares and regret in
Sect.~5. & EDA\\
Pilot-B evidence (RQ2) & Per-world method contrasts, gate decision and objective-weight vectors
in Sects.~5 and 7. & PBA\\
Observed selection maps (DQ1) & All audited modal identities, shares, unique-choice counts and
pairwise agreement summaries in Sect.~5. & PCD\\
Decision operators and gate & Exact weighting, tie, regret and hard-stop rules in Sects.~6--7.
& EDA\\
Unexecuted primary plan & Clearly labelled archival protocols in Sects.~8--12. & ANE\\
Integrity and closure & Content digests, software tests and permanent-stop record in
Sects.~13--14. & EDA\\
\bottomrule
\end{tabularx}
\end{table}

\section{Detailed Benchmark Construction}

\subsection{Criteria and Capability Pathways}

The six alternatives in the main manuscript were evaluated on the ten technology-agnostic
criteria in Table~\ref{supp:tab:criteria}. The original lifecycle-burden criterion was a cost;
all other criteria were benefits.

\begin{table}[t]
\centering
\caption{Controlled evaluation criteria.}
\label{supp:tab:criteria}
\small
\begin{tabularx}{\linewidth}{@{}c X c@{}}
\toprule
Code & Criterion & Original direction\\
\midrule
$C_1$ & Mobility-efficiency improvement & Benefit\\
$C_2$ & Travel-time reliability improvement & Benefit\\
$C_3$ & Person-throughput improvement & Benefit\\
$C_4$ & Safety-surrogate improvement & Benefit\\
$C_5$ & Non-recurring-congestion resilience & Benefit\\
$C_6$ & Environmental-efficiency improvement & Benefit\\
$C_7$ & User and accessibility benefit & Benefit\\
$C_8$ & Infrastructure and data readiness & Benefit\\
$C_9$ & Interoperability and scalability & Benefit\\
$C_{10}$ & Annualised lifecycle deployment burden & Cost\\
\bottomrule
\end{tabularx}
\end{table}

The pre-specified capability mask in Table~\ref{supp:tab:mask} defines synthetic functional
pathways, not empirical impact magnitudes. Direct (D), indirect (I) and zero (0) pathways set
broad response-amplitude bands.

\begin{table}[t]
\centering
\caption{Pre-specified structural capability mask.}
\label{supp:tab:mask}
\scriptsize
\begin{tabular}{@{}lccccccc ccc@{}}
\toprule
 & $C_1$ & $C_2$ & $C_3$ & $C_4$ & $C_5$ & $C_6$ & $C_7$ & Ready & Interop. & Burden\\
\midrule
ATSC      & D & D & D & I & I & I & I & M    & H & M\\
TSP       & I & D & D & I & 0 & I & D & M    & H & L--M\\
TIDM      & I & D & I & D & D & I & I & M    & H & M\\
RMTI--DRG & I & I & I & I & D & I & D & L--M & H & L--M\\
V2X--CS   & I & I & 0 & D & I & I & I & H    & H & H\\
DSM       & D & D & I & D & D & I & I & M--H & H & M--H\\
\bottomrule
\end{tabular}
\end{table}

\subsection{Gaussian-Copula Contexts and Opportunity Functions}

For demand, incident, person-movement, environmental and readiness factors, generate
$z_{0,s},z_{1,s},\ldots,z_{5,s}\overset{iid}{\sim}\mathcal N(0,1)$ and set
\begin{equation}
\widetilde h_{k,s}^{(\rho)}
=\sqrt{\rho}\,z_{0,s}+\sqrt{1-\rho}\,z_{k,s},
\qquad h_{k,s}^{(\rho)}=\Phi(\widetilde h_{k,s}^{(\rho)}).
\label{supp:eq:copula}
\end{equation}
The same base draws were reused across dependence levels, while estimation and external TEST
streams used disjoint namespaces. The first seven criterion opportunities were
\begin{align}
o_{s1}&=h_D, & o_{s2}&=0.60h_D+0.40h_I, & o_{s3}&=0.60h_D+0.40h_T,\nonumber\\
o_{s4}&=0.50h_D+0.50h_I, & o_{s5}&=h_I, & o_{s6}&=0.60h_D+0.40h_E,\nonumber\\
o_{s7}&=0.40h_D+0.60h_T.&&&
\label{supp:eq:opportunities}
\end{align}

Each replication seed $r$ defined one benchmark world $W_r$: technology parameters, a master
context stream and all common-random-number streams. The conditional method target had been
\begin{equation}
\Psi_m(x)=\E_W\!\left[\E_{H,\varepsilon}
\{L_m(x;W,H,\varepsilon)\mid W\}\right].
\label{supp:eq:estimand}
\end{equation}
The world, rather than an alternative--context row, was the intended statistical unit. The
primary worlds needed to estimate this quantity were never accessed after the gate fired.

\subsection{Historical and Frozen Response Layers}

For $C_1$--$C_7$, capability amplitudes were drawn once per world:
\begin{equation}
\theta_{aj}\sim
\begin{cases}
0,&M_{aj}=0,\\
\mathcal U(0.05,0.20),&M_{aj}=I,\\
\mathcal U(0.20,0.40),&M_{aj}=D.
\end{cases}
\label{supp:eq:effectbands}
\end{equation}
The historical v2.1 response was
\begin{equation}
x_{asj}=\clip\!\left(\theta_{aj}o_{sj}+\sigma_x\xi_{asj},0,1\right).
\label{supp:eq:v21}
\end{equation}
The deployment classes were
\begin{equation}
L\sim U(0.25,0.40),\qquad
M\sim U(0.45,0.60),\qquad
H\sim U(0.65,0.80).
\label{supp:eq:deploymentclasses}
\end{equation}
Compound classes were equal-probability mixtures. With readiness
requirement $r_a$, interoperability $\iota_a$ and burden $\kappa_a$,
\begin{align}
x_{as8}&=\clip\!\left(1-\max(0,r_a-h_{R,s})+\eta_{as8},0,1\right),\nonumber\\
x_{as9}&=\clip\!\left(\iota_a+0.10h_{R,s}+\eta_{as9},0,1\right),\nonumber\\
x_{as,10}&=\clip\!\left(\kappa_a+0.20(1-h_{R,s})r_a+\eta_{as,10},0,1\right).
\label{supp:eq:c8c10}
\end{align}
Benefit orientation used $g_{asj}=x_{asj}$ for $j\le9$ and
$g_{as,10}=1-x_{as,10}$. At $\sigma_x=0$, the first seven criteria reduced to
$g_{asj}=\theta_{aj}o_{sj}$, which explains the exact normalization cancellation in the main
paper.

The frozen D29 kernel was
\begin{equation}
g_{asj}=\theta_{aj}o_{sj}+0.05v_{aj}o_{sj}(1-o_{sj}),\qquad j\le7.
\label{supp:eq:d29}
\end{equation}
For $m_j$ active alternatives, its ordered contrast grid was
$v^{grid}_{\ell j}=-1+2(\ell-1)/(m_j-1)$. After sorting active IDs, one deterministic NumPy
permutation used
\begin{equation}
\texttt{SeedSequence([r, 2010, j])}.
\label{supp:eq:d29seedsequence}
\end{equation}
No outcome or
winner information entered this assignment. Active $\theta_{aj}\ge0.05$ ensured
$\partial g/\partial o=\theta_{aj}+0.05v_{aj}(1-2o)\ge0$. Structural-zero pathways remained
zero, and $C_8$--$C_{10}$ retained Eq.~\eqref{supp:eq:c8c10}.

For $o_{sj}>0$, within-context vector normalization over active alternatives gives the exact
geometry
\begin{equation}
r_{asj}=
\frac{\theta_{aj}+0.05v_{aj}(1-o_{sj})}
{\sqrt{\sum_{b\in\mathcal A_j}
[\theta_{bj}+0.05v_{bj}(1-o_{sj})]^2}}.
\label{supp:eq:d29normalized}
\end{equation}
Thus D29 removes the v2.1 cancellation, but each criterion-specific normalized profile remains
on a one-parameter curve indexed by $o_{sj}$. Because opportunities differ across criteria,
Eq.~\eqref{supp:eq:d29normalized} alone does not prove that a fixed-weight MOORA winner is
context invariant. At Pilot-B closure, fixed-weight invariance was therefore an untested
hypothesis rather than an explanation of the stop. The subsequently frozen descriptive
selection-map diagnostic reported below resolves this question for the seven observed vectors
only; it does not establish a causal effect of Eq.~\eqref{supp:eq:d29normalized} or characterize
unobserved weight vectors.

\section{Oracle, Learning and Background Details}

\subsection{Oracle Utility and Closed-Form Attribution}

The monotone transform, coefficient vector and interaction graph were
\begin{align}
q_j(g)&=\frac{\log(1+2g)}{\log3},\nonumber\\
\boldsymbol\alpha&=(0.11,0.04,0.05,0.07,0.10,0.14,0.12,0.16,0.08,0.13),\nonumber\\
\mathcal E&=\{(1,2),(3,7),(4,5),(6,10),(8,9)\},\qquad \beta_{jk}=0.20.
\label{supp:eq:oraclecomponents}
\end{align}
For $z_j=q_j(g_j)$, the empirical FIT-background moments were
$\mu_j=\E_{bg}(z_j)$ and $\mu_{jk}=\E_{bg}(z_jz_k)$; dependence was retained by not replacing
$\mu_{jk}$ with $\mu_j\mu_k$. The main-effect contribution was
$\phi_j^{main}=\alpha_j(z_j-\mu_j)$. For pair $\beta_{jk}z_jz_k$, feature $j$ received
\begin{equation}
\phi_j^{(jk)}=\frac{\beta_{jk}}{2}
\{z_j\mu_k-\mu_{jk}+z_jz_k-\mu_jz_k\},
\label{supp:eq:pairshap}
\end{equation}
with the symmetric expression for $k$. The closed-form attribution was
\begin{equation}
\phi_j^\star(\mathbf g)=\frac{1}{1+\lambda}
\left[\alpha_j(z_j-\mu_j)+\lambda\!\sum_{k:(j,k)\in\mathcal E}\phi_j^{(jk)}\right],
\label{supp:eq:closedshap}
\end{equation}
with symmetric handling when $j$ was second in a pair. Exhaustive coalition enumeration was a
unit-test reference. Global weights normalized mean absolute WEIGHT-row attributions; total
importance at most $10^{-12}$ made the vector undefined, with no Equal substitution.

\subsection{XGBoost and TreeSHAP Freeze}

XGBoost was selected once by context-grouped cross-validation on development FIT rows at the
reference condition. Candidate 45 attained mean RMSE $0.017757058913365983$ with
\begin{equation}
\begin{split}
&n_{\mathrm{estimators}}=600,\quad d_{\max}=2,\quad \eta=0.05,\\
&\text{subsample}=0.8,\quad \text{colsample\_bytree}=0.8,\quad
\text{min\_child\_weight}=5,\\
&\text{reg\_alpha}=0,\quad \text{reg\_lambda}=1.
\end{split}
\label{supp:eq:xgbparams}
\end{equation}
The objective was \texttt{reg:squarederror}, with one computational thread and estimator
random-state namespace 82002. The selected model was then fixed. Interventional TreeSHAP used 100 deterministic FIT
identities selected with row-priority namespace 83001 and evaluated mean absolute attributions
on WEIGHT. External TEST did not enter calibration, background selection or weighting.

\subsection{Matched B100 Oracle Diagnostic}

The primary oracle used all eligible FIT rows for $\mu_j$ and $\mu_{jk}$. The secondary B100
diagnostic changed only these moments to the exact 100 frozen TreeSHAP identities; WEIGHT rows,
the oracle functional and every other setting were unchanged. It compared full-FIT and B100
weight vectors through mean absolute difference, total variation, maximum absolute difference,
Spearman correlation, Top-3/Top-5 overlap and maximum marginal/joint moment changes. Undefined
vectors remained undefined. The only authorized pre-primary execution was seed 21001 at
$N=250$, $\rho=0.4$ and $\lambda=0.5$; no $Y$, SHAP package call or MCDM operation was involved.

\section{Detailed Structural Calibration and Candidate Screening}

\subsection{Structural and Reachability Metrics}

Let $G_j\in\mathbb R^{S\times A_j}$ denote the active-alternative response matrix for criterion
$j$ at zero measurement noise. If its singular values are
$\sigma_{1j}\ge\sigma_{2j}\ge\cdots$ and
$E_j=\sum_\ell\sigma_{\ell j}^2$, the rank-one residual was
\begin{equation}
NS_j=1-\frac{\sigma_{1j}^2}{E_j},
\qquad E_j>10^{-24}.
\label{supp:eq:NS}
\end{equation}
No epsilon was inserted into this ratio; structural-zero matrices were reported as undefined.
For each jointly positive active alternative pair,
\begin{equation}
LRV_{ab,j}=\SD_s\!\left[\log(g_{asj}/g_{bsj})\right],
\qquad
LRV50_j=\median_{a<b}LRV_{ab,j},
\label{supp:eq:LRV}
\end{equation}
using population SD. For normalization family $m$,
\begin{equation}
NSV_{j,m}=\left\{S^{-1}\sum_s
\left\|\mathbf r_{sj}^{(m)}-\bar{\mathbf r}_j^{(m)}\right\|_2^2\right\}^{1/2}.
\label{supp:eq:NSV}
\end{equation}
Vector normalization was the primary structural gate; sum, max and min--max variants were
robustness diagnostics.

Reachability reflected one idiosyncratic Gaussian-copula shock while holding the shared and all
other idiosyncratic shocks fixed. With opportunity functions recomputed after reflection,
\begin{equation}
RE_{kj}=\median_{s,a\in\mathcal A_j}|g_{asj}^{(-k)}-g_{asj}|,
\qquad
SRE_{kj}=\frac{RE_{kj}}{\IQR(g_j)+10^{-12}}.
\label{supp:eq:SRE}
\end{equation}

\subsection{Scientific Positive Control}

For active $C_1$--$C_7$ pathways, the dedicated positive control was
\begin{equation}
g_{asj}^{PC,\ell}=\left[\theta_{aj}+0.05c_a^{(\ell)}(2o_{sj}-1)\right]o_{sj},
\label{supp:eq:positivecontrol}
\end{equation}
with contrast base $(-1,-0.6,-0.2,0.2,0.6,1)$ and all six cyclic rotations. Structural-zero
pathways remained zero and $C_8$--$C_{10}$ were unchanged. Rotation medians were formed within
development worlds 21001--21005 and world medians were then used as scientific references.

\begin{table}[t]
\centering
\caption{Frozen D2.7 criterion-level scientific references.}
\label{supp:tab:d27layera}
\small
\begin{tabular}{@{}ccc@{}}
\toprule
Criterion & $T^{sci}_{LRV,j}$ & $T^{sci}_{NSV,j}$\\
\midrule
$C_1$ & 0.185606 & 0.101685\\
$C_2$ & 0.107107 & 0.069436\\
$C_3$ & 0.163029 & 0.064900\\
$C_4$ & 0.115796 & 0.069423\\
$C_5$ & 0.127964 & 0.073681\\
$C_6$ & 0.194899 & 0.129935\\
$C_7$ & 0.144080 & 0.073092\\
\bottomrule
\end{tabular}
\end{table}

\begin{table}[t]
\centering
\caption{Frozen D2.7 latent-pathway reachability references.}
\label{supp:tab:d27sre}
\small
\begin{tabular}{@{}lc@{}}
\toprule
Pathway & $T^{sci}_{SRE,kj}$\\
\midrule
$h_D\rightarrow C_1$ & 0.556245\\
$h_D\rightarrow C_2$ & 0.334437\\
$h_D\rightarrow C_3$ & 0.284878\\
$h_D\rightarrow C_4$ & 0.292844\\
$h_D\rightarrow C_6$ & 0.431612\\
$h_D\rightarrow C_7$ & 0.252576\\
$h_E\rightarrow C_6$ & 0.296869\\
$h_I\rightarrow C_2$ & 0.212043\\
$h_I\rightarrow C_4$ & 0.294905\\
$h_I\rightarrow C_5$ & 0.455271\\
$h_T\rightarrow C_3$ & 0.186995\\
$h_T\rightarrow C_7$ & 0.373380\\
\bottomrule
\end{tabular}
\end{table}

The paired positive-control uplift in mean normalized constant-modal regret was
$(0.001485,0.020022,-0.013836,0.015534,0.006161)$ for worlds 21001--21005. Because the activation
rule required positive uplift in all five worlds, no Layer-B scientific regret reference was
defined and no post-hoc value was introduced.

\subsection{Numerical Null and Candidate Outcomes}

Numerical thresholds were calibrated from 2000 exact rank-one replicates for each of five- and
six-alternative active sets:
\begin{equation}
T^{num}_{NS}=10^{-12},\qquad
T^{num}_{LRV}=10^{-10},\qquad
T^{num}_{NSV}=3.2506084454\times10^{-8}.
\label{supp:eq:nullthresholds}
\end{equation}
All 35 historical v2.1 world--criterion rows fell below every threshold. For active pathways
\(j\leq7\), the D4-F1 curvature family was
\(g_{asj}^{D4F1}(\chi)=\theta_{aj}o_{sj}[1+\chi v_{aj}(1-o_{sj})]\), with structural-zero
pathways fixed at zero. It cleared some structural references at higher curvature but lost reachability; no
$\chi\in\{0.1,\ldots,1.0\}$ passed all gates, so the family was rejected.

The later exact D2.9 audit evaluated all $6^5=7776$ prespecified positive-control profile
assignments. None met the old simultaneous binary conjunction. This established non-viability
for that calibrated design, not mathematical impossibility. Under the corrected prospective
rule, 54 candidates were eligible. D29 was retained solely for protocol continuity and minimal
intervention, not scientific preference, downstream performance, SHAP pattern or ITS identity.

The count of 54 is composed rather than directly observed in a single artefact, and is
decomposed here to prevent confusion. The nonselective retrospective re-adjudication examined 54
historical parameter points and returned 44 \texttt{VERIFIED\_\allowbreak ELIGIBLE},
0 \texttt{VERIFIED\_\allowbreak INELIGIBLE} and 10
\texttt{INSUFFICIENT\_\allowbreak FROZEN\_\allowbreak EVIDENCE}. The 10
unresolved points were all v2.2 points whose frozen artefact schema did not record C8--C10
invariance. A separate supplemental resolution then established that invariance from the
contemporaneous frozen source contract, and all 10 satisfied the corrected hard constraints.
The frozen retention artefact therefore records an eligible universe of
$44+10=54$ specifications. Eligibility is explicitly non-unique: the retention of D29 is a
minimal-intervention and provenance-stability decision, not a superiority, uniqueness or
statistical-selection claim, and no ranking or re-optimization over the 54 points was
performed.

\section{Detailed Development Audit Results}

\subsection{D29 Layer-B Decision Geometry}

The one-shot characterization used development worlds 21001--21005, the complete $N=1000$
FIT+WEIGHT pool, $\rho=0.4$, $\sigma_x=0$ and $\lambda=0.5$. It generated no $Y$ and executed
no learned model, SHAP package or MCDM method. Regret used the exact context-specific oracle
range and was undefined only if that range was at most $10^{-12}$; no context was undefined.

\begin{table}[t]
\centering
\caption{Complete frozen D29 Layer-B development record.}
\label{supp:tab:d29layerb}
\scriptsize
\resizebox{\linewidth}{!}{%
\begin{tabular}{@{}r r r l r r r r@{}}
\toprule
World & Orderings & Winners & Modal & Modal share & Mean regret & Median regret & $P_{95}$ regret\\
\midrule
21001 & 23 & 2 & $A_6$ & 0.983 & 0.00042018 & 0 & 0\\
21002 & 78 & 5 & $A_1$ & 0.519 & 0.05992337 & 0 & 0.29155475\\
21003 & 24 & 4 & $A_2$ & 0.745 & 0.04523642 & 0 & 0.28747533\\
21004 & 17 & 3 & $A_2$ & 0.827 & 0.01734351 & 0 & 0.12641643\\
21005 & 38 & 4 & $A_2$ & 0.809 & 0.01606539 & 0 & 0.12654230\\
\bottomrule
\end{tabular}}
\end{table}

Across worlds, the medians were 24 orderings, four winners, modal share $0.809$, mean regret
$0.01734351$, median regret zero and $P_{95}$ regret $0.12654230$. At the directly comparable
world 21001, the historical v2.1 record had 19 orderings, two winners, modal share $0.962$ and
mean always-modal regret $0.0011106$ under its frozen $D_s+10^{-12}$ denominator. D29 had 23
orderings, two winners, modal share $0.983$ and mean regret $0.00042018$ under exact $D_s$.
The comparison therefore supports neither universal degeneration nor universal improvement.

\subsection{XGBoost and TreeSHAP Development Checks}

The selected XGBoost hyperparameters are reported in Eq.~\eqref{supp:eq:xgbparams}; candidate
number and grouped-CV RMSE are given above. At world 21001 and the reference condition, total mean-absolute
TreeSHAP importance was $0.07006131011183041$ and maximum local-accuracy error was
$6.45\times10^{-7}$. The development-only weight vector for $(C_1,\ldots,C_{10})$ was
\begin{equation}
\begin{split}
&(0.124860,\ 0.113194,\ 0.069426,\ 0.025945,\ 0.155169,\\
&\qquad 0.078339,\ 0.071621,\ 0.265059,\ 0.020729,\ 0.075658).
\end{split}
\label{supp:eq:devshapweights}
\end{equation}
These observations did not tune comparator definitions, MCDM rules or winner semantics.

\subsection{B100-Matched Background Result}

Both oracle vectors were defined, and the 100 identities exactly matched the frozen TreeSHAP
membership. Total mean-absolute importance was $0.06961866590$ under full-FIT moments and
$0.06992189186$ under B100 moments.

\begin{table}[t]
\centering
\caption{Development-only oracle-attribution weights under full-FIT and matched B100 moments.}
\label{supp:tab:b100weights}
\small
\begin{tabular}{@{}lrrr@{}}
\toprule
Criterion & Full-FIT & B100 & Absolute difference\\
\midrule
$C_1$  & 0.09663438 & 0.09712037 & 0.00048600\\
$C_2$  & 0.04849942 & 0.04854063 & 0.00004122\\
$C_3$  & 0.05910838 & 0.05913637 & 0.00002799\\
$C_4$  & 0.06665455 & 0.06559847 & 0.00105608\\
$C_5$  & 0.14880967 & 0.14667297 & 0.00213669\\
$C_6$  & 0.09368358 & 0.09388953 & 0.00020596\\
$C_7$  & 0.07615062 & 0.07594850 & 0.00020212\\
$C_8$  & 0.24766693 & 0.25232065 & 0.00465372\\
$C_9$  & 0.04819959 & 0.04935317 & 0.00115358\\
$C_{10}$ & 0.11459290 & 0.11141932 & 0.00317357\\
\bottomrule
\end{tabular}
\end{table}

The weight-vector discrepancies were
\begin{equation}
\operatorname{MAE}_w=0.0013136928,\quad
TV_w=0.0065684640,\quad
\max_j|\Delta w_j|=0.0046537223.
\label{supp:eq:b100distances}
\end{equation}
Spearman correlation was $0.98787879$, and Top-3 and Top-5 overlaps were both $1.0$.
Maximum absolute marginal- and joint-moment changes were $0.0098907938$ and $0.0041811926$,
respectively. Equal fallback was not used. These results show small
finite-background sensitivity at this one development reference, not uniform robustness.

\subsection{World-Resolved Pilot-B Oracle and SHAP Comparison}

Table~\ref{supp:tab:pilotboracleshap} retains the full comparison moved from the main paper.
Each entry is an advantage over the world-specific MajorityWinner reference; positive values
favour the structured method.

\begin{table}[t]
\centering
\caption{Per-world closed-form oracle-attribution and SHAP advantages over MajorityWinner.}
\label{supp:tab:pilotboracleshap}
\small
\begin{tabular}{@{}crrrr@{}}
\toprule
World & Oracle $\Delta$Top-1 & SHAP $\Delta$Top-1 & Oracle $\Delta$regret & SHAP $\Delta$regret\\
\midrule
21001 & $ 0.000$ & $ 0.000$ & $ 0.000000$ & $ 0.000000$\\
21002 & $ 0.035$ & $ 0.030$ & $ 0.001538$ & $ 0.003064$\\
21003 & $-0.465$ & $ 0.020$ & $-0.143522$ & $ 0.006717$\\
21004 & $-0.010$ & $ 0.000$ & $-0.002119$ & $ 0.000000$\\
21005 & $-0.380$ & $-0.035$ & $-0.090328$ & $-0.002986$\\
\bottomrule
\end{tabular}
\end{table}

The closed-form method exactly implements the frozen global interventional attribution
functional; it is not an exact representation of context-specific oracle utility. Thus its
failure to escape MajorityWinner is evidence about the frozen global-compression--MOORA
pipeline, not a numerical error in the attribution formula.

\subsection{Illustrative Frozen Objective-Weight Vectors}

Table~\ref{supp:tab:objectivevectors} directly transcribes the CRITIC and Entropy vectors stored
for development world 21001; no new transformation of TEST decisions was performed. The display
cannot alter the frozen gate or establish a performance mechanism. Both vectors were defined
without Equal fallback, and neither has its largest component on C8 or C10. Thus concentration
on those two criteria is not supported in this world. The table is deliberately not combined
with five-world median performance to infer cross-world geometry.

\begin{table}[t]
\centering
\caption{Frozen objective-weight vectors for development world 21001.}
\label{supp:tab:objectivevectors}
\scriptsize
\setlength{\tabcolsep}{2.8pt}
\begin{tabular}{@{}lrrrrrrrrrr@{}}
\toprule
Method & C1 & C2 & C3 & C4 & C5 & C6 & C7 & C8 & C9 & C10\\
\midrule
CRITIC & .079 & .075 & .101 & .093 & .150 & .099 & .097 & .111 & .084 & .111\\
Entropy & .131 & .114 & .175 & .099 & .236 & .064 & .084 & .027 & .032 & .038\\
\bottomrule
\end{tabular}
\end{table}

The aggregate coincidence in world 21003 did not, by itself, identify a common decision map.
Equal, CRITIC and Entropy each recorded Top-1 accuracy $0.190$ and mean normalized regret
$0.316758$, but their mean Kendall $\tau_b$ values were respectively $0.5493$, $0.5427$ and
$0.5113$. Top-1 accuracy and normalized regret depend on the selected alternative, whereas
Kendall $\tau_b$ depends on the complete ranking. The $\tau_b$ differences therefore establish
different lower-order ranking behaviour on average but are logically uninformative about
argmax equality. That ambiguity motivated the separately frozen diagnostic below.

\subsection{Frozen Post-Closure Fixed-Weight Selection-Map Diagnostic (DQ1)}
\label{supp:sec:fixedselection}

\noindent\textbf{Evidential status: post-closure, descriptive, non-confirmatory.} This
subsection answers DQ1 only. It cannot modify, reinterpret, rescue or reopen the Pilot-B gate,
it defines no pass/fail criterion, and none of its content is confirmatory evidence for RQ1 or
RQ2.

The diagnostic was specified only after the Pilot-B result had been independently audited,
frozen and permanently closed. Its protocol was committed before any context-level
fixed-weight result was reconstructed. It used worlds 21001--21005 at the same reference
condition and the 200 frozen external TEST identities per world. The six structured vectors
were read from the immutable Pilot-B result without re-estimation; Equal was exactly
$(0.1,\ldots,0.1)$. D29 TEST decision matrices and oracle utilities were reconstructed only to
validate the archived aggregates. No target $Y$, model fit/prediction, SHAP, PI, RandomWeights,
bootstrap or protected seed was executed.

For each world, method and context the result stores the selected alternative; all six MOORA
scores and ranks; oracle utilities, ranks and selected alternative; Top-1 match; normalized
regret contribution; and agreement with the frozen MajorityWinner identity. Exactly
$5\times7\times200=7000$ records were independently audited. The audit reconstructed every
selection, anchor-based rank and regret directly from the stored full-precision vectors,
matched all 35 archived method--world aggregates within $10^{-12}$, and rebuilt all 105
pairwise selection-agreement summaries without executing the generator.

Table~\ref{supp:tab:fixedselectionmaps} reports each method's modal selected alternative,
modal share and number of distinct selected alternatives. Entries use the format
``alternative/\allowbreak share/\allowbreak $k$'', where $k=1$ denotes
context-invariant selection in that world.
MajorityWinner (MW) is a single frozen alternative by construction.

\begin{table}[t]
\centering
\caption{Audited selection maps for the seven observed fixed vectors. Each method entry is
modal alternative/modal share/number of distinct selected alternatives over 200 TEST contexts.}
\label{supp:tab:fixedselectionmaps}
\scriptsize
\setlength{\tabcolsep}{2.1pt}
\begin{tabular}{@{}rclllllll@{}}
\toprule
World & MW & Oracle & SHAP & PI & Ridge+ & CRITIC & Entropy & Equal\\
\midrule
21001 & A6 & A6/1.000/1 & A6/1.000/1 & A6/.980/2 & A6/1.000/1 & A6/1.000/1 & A6/1.000/1 & A6/1.000/1\\
21002 & A1 & A1/.700/3 & A1/.950/2 & A1/.630/2 & A1/.875/3 & A3/.800/3 & A1/1.000/1 & A1/1.000/1\\
21003 & A2 & A6/.720/4 & A2/.690/3 & A2/.530/3 & A6/.495/3 & A6/1.000/1 & A6/1.000/1 & A6/1.000/1\\
21004 & A2 & A2/.990/2 & A2/1.000/1 & A2/.775/4 & A2/.990/2 & A3/.750/3 & A3/.985/3 & A2/1.000/1\\
21005 & A2 & A6/.350/4 & A2/.790/3 & A3/.675/5 & A6/.375/4 & A6/.915/3 & A6/.800/3 & A2/1.000/1\\
\bottomrule
\end{tabular}
\end{table}

Equal was context invariant in all five worlds and exactly matched MajorityWinner across all
200 contexts in worlds 21001, 21002, 21004 and 21005. In world 21003, Equal instead selected A6
throughout while MajorityWinner was A2. The numbers of constant-selection worlds were
Oracle 1, SHAP 2, PI 0, Ridge+ 1, CRITIC 2, Entropy 3 and Equal 5; exact all-context
MajorityWinner agreement occurred in 1, 2, 0, 1, 1, 2 and 4 worlds, respectively.

The world-21003 ambiguity is now directly resolved: Equal, CRITIC and Entropy all selected A6
in every TEST context. Their different mean $\tau_b$ values therefore reflect differences below
the argmax, not different selected alternatives. More broadly, the heterogeneous constancy
counts show that fixed-vector MOORA was neither universally context invariant nor uniformly
context responsive over the seven observed vectors. The result supports coarse
observed-vector decision regions as a descriptive hypothesis, but cannot establish the causal
mechanism of the Pilot-B stop, characterize unobserved regions of the weight simplex, rank
methods or generalize beyond these five development worlds.

The new artefact makes paired context resampling technically possible, but the prospectively
frozen diagnostic explicitly excluded bootstrap and defined no inferential or pass/fail role.
No resampling was executed. Any future paired sensitivity analysis requires its own protocol,
frozen before inspecting bootstrap output, and cannot revise the closed Pilot-B decision.

\section{Detailed Weighting and Decision Operators}

The primary information partitions were fixed: Oracle, SHAP, PI, CRITIC and Entropy used WEIGHT;
Ridge+ used FIT; no method used external TEST for fitting or weight estimation. Oracle and SHAP
normalized mean absolute attributions. SHAP was declared undefined if total importance was at
most $10^{-12}$ and was never replaced by Equal weights.

CRITIC applied criterion-wise min--max scaling, population SD and Pearson conflict over
nonconstant columns. Entropy applied the same min--max scaling followed by Shannon
diversification with $0\log0=0$. Ridge+ standardized predictors and target with fold-local
population moments, solved non-negative L2-augmented least squares without an intercept, and
used the frozen grid
\begin{equation}
\tau\in\{10^{-4},10^{-3},10^{-2},10^{-1},1,10,100\}.
\label{supp:eq:ridgegrid}
\end{equation}
Five-fold context-grouped original-scale RMSE selected $\tau$, with the smallest value as tie
break.

Permutation importance kept the fitted XGBoost model fixed. For each criterion, it copied the
complete six-alternative block from source context $\pi(d)$ to destination context $d$ under 20
unique derangements. Their frozen generator was
\begin{equation}
\texttt{SeedSequence([replication, 84001, N])}.
\label{supp:eq:piseedsequence}
\end{equation}
The same
derangements were reused across criteria and factor levels for a given seed and $N$. Importance
was permuted minus baseline WEIGHT MSE; negative means were truncated before normalization and
repeat dispersion used sample SD. CRITIC, Entropy, Ridge+ and PI used Equal weights only under
their frozen zero-information rules, with fallback recorded.

Equal was $(0.1,\ldots,0.1)$. RandomWeights used 200
$\operatorname{Dirichlet}(1,\ldots,1)$ vectors per world from namespace 81001 and reused them
across conditions. MajorityWinner estimated the oracle modal winner on WEIGHT and applied that
constant choice to TEST. MajorityWinner was intentionally oracle-informed: it was the
context-free oracle-modal alternative estimated on WEIGHT and evaluated on TEST, and it was
optimal only for empirical oracle Top-1 frequency among constant choices. It is not claimed to
minimize regret or to constitute an operational transportation policy. It was neither a
deployable competitor nor a universally optimal choice, and it was not stakeholder-derived or
normative. The regret arm of the frozen gate uses this same
prospectively frozen Top-1-derived constant reference rather than a separately optimized
regret-minimizing constant policy; a future benchmark could prospectively define metric-matched
constant references. The frozen gate itself is unchanged by this observation. Direct-XGBoost was a parallel
predictive-decision reference rather than a weighting
method. MOORA used within-context vector normalization and a weighted sum; TOPSIS used the same
normalized benefit matrix as a conditional aggregation robustness operator.

\subsection{Frozen Tie Semantics and Archived Decision Margin}

Within each context, alternatives were sorted by descending raw score and ascending
\texttt{alternative\_id}. The highest unassigned score anchored a tie group; every remaining
alternative within absolute distance $10^{-12}$ of that anchor joined the group and received
the average of the occupied one-based rank positions. Ties were not chained through adjacent
near-equalities. Deterministic Top-1 used
\begin{equation}
\mathcal T_s=\{a:S_{\max,s}-S_{as}\le10^{-12}\}
\label{supp:eq:top1tieset}
\end{equation}
and selected the smallest alternative ID in $\mathcal T_s$. Relative tolerance was zero. These
semantics applied to oracle, MOORA, TOPSIS and Direct-XGBoost scores.

The conditional primary plan also defined the guarded oracle decision margin
\begin{equation}
M_s=\frac{U_{(1),s}^\star-U_{(2),s}^\star}
{\max_aU_{as}^\star-\min_aU_{as}^\star+10^{-12}}.
\label{supp:eq:margin}
\end{equation}
Margin-conditioned analyses were not executed because Pilot B stopped the primary study.

\section{Frozen Pilot-B Gate and Complete Decision}

Pilot B reused development worlds 21001--21005 at the frozen reference
\begin{equation}
(N,c,\rho,\lambda)=(250,0.30,0.4,0.5).
\label{supp:eq:pilotbreference}
\end{equation}
Each world used 200 external TEST contexts and 200 common
RandomWeights draws per world. It was a protocol-level adequacy screen, not independent
population validation. A structured method was eligible only if it was defined without Equal
fallback in all five worlds and Kendall's $\tau_b$ was defined in at least 190 TEST contexts per
world; at least 190 RandomWeights draws also had to be eligible. Cross-world advantages used
medians and required strict positivity in at least four worlds.

The five independently sufficient hard stops were: incomplete coverage; no eligible method
clearing median $\Delta\tau_b=0.05$ or $\Delta$mean-regret $=0.01$ versus RandomWeights; no
eligible method clearing median $\Delta$Top-1 $=0.02$ or $\Delta$mean-regret $=0.01$ versus
MajorityWinner; pooled oracle modal share at least $0.90$; or the frozen random-weight
insensitivity criterion, including its $0.90$ invariant fraction and joint width limits
$0.05$ ($\tau_b$), $0.02$ (Top-1) and $0.01$ (mean regret) in at least four worlds. Passage
required every stop to be false.

\begin{table}[t]
\centering
\caption{Complete frozen Pilot-B gate decision.}
\label{supp:tab:pilotbgates}
\small
\begin{tabularx}{\linewidth}{@{}lXc@{}}
\toprule
Gate & Audited result & Stop\\
\midrule
Coverage & Six structured methods eligible without fallback & No\\
Random equivalence & Four of six methods separated from RandomWeights & No\\
Modal-winner effective tie & Zero of six methods escaped MajorityWinner & \textbf{Yes}\\
Oracle-winner dominance & Pooled modal share $0.526<0.90$ & No\\
Weight insensitivity & Invariant fraction $0.004<0.90$; width-qualified worlds $0<4$ & No\\
\midrule
Overall OR rule & At least one frozen stop fired & \textbf{Yes}\\
\bottomrule
\end{tabularx}
\end{table}

The best strictly positive consistency count was two of five worlds for either Top-1 or mean
regret, and the best cross-world medians were zero. The independently reconstructed decision
therefore barred the protected factorial. This is a negative result about the adequacy of the
frozen comparison at one condition, not evidence that SHAP, weights or MCDM fail universally.

\section{Archived Stage-Wise Fidelity Ladder}

The conditional primary analysis contained nine representations:
\begin{enumerate}
\item \texttt{OracleUtility}, the full interacting oracle;
\item \texttt{OracleMainEffectReference}, with the interaction removed;
\item \texttt{OracleGlobalAttributionNonlinearQ};
\item \texttt{SHAPGlobalAttributionNonlinearQ};
\item \texttt{OracleGlobalAttributionLinearG};
\item \texttt{SHAPGlobalAttributionLinearG};
\item \texttt{OracleWeightMOORA};
\item \texttt{SHAPWeightMOORA};
\item \texttt{DirectXGBoost}, a parallel reference rather than a sequential stage.
\end{enumerate}
Eight paired contrasts isolated interaction removal, oracle global compression, SHAP estimation
on the nonlinear scale, nonlinear-to-linear transformation under each weight vector, MOORA
transformation under each vector and post-MOORA Oracle-versus-SHAP attribution. Kendall's
$\tau_b$ and normalized oracle regret were specified for every representation. The ladder was
diagnostic, not an additive error decomposition, and was not executed because Pilot B stopped
the primary study.

\section{Conditional Primary Factorial}

The four planned primary factors were
\begin{align}
N&\in\{25,50,100,250,1000\}, &
c&\in\{0.10,0.30,0.60\},\label{supp:eq:factors1}\\
\rho&\in\{0,0.4,0.8\}, &
\lambda&\in\{0,0.5,1\}.
\label{supp:eq:factors2}
\end{align}
The resulting $5\times3\times3\times3=135$ conditions were to be repeated over 30
prespecified primary benchmark-instance seeds, giving 4050 seed--condition replications. The
reference condition was
\begin{equation}
(N,c,\rho,\lambda)=(250,0.30,0.4,0.5).
\label{supp:eq:reference}
\end{equation}

\begin{table}[t]
\centering
\caption{Archived nested estimation partitions and fixed TEST size.}
\label{supp:tab:partitions}
\small
\begin{tabular}{@{}rrrr@{}}
\toprule
$N$ & FIT & WEIGHT & external TEST\\
\midrule
25 & 20 & 5 & 200\\
50 & 40 & 10 & 200\\
100 & 80 & 20 & 200\\
250 & 200 & 50 & 200\\
1000 & 800 & 200 & 200\\
\bottomrule
\end{tabular}
\end{table}

Technology-response parameters and common-random-number streams were to be sampled once per
seed and reused across factor levels. For each $(\rho,\lambda)$ cell, the oracle signal scale
$s_U$ was to be calculated on the complete 1000-context FIT+WEIGHT master pool so that changing
$N$ would not change the imposed noise-to-signal ratio. Realized signal-to-noise ratio was to
be logged for every cell. FIT, WEIGHT and external TEST partitions were to remain
context-disjoint, and no TEST row was to be used for model fitting or weight estimation.

\section{Archived Alpha-Dispersion Stress Analysis}

The conditional primary analysis retained the frozen heterogeneous coefficient vector. At the
reference condition only, three additional coefficient structures had been fixed: balanced,
dispersion-aligned and dispersion-anti-aligned. The mapping was obtained from an independent
200-world technology ensemble using
\begin{equation}
D_{j,w}^{attr}
=\frac{1}{M_w}\sum_i\left|q_j(g_{ijw})-\bar q_{j,w}\right|,
\qquad
\bar D_j^{attr}=\E_W[D_{j,W}^{attr}].
\label{supp:eq:qmad}
\end{equation}
The complete ordering agreed between 150 and 200 worlds:
\begin{equation}
C_8>C_{10}>C_5>C_2>C_3>C_4>C_1>C_7>C_6>C_9.
\label{supp:eq:d3order}
\end{equation}

\begin{table}[t]
\centering
\caption{Frozen 200-world mean attribution-scale dispersion.}
\label{supp:tab:d3dispersion}
\small
\begin{tabular}{@{}lc@{}}
\toprule
Criterion & Mean dispersion\\
\midrule
$C_8$ & 0.122873493\\
$C_{10}$ & 0.119403289\\
$C_5$ & 0.110114266\\
$C_2$ & 0.094820283\\
$C_3$ & 0.090482741\\
$C_4$ & 0.088452971\\
$C_1$ & 0.087135074\\
$C_7$ & 0.078706704\\
$C_6$ & 0.049591336\\
$C_9$ & 0.028266887\\
\bottomrule
\end{tabular}
\end{table}

Balanced coefficients were $\alpha_j=0.10$. The aligned and anti-aligned mappings were
\begin{align}
\boldsymbol\alpha^{aligned}:
&(C_8,C_{10},C_5,C_2,C_3,C_4,C_1,C_7,C_6,C_9)
\nonumber\\
&\mapsto(0.16,0.14,0.13,0.12,0.11,0.10,0.08,0.07,0.05,0.04),
\label{supp:eq:alphaaligned}\\
\boldsymbol\alpha^{anti}:
&(C_9,C_6,C_7,C_1,C_4,C_3,C_2,C_5,C_{10},C_8)
\nonumber\\
&\mapsto(0.16,0.14,0.13,0.12,0.11,0.10,0.08,0.07,0.05,0.04).
\label{supp:eq:alphaanti}
\end{align}
The frozen primary heterogeneous assignment had weak rank association with the dispersion
ordering ($\rho_S=0.1273$, $\tau_b=0.0222$), so it was not changed. Although the D3 mapping
itself was a completed development artefact, the planned primary stress comparisons were never
executed.

\section{Archived Bootstrap and Robustness Plan}

At the reference condition, $B=200$ context-level bootstrap replications were to resample FIT
and WEIGHT contexts separately, refit the model, recompute attributions and weights, and
evaluate decisions on the fixed original external TEST contexts. Confidence intervals were to
be reported for criterion weights, total-variation weight error, Kendall's $\tau_b$, Top-1
accuracy and normalized oracle regret. No cross-context probability that a particular ITS was
ranked first was to be interpreted as a global deployment probability.

A 50-oracle ensemble was to leave the frozen D29 technology-response kernel unchanged while
varying oracle coefficient vectors, nonlinearities and semantically admissible interaction
pairs. Relative outcome noise was to be recalibrated to each oracle variant's signal standard
deviation. This ensemble would not replace the response generator and therefore would not
reintroduce the historical v2.1 separability defect through oracle-only perturbations.

Two distribution-shift analyses were specified: holding out the upper 20\% of demand pressure
and holding out the lower 20\% of infrastructure readiness. Secondary robustness analyses were
to include criterion removal, $\pm10/20/30\%$ SHAP-weight perturbations,
leave-one-alternative-out analysis and MOORA--TOPSIS comparison.

\section{Archived Statistical Analysis Plan}

Factor effects were to be summarized with repeated-seed regression or mixed-effects models,
emphasizing effect sizes and uncertainty intervals rather than binary significance alone.
Method comparisons were paired because every method would share the same benchmark instance.
The 30 primary seeds were to be treated as benchmark-instance clusters; alternative--context
rows were not to be interpreted as independent superpopulation replications.

Method availability was to be reported explicitly. For each factorial cell, undefined
Oracle/SHAP weight rates and recorded Equal-fallback rates were to accompany conditional
decision-fidelity summaries. Pairwise contrasts were to use only benchmark instances on which
both methods were defined, with coverage and fallback rates reported alongside them. Undefined
Oracle or SHAP vectors were never to be silently replaced by Equal weights.

Planned outputs included predictive-performance and Direct-XGBoost references, weight-recovery
tables, decision-fidelity summaries, factorial effect summaries, alpha-alignment stress results,
bootstrap intervals, validity-boundary figures and the prespecified stage-wise fidelity ladder.
None of these primary outputs exists under v1.

\section{Reproducibility and Content-Integrity Record}
\label{supp:sec:integrity}

Table~\ref{supp:tab:digests} gives non-identifying truncated content digests. The selected B100
background exactly matched the frozen TreeSHAP identity set.

\begin{table}[t]
\centering
\caption{Truncated, non-identifying content digests of the frozen artefacts underlying the
reported numbers.}
\label{supp:tab:digests}
\small
\begin{tabular}{@{}ll@{}}
\toprule
Artifact & Truncated SHA-256\\
\midrule
Oracle B100 result & \texttt{7f9864bd...439c9fb0}\\
TreeSHAP background identities & \texttt{3184c675...1646eda1}\\
Pilot-B result & \texttt{05072857...9f787a88}\\
Fixed-weight selection-map result & \texttt{a1f293b6...d7e20f4e}\\
\bottomrule
\end{tabular}
\end{table}

The full anonymous provenance package records, in order, the
threshold freeze, evaluator/executor freeze, scientifically invariant direct-file loader repair,
audited Pilot-B result freeze, immutable closure reference, post-closure diagnostic protocol
freeze, implementation freeze and independently audited result freeze. The worktree was clean
after each closure point, and the final complete test run reported 836 passes with three
pre-existing external plotting warnings.
\ifanon
Exact repository identifiers are intentionally omitted here to protect anonymous review.
\else
Exact repository identifiers accompany the public reproducibility archive.
\fi
These facts document software and artefact integrity; they are not statistical evidence.

\section{Prospective Stop and Permanent Closure}

Pilot B was governed by five independently sufficient hard stops: coverage failure, failure to
separate RandomWeights, failure to escape MajorityWinner, pooled oracle-winner dominance and
weight insensitivity. Passage required every flag to be false. In the frozen result, only the
modal-winner effective-tie flag was true, but the overall OR rule therefore barred every
analysis in this supplement.

The immutable version-control, digest and test record appears in
Section~\ref{supp:sec:integrity}. Pilot B v1 is closed and must not be rerun. Any later study
requires a new specification and version reference, explicit disclosure of the v1 failure and
development-only design work; it cannot alter the v1 thresholds, result or closure reference.
The fixed-weight selection-map diagnostic in Section~\ref{supp:sec:fixedselection} is such a
separately frozen descriptive layer; it neither reopened Pilot B nor authorized any protected
analysis.

\end{document}
"
% ============================================================
% ChkTeX: intentional en dashes, literal digest ellipses, label placement,
% acronym punctuation and LNCS name/bibliography conventions.
% chktex-file 8
% chktex-file 11
% chktex-file 13
% chktex-file 24
% chktex-file 36
% chktex-file 38

\documentclass[runningheads,orivec]{llncs}

\usepackage[T1]{fontenc}
\usepackage[utf8]{inputenc}
\usepackage{amsmath,amssymb,amsfonts}
\usepackage{booktabs}
\usepackage{graphicx}
\usepackage{tabularx}
\usepackage{hyperref}
\hypersetup{hidelinks}

\newcommand{\E}{\mathbb{E}}
\newcommand{\median}{\operatorname{median}}
\newcommand{\IQR}{\operatorname{IQR}}
\newcommand{\SD}{\operatorname{SD}}
\newcommand{\clip}{\operatorname{clip}}
\newcommand{\XGB}{XGBoost}

% ---- build-mode switch: anonymous by default -----------------
\providecommand{\CAMERAREADY}{0}
\newif\ifanon
\anontrue
\ifnum\CAMERAREADY=1 \anonfalse \fi
% --------------------------------------------------------------

\begin{document}
\emergencystretch=1.5em

\title{Supplementary Methods and Audit Archive for\texorpdfstring{\\}{ }
``From Predictive Importance to Decision Value: Evaluating XAI-Based Weights for
ITS Prioritisation''}
\titlerunning{Supplementary Benchmark-Adequacy Record}

\ifanon
  \author{Anonymous Author(s)}
  \authorrunning{Anonymous Author(s)}
  \institute{Anonymous Institution(s)}
\else
  \author{Mohamed Yassine Neifar\inst{1,2} \and
  Ahmed Frikha\inst{2} \and
  Nadir Farhi\inst{1}}
  \authorrunning{Neifar et al.}
  \institute{Universit\'e Gustave Eiffel, COSYS--GRETTIA,
  Champs-sur-Marne, F-77447 Marne-la-Vall\'ee Cedex 2, France
  \and
  Institut Sup\'erieur de Gestion Industrielle de Sfax (ISGIS),
  OLID, University of Sfax, Sfax 3038, Tunisia}
\fi

\maketitle

\section{Status and Interpretation Boundary}

This supplement reports executed development audits and preserves the detailed primary and
robustness analyses that had been specified before the prospectively frozen Pilot-B
benchmark-adequacy gate was executed. The latter analyses were conditional on all five Pilot-B
hard-stop flags being false. The modal-winner effective-tie stop fired, so the primary factorial
was not authorized. Primary seeds 11001--11030 and reserve seeds 30001--30005 were never
accessed; all sections explicitly labelled ``Archived'' remain unexecuted protocols rather than
results.

The closed v1 state is anchored by an immutable version-control reference.
\ifanon
Exact repository identifiers and complete digests are withheld from this anonymous reviewer
file; they are available confidentially to editors and will accompany the public
reproducibility archive after review.
\else
Exact repository identifiers and complete digests accompany the public reproducibility archive.
\fi
This document is therefore a methods, audit-results and archival-protocol supplement. It is not
authorization to execute the protected design.
After permanent Pilot-B closure, one separately frozen development-only descriptive diagnostic
reconstructed the context-level MOORA selection maps of the seven observed archived fixed
vectors. It used no primary/reserve world, generated no target outcome, fitted or queried no
model, re-estimated no weight and could not alter the gate. Its protocol, implementation,
independent audit and result were frozen as a distinct post-closure layer.
Methodological citations are provided in the main manuscript and are not duplicated here.

\subsection*{Question Structure and Evidential Status}

The main manuscript states two prospective research questions and one separate post-closure
descriptive question. RQ1 asks whether structural diagnostics can reveal that apparent
contextual variation disappears under the normalization used by the downstream MCDM operator.
RQ2 asks whether, after that structural defect is corrected, structured global-weight pipelines
provide sufficient incremental decision value over random-weight and constant-decision
references to justify protected primary evaluation. Both are addressed by the prospectively
structured benchmark-development record and the frozen Pilot-B result reported here.

DQ1 asks what the archived context-level selection maps reveal about the observed decision
behaviour of the seven frozen fixed-weight vectors. DQ1 is post-closure and descriptive. It is
non-confirmatory, it carries no pass/fail role, and it cannot modify, reinterpret or reopen the
Pilot-B gate. Section~\ref{supp:sec:fixedselection} reports it under that label.

Throughout this supplement, \emph{each replication seed defines one benchmark world}, and the
world is the statistical unit. Where a numbered identifier such as 21001 appears it names a
world; the phrase ``replication seed'' is retained only when the frozen artefact field or
protocol clause is itself named that way. \emph{Validity} denotes conceptual scientific
properties of the benchmark; \emph{adequacy} denotes the operational prospectively frozen gate
used to decide whether protected primary execution was warranted.

\subsection*{Reader Map}

The conference-facing main manuscript follows a problem--method--result--meaning structure and
deliberately omits most development chronology. Table~\ref{supp:tab:readermap} directs readers
to the corresponding detailed evidence. This separation changes presentation only: the frozen
scientific record, Pilot-B decision and post-closure interpretation boundary are unchanged.

\begin{table}[t]
\centering
\caption{Navigation from the concise main argument to the detailed reproducibility record,
with the evidential status of each block. Status codes are: EDA, executed development audit;
PBA, Pilot-B confirmatory adequacy result; PCD, post-closure descriptive diagnostic; ANE,
archived and not executed.}
\label{supp:tab:readermap}
\small
\begin{tabularx}{\linewidth}{@{}lXc@{}}
\toprule
\textbf{Main-paper element} & \textbf{Detailed supplementary evidence} & \textbf{Status}\\
\midrule
Controlled ITS laboratory & Criterion definitions, capability pathways, copula contexts and
opportunity functions in Sect.~2. & EDA\\
Oracle and global attribution & Closed-form functional, joint-background moments, \XGB,
TreeSHAP and B100 checks in Sect.~3. & EDA\\
Structural correction & Calibration, positive control, numerical null and candidate screening
in Sect.~4. & EDA\\
D29 decision geometry & World-resolved oracle orderings, winners, modal shares and regret in
Sect.~5. & EDA\\
Pilot-B evidence (RQ2) & Per-world method contrasts, gate decision and objective-weight vectors
in Sects.~5 and 7. & PBA\\
Observed selection maps (DQ1) & All audited modal identities, shares, unique-choice counts and
pairwise agreement summaries in Sect.~5. & PCD\\
Decision operators and gate & Exact weighting, tie, regret and hard-stop rules in Sects.~6--7.
& EDA\\
Unexecuted primary plan & Clearly labelled archival protocols in Sects.~8--12. & ANE\\
Integrity and closure & Content digests, software tests and permanent-stop record in
Sects.~13--14. & EDA\\
\bottomrule
\end{tabularx}
\end{table}

\section{Detailed Benchmark Construction}

\subsection{Criteria and Capability Pathways}

The six alternatives in the main manuscript were evaluated on the ten technology-agnostic
criteria in Table~\ref{supp:tab:criteria}. The original lifecycle-burden criterion was a cost;
all other criteria were benefits.

\begin{table}[t]
\centering
\caption{Controlled evaluation criteria.}
\label{supp:tab:criteria}
\small
\begin{tabularx}{\linewidth}{@{}c X c@{}}
\toprule
Code & Criterion & Original direction\\
\midrule
$C_1$ & Mobility-efficiency improvement & Benefit\\
$C_2$ & Travel-time reliability improvement & Benefit\\
$C_3$ & Person-throughput improvement & Benefit\\
$C_4$ & Safety-surrogate improvement & Benefit\\
$C_5$ & Non-recurring-congestion resilience & Benefit\\
$C_6$ & Environmental-efficiency improvement & Benefit\\
$C_7$ & User and accessibility benefit & Benefit\\
$C_8$ & Infrastructure and data readiness & Benefit\\
$C_9$ & Interoperability and scalability & Benefit\\
$C_{10}$ & Annualised lifecycle deployment burden & Cost\\
\bottomrule
\end{tabularx}
\end{table}

The pre-specified capability mask in Table~\ref{supp:tab:mask} defines synthetic functional
pathways, not empirical impact magnitudes. Direct (D), indirect (I) and zero (0) pathways set
broad response-amplitude bands.

\begin{table}[t]
\centering
\caption{Pre-specified structural capability mask.}
\label{supp:tab:mask}
\scriptsize
\begin{tabular}{@{}lccccccc ccc@{}}
\toprule
 & $C_1$ & $C_2$ & $C_3$ & $C_4$ & $C_5$ & $C_6$ & $C_7$ & Ready & Interop. & Burden\\
\midrule
ATSC      & D & D & D & I & I & I & I & M    & H & M\\
TSP       & I & D & D & I & 0 & I & D & M    & H & L--M\\
TIDM      & I & D & I & D & D & I & I & M    & H & M\\
RMTI--DRG & I & I & I & I & D & I & D & L--M & H & L--M\\
V2X--CS   & I & I & 0 & D & I & I & I & H    & H & H\\
DSM       & D & D & I & D & D & I & I & M--H & H & M--H\\
\bottomrule
\end{tabular}
\end{table}

\subsection{Gaussian-Copula Contexts and Opportunity Functions}

For demand, incident, person-movement, environmental and readiness factors, generate
$z_{0,s},z_{1,s},\ldots,z_{5,s}\overset{iid}{\sim}\mathcal N(0,1)$ and set
\begin{equation}
\widetilde h_{k,s}^{(\rho)}
=\sqrt{\rho}\,z_{0,s}+\sqrt{1-\rho}\,z_{k,s},
\qquad h_{k,s}^{(\rho)}=\Phi(\widetilde h_{k,s}^{(\rho)}).
\label{supp:eq:copula}
\end{equation}
The same base draws were reused across dependence levels, while estimation and external TEST
streams used disjoint namespaces. The first seven criterion opportunities were
\begin{align}
o_{s1}&=h_D, & o_{s2}&=0.60h_D+0.40h_I, & o_{s3}&=0.60h_D+0.40h_T,\nonumber\\
o_{s4}&=0.50h_D+0.50h_I, & o_{s5}&=h_I, & o_{s6}&=0.60h_D+0.40h_E,\nonumber\\
o_{s7}&=0.40h_D+0.60h_T.&&&
\label{supp:eq:opportunities}
\end{align}

Each replication seed $r$ defined one benchmark world $W_r$: technology parameters, a master
context stream and all common-random-number streams. The conditional method target had been
\begin{equation}
\Psi_m(x)=\E_W\!\left[\E_{H,\varepsilon}
\{L_m(x;W,H,\varepsilon)\mid W\}\right].
\label{supp:eq:estimand}
\end{equation}
The world, rather than an alternative--context row, was the intended statistical unit. The
primary worlds needed to estimate this quantity were never accessed after the gate fired.

\subsection{Historical and Frozen Response Layers}

For $C_1$--$C_7$, capability amplitudes were drawn once per world:
\begin{equation}
\theta_{aj}\sim
\begin{cases}
0,&M_{aj}=0,\\
\mathcal U(0.05,0.20),&M_{aj}=I,\\
\mathcal U(0.20,0.40),&M_{aj}=D.
\end{cases}
\label{supp:eq:effectbands}
\end{equation}
The historical v2.1 response was
\begin{equation}
x_{asj}=\clip\!\left(\theta_{aj}o_{sj}+\sigma_x\xi_{asj},0,1\right).
\label{supp:eq:v21}
\end{equation}
The deployment classes were
\begin{equation}
L\sim U(0.25,0.40),\qquad
M\sim U(0.45,0.60),\qquad
H\sim U(0.65,0.80).
\label{supp:eq:deploymentclasses}
\end{equation}
Compound classes were equal-probability mixtures. With readiness
requirement $r_a$, interoperability $\iota_a$ and burden $\kappa_a$,
\begin{align}
x_{as8}&=\clip\!\left(1-\max(0,r_a-h_{R,s})+\eta_{as8},0,1\right),\nonumber\\
x_{as9}&=\clip\!\left(\iota_a+0.10h_{R,s}+\eta_{as9},0,1\right),\nonumber\\
x_{as,10}&=\clip\!\left(\kappa_a+0.20(1-h_{R,s})r_a+\eta_{as,10},0,1\right).
\label{supp:eq:c8c10}
\end{align}
Benefit orientation used $g_{asj}=x_{asj}$ for $j\le9$ and
$g_{as,10}=1-x_{as,10}$. At $\sigma_x=0$, the first seven criteria reduced to
$g_{asj}=\theta_{aj}o_{sj}$, which explains the exact normalization cancellation in the main
paper.

The frozen D29 kernel was
\begin{equation}
g_{asj}=\theta_{aj}o_{sj}+0.05v_{aj}o_{sj}(1-o_{sj}),\qquad j\le7.
\label{supp:eq:d29}
\end{equation}
For $m_j$ active alternatives, its ordered contrast grid was
$v^{grid}_{\ell j}=-1+2(\ell-1)/(m_j-1)$. After sorting active IDs, one deterministic NumPy
permutation used
\begin{equation}
\texttt{SeedSequence([r, 2010, j])}.
\label{supp:eq:d29seedsequence}
\end{equation}
No outcome or
winner information entered this assignment. Active $\theta_{aj}\ge0.05$ ensured
$\partial g/\partial o=\theta_{aj}+0.05v_{aj}(1-2o)\ge0$. Structural-zero pathways remained
zero, and $C_8$--$C_{10}$ retained Eq.~\eqref{supp:eq:c8c10}.

For $o_{sj}>0$, within-context vector normalization over active alternatives gives the exact
geometry
\begin{equation}
r_{asj}=
\frac{\theta_{aj}+0.05v_{aj}(1-o_{sj})}
{\sqrt{\sum_{b\in\mathcal A_j}
[\theta_{bj}+0.05v_{bj}(1-o_{sj})]^2}}.
\label{supp:eq:d29normalized}
\end{equation}
Thus D29 removes the v2.1 cancellation, but each criterion-specific normalized profile remains
on a one-parameter curve indexed by $o_{sj}$. Because opportunities differ across criteria,
Eq.~\eqref{supp:eq:d29normalized} alone does not prove that a fixed-weight MOORA winner is
context invariant. At Pilot-B closure, fixed-weight invariance was therefore an untested
hypothesis rather than an explanation of the stop. The subsequently frozen descriptive
selection-map diagnostic reported below resolves this question for the seven observed vectors
only; it does not establish a causal effect of Eq.~\eqref{supp:eq:d29normalized} or characterize
unobserved weight vectors.

\section{Oracle, Learning and Background Details}

\subsection{Oracle Utility and Closed-Form Attribution}

The monotone transform, coefficient vector and interaction graph were
\begin{align}
q_j(g)&=\frac{\log(1+2g)}{\log3},\nonumber\\
\boldsymbol\alpha&=(0.11,0.04,0.05,0.07,0.10,0.14,0.12,0.16,0.08,0.13),\nonumber\\
\mathcal E&=\{(1,2),(3,7),(4,5),(6,10),(8,9)\},\qquad \beta_{jk}=0.20.
\label{supp:eq:oraclecomponents}
\end{align}
For $z_j=q_j(g_j)$, the empirical FIT-background moments were
$\mu_j=\E_{bg}(z_j)$ and $\mu_{jk}=\E_{bg}(z_jz_k)$; dependence was retained by not replacing
$\mu_{jk}$ with $\mu_j\mu_k$. The main-effect contribution was
$\phi_j^{main}=\alpha_j(z_j-\mu_j)$. For pair $\beta_{jk}z_jz_k$, feature $j$ received
\begin{equation}
\phi_j^{(jk)}=\frac{\beta_{jk}}{2}
\{z_j\mu_k-\mu_{jk}+z_jz_k-\mu_jz_k\},
\label{supp:eq:pairshap}
\end{equation}
with the symmetric expression for $k$. The closed-form attribution was
\begin{equation}
\phi_j^\star(\mathbf g)=\frac{1}{1+\lambda}
\left[\alpha_j(z_j-\mu_j)+\lambda\!\sum_{k:(j,k)\in\mathcal E}\phi_j^{(jk)}\right],
\label{supp:eq:closedshap}
\end{equation}
with symmetric handling when $j$ was second in a pair. Exhaustive coalition enumeration was a
unit-test reference. Global weights normalized mean absolute WEIGHT-row attributions; total
importance at most $10^{-12}$ made the vector undefined, with no Equal substitution.

\subsection{XGBoost and TreeSHAP Freeze}

XGBoost was selected once by context-grouped cross-validation on development FIT rows at the
reference condition. Candidate 45 attained mean RMSE $0.017757058913365983$ with
\begin{equation}
\begin{split}
&n_{\mathrm{estimators}}=600,\quad d_{\max}=2,\quad \eta=0.05,\\
&\text{subsample}=0.8,\quad \text{colsample\_bytree}=0.8,\quad
\text{min\_child\_weight}=5,\\
&\text{reg\_alpha}=0,\quad \text{reg\_lambda}=1.
\end{split}
\label{supp:eq:xgbparams}
\end{equation}
The objective was \texttt{reg:squarederror}, with one computational thread and estimator
random-state namespace 82002. The selected model was then fixed. Interventional TreeSHAP used 100 deterministic FIT
identities selected with row-priority namespace 83001 and evaluated mean absolute attributions
on WEIGHT. External TEST did not enter calibration, background selection or weighting.

\subsection{Matched B100 Oracle Diagnostic}

The primary oracle used all eligible FIT rows for $\mu_j$ and $\mu_{jk}$. The secondary B100
diagnostic changed only these moments to the exact 100 frozen TreeSHAP identities; WEIGHT rows,
the oracle functional and every other setting were unchanged. It compared full-FIT and B100
weight vectors through mean absolute difference, total variation, maximum absolute difference,
Spearman correlation, Top-3/Top-5 overlap and maximum marginal/joint moment changes. Undefined
vectors remained undefined. The only authorized pre-primary execution was seed 21001 at
$N=250$, $\rho=0.4$ and $\lambda=0.5$; no $Y$, SHAP package call or MCDM operation was involved.

\section{Detailed Structural Calibration and Candidate Screening}

\subsection{Structural and Reachability Metrics}

Let $G_j\in\mathbb R^{S\times A_j}$ denote the active-alternative response matrix for criterion
$j$ at zero measurement noise. If its singular values are
$\sigma_{1j}\ge\sigma_{2j}\ge\cdots$ and
$E_j=\sum_\ell\sigma_{\ell j}^2$, the rank-one residual was
\begin{equation}
NS_j=1-\frac{\sigma_{1j}^2}{E_j},
\qquad E_j>10^{-24}.
\label{supp:eq:NS}
\end{equation}
No epsilon was inserted into this ratio; structural-zero matrices were reported as undefined.
For each jointly positive active alternative pair,
\begin{equation}
LRV_{ab,j}=\SD_s\!\left[\log(g_{asj}/g_{bsj})\right],
\qquad
LRV50_j=\median_{a<b}LRV_{ab,j},
\label{supp:eq:LRV}
\end{equation}
using population SD. For normalization family $m$,
\begin{equation}
NSV_{j,m}=\left\{S^{-1}\sum_s
\left\|\mathbf r_{sj}^{(m)}-\bar{\mathbf r}_j^{(m)}\right\|_2^2\right\}^{1/2}.
\label{supp:eq:NSV}
\end{equation}
Vector normalization was the primary structural gate; sum, max and min--max variants were
robustness diagnostics.

Reachability reflected one idiosyncratic Gaussian-copula shock while holding the shared and all
other idiosyncratic shocks fixed. With opportunity functions recomputed after reflection,
\begin{equation}
RE_{kj}=\median_{s,a\in\mathcal A_j}|g_{asj}^{(-k)}-g_{asj}|,
\qquad
SRE_{kj}=\frac{RE_{kj}}{\IQR(g_j)+10^{-12}}.
\label{supp:eq:SRE}
\end{equation}

\subsection{Scientific Positive Control}

For active $C_1$--$C_7$ pathways, the dedicated positive control was
\begin{equation}
g_{asj}^{PC,\ell}=\left[\theta_{aj}+0.05c_a^{(\ell)}(2o_{sj}-1)\right]o_{sj},
\label{supp:eq:positivecontrol}
\end{equation}
with contrast base $(-1,-0.6,-0.2,0.2,0.6,1)$ and all six cyclic rotations. Structural-zero
pathways remained zero and $C_8$--$C_{10}$ were unchanged. Rotation medians were formed within
development worlds 21001--21005 and world medians were then used as scientific references.

\begin{table}[t]
\centering
\caption{Frozen D2.7 criterion-level scientific references.}
\label{supp:tab:d27layera}
\small
\begin{tabular}{@{}ccc@{}}
\toprule
Criterion & $T^{sci}_{LRV,j}$ & $T^{sci}_{NSV,j}$\\
\midrule
$C_1$ & 0.185606 & 0.101685\\
$C_2$ & 0.107107 & 0.069436\\
$C_3$ & 0.163029 & 0.064900\\
$C_4$ & 0.115796 & 0.069423\\
$C_5$ & 0.127964 & 0.073681\\
$C_6$ & 0.194899 & 0.129935\\
$C_7$ & 0.144080 & 0.073092\\
\bottomrule
\end{tabular}
\end{table}

\begin{table}[t]
\centering
\caption{Frozen D2.7 latent-pathway reachability references.}
\label{supp:tab:d27sre}
\small
\begin{tabular}{@{}lc@{}}
\toprule
Pathway & $T^{sci}_{SRE,kj}$\\
\midrule
$h_D\rightarrow C_1$ & 0.556245\\
$h_D\rightarrow C_2$ & 0.334437\\
$h_D\rightarrow C_3$ & 0.284878\\
$h_D\rightarrow C_4$ & 0.292844\\
$h_D\rightarrow C_6$ & 0.431612\\
$h_D\rightarrow C_7$ & 0.252576\\
$h_E\rightarrow C_6$ & 0.296869\\
$h_I\rightarrow C_2$ & 0.212043\\
$h_I\rightarrow C_4$ & 0.294905\\
$h_I\rightarrow C_5$ & 0.455271\\
$h_T\rightarrow C_3$ & 0.186995\\
$h_T\rightarrow C_7$ & 0.373380\\
\bottomrule
\end{tabular}
\end{table}

The paired positive-control uplift in mean normalized constant-modal regret was
$(0.001485,0.020022,-0.013836,0.015534,0.006161)$ for worlds 21001--21005. Because the activation
rule required positive uplift in all five worlds, no Layer-B scientific regret reference was
defined and no post-hoc value was introduced.

\subsection{Numerical Null and Candidate Outcomes}

Numerical thresholds were calibrated from 2000 exact rank-one replicates for each of five- and
six-alternative active sets:
\begin{equation}
T^{num}_{NS}=10^{-12},\qquad
T^{num}_{LRV}=10^{-10},\qquad
T^{num}_{NSV}=3.2506084454\times10^{-8}.
\label{supp:eq:nullthresholds}
\end{equation}
All 35 historical v2.1 world--criterion rows fell below every threshold. For active pathways
\(j\leq7\), the D4-F1 curvature family was
\(g_{asj}^{D4F1}(\chi)=\theta_{aj}o_{sj}[1+\chi v_{aj}(1-o_{sj})]\), with structural-zero
pathways fixed at zero. It cleared some structural references at higher curvature but lost reachability; no
$\chi\in\{0.1,\ldots,1.0\}$ passed all gates, so the family was rejected.

The later exact D2.9 audit evaluated all $6^5=7776$ prespecified positive-control profile
assignments. None met the old simultaneous binary conjunction. This established non-viability
for that calibrated design, not mathematical impossibility. Under the corrected prospective
rule, 54 candidates were eligible. D29 was retained solely for protocol continuity and minimal
intervention, not scientific preference, downstream performance, SHAP pattern or ITS identity.

The count of 54 is composed rather than directly observed in a single artefact, and is
decomposed here to prevent confusion. The nonselective retrospective re-adjudication examined 54
historical parameter points and returned 44 \texttt{VERIFIED\_\allowbreak ELIGIBLE},
0 \texttt{VERIFIED\_\allowbreak INELIGIBLE} and 10
\texttt{INSUFFICIENT\_\allowbreak FROZEN\_\allowbreak EVIDENCE}. The 10
unresolved points were all v2.2 points whose frozen artefact schema did not record C8--C10
invariance. A separate supplemental resolution then established that invariance from the
contemporaneous frozen source contract, and all 10 satisfied the corrected hard constraints.
The frozen retention artefact therefore records an eligible universe of
$44+10=54$ specifications. Eligibility is explicitly non-unique: the retention of D29 is a
minimal-intervention and provenance-stability decision, not a superiority, uniqueness or
statistical-selection claim, and no ranking or re-optimization over the 54 points was
performed.

\section{Detailed Development Audit Results}

\subsection{D29 Layer-B Decision Geometry}

The one-shot characterization used development worlds 21001--21005, the complete $N=1000$
FIT+WEIGHT pool, $\rho=0.4$, $\sigma_x=0$ and $\lambda=0.5$. It generated no $Y$ and executed
no learned model, SHAP package or MCDM method. Regret used the exact context-specific oracle
range and was undefined only if that range was at most $10^{-12}$; no context was undefined.

\begin{table}[t]
\centering
\caption{Complete frozen D29 Layer-B development record.}
\label{supp:tab:d29layerb}
\scriptsize
\resizebox{\linewidth}{!}{%
\begin{tabular}{@{}r r r l r r r r@{}}
\toprule
World & Orderings & Winners & Modal & Modal share & Mean regret & Median regret & $P_{95}$ regret\\
\midrule
21001 & 23 & 2 & $A_6$ & 0.983 & 0.00042018 & 0 & 0\\
21002 & 78 & 5 & $A_1$ & 0.519 & 0.05992337 & 0 & 0.29155475\\
21003 & 24 & 4 & $A_2$ & 0.745 & 0.04523642 & 0 & 0.28747533\\
21004 & 17 & 3 & $A_2$ & 0.827 & 0.01734351 & 0 & 0.12641643\\
21005 & 38 & 4 & $A_2$ & 0.809 & 0.01606539 & 0 & 0.12654230\\
\bottomrule
\end{tabular}}
\end{table}

Across worlds, the medians were 24 orderings, four winners, modal share $0.809$, mean regret
$0.01734351$, median regret zero and $P_{95}$ regret $0.12654230$. At the directly comparable
world 21001, the historical v2.1 record had 19 orderings, two winners, modal share $0.962$ and
mean always-modal regret $0.0011106$ under its frozen $D_s+10^{-12}$ denominator. D29 had 23
orderings, two winners, modal share $0.983$ and mean regret $0.00042018$ under exact $D_s$.
The comparison therefore supports neither universal degeneration nor universal improvement.

\subsection{XGBoost and TreeSHAP Development Checks}

The selected XGBoost hyperparameters are reported in Eq.~\eqref{supp:eq:xgbparams}; candidate
number and grouped-CV RMSE are given above. At world 21001 and the reference condition, total mean-absolute
TreeSHAP importance was $0.07006131011183041$ and maximum local-accuracy error was
$6.45\times10^{-7}$. The development-only weight vector for $(C_1,\ldots,C_{10})$ was
\begin{equation}
\begin{split}
&(0.124860,\ 0.113194,\ 0.069426,\ 0.025945,\ 0.155169,\\
&\qquad 0.078339,\ 0.071621,\ 0.265059,\ 0.020729,\ 0.075658).
\end{split}
\label{supp:eq:devshapweights}
\end{equation}
These observations did not tune comparator definitions, MCDM rules or winner semantics.

\subsection{B100-Matched Background Result}

Both oracle vectors were defined, and the 100 identities exactly matched the frozen TreeSHAP
membership. Total mean-absolute importance was $0.06961866590$ under full-FIT moments and
$0.06992189186$ under B100 moments.

\begin{table}[t]
\centering
\caption{Development-only oracle-attribution weights under full-FIT and matched B100 moments.}
\label{supp:tab:b100weights}
\small
\begin{tabular}{@{}lrrr@{}}
\toprule
Criterion & Full-FIT & B100 & Absolute difference\\
\midrule
$C_1$  & 0.09663438 & 0.09712037 & 0.00048600\\
$C_2$  & 0.04849942 & 0.04854063 & 0.00004122\\
$C_3$  & 0.05910838 & 0.05913637 & 0.00002799\\
$C_4$  & 0.06665455 & 0.06559847 & 0.00105608\\
$C_5$  & 0.14880967 & 0.14667297 & 0.00213669\\
$C_6$  & 0.09368358 & 0.09388953 & 0.00020596\\
$C_7$  & 0.07615062 & 0.07594850 & 0.00020212\\
$C_8$  & 0.24766693 & 0.25232065 & 0.00465372\\
$C_9$  & 0.04819959 & 0.04935317 & 0.00115358\\
$C_{10}$ & 0.11459290 & 0.11141932 & 0.00317357\\
\bottomrule
\end{tabular}
\end{table}

The weight-vector discrepancies were
\begin{equation}
\operatorname{MAE}_w=0.0013136928,\quad
TV_w=0.0065684640,\quad
\max_j|\Delta w_j|=0.0046537223.
\label{supp:eq:b100distances}
\end{equation}
Spearman correlation was $0.98787879$, and Top-3 and Top-5 overlaps were both $1.0$.
Maximum absolute marginal- and joint-moment changes were $0.0098907938$ and $0.0041811926$,
respectively. Equal fallback was not used. These results show small
finite-background sensitivity at this one development reference, not uniform robustness.

\subsection{World-Resolved Pilot-B Oracle and SHAP Comparison}

Table~\ref{supp:tab:pilotboracleshap} retains the full comparison moved from the main paper.
Each entry is an advantage over the world-specific MajorityWinner reference; positive values
favour the structured method.

\begin{table}[t]
\centering
\caption{Per-world closed-form oracle-attribution and SHAP advantages over MajorityWinner.}
\label{supp:tab:pilotboracleshap}
\small
\begin{tabular}{@{}crrrr@{}}
\toprule
World & Oracle $\Delta$Top-1 & SHAP $\Delta$Top-1 & Oracle $\Delta$regret & SHAP $\Delta$regret\\
\midrule
21001 & $ 0.000$ & $ 0.000$ & $ 0.000000$ & $ 0.000000$\\
21002 & $ 0.035$ & $ 0.030$ & $ 0.001538$ & $ 0.003064$\\
21003 & $-0.465$ & $ 0.020$ & $-0.143522$ & $ 0.006717$\\
21004 & $-0.010$ & $ 0.000$ & $-0.002119$ & $ 0.000000$\\
21005 & $-0.380$ & $-0.035$ & $-0.090328$ & $-0.002986$\\
\bottomrule
\end{tabular}
\end{table}

The closed-form method exactly implements the frozen global interventional attribution
functional; it is not an exact representation of context-specific oracle utility. Thus its
failure to escape MajorityWinner is evidence about the frozen global-compression--MOORA
pipeline, not a numerical error in the attribution formula.

\subsection{Illustrative Frozen Objective-Weight Vectors}

Table~\ref{supp:tab:objectivevectors} directly transcribes the CRITIC and Entropy vectors stored
for development world 21001; no new transformation of TEST decisions was performed. The display
cannot alter the frozen gate or establish a performance mechanism. Both vectors were defined
without Equal fallback, and neither has its largest component on C8 or C10. Thus concentration
on those two criteria is not supported in this world. The table is deliberately not combined
with five-world median performance to infer cross-world geometry.

\begin{table}[t]
\centering
\caption{Frozen objective-weight vectors for development world 21001.}
\label{supp:tab:objectivevectors}
\scriptsize
\setlength{\tabcolsep}{2.8pt}
\begin{tabular}{@{}lrrrrrrrrrr@{}}
\toprule
Method & C1 & C2 & C3 & C4 & C5 & C6 & C7 & C8 & C9 & C10\\
\midrule
CRITIC & .079 & .075 & .101 & .093 & .150 & .099 & .097 & .111 & .084 & .111\\
Entropy & .131 & .114 & .175 & .099 & .236 & .064 & .084 & .027 & .032 & .038\\
\bottomrule
\end{tabular}
\end{table}

The aggregate coincidence in world 21003 did not, by itself, identify a common decision map.
Equal, CRITIC and Entropy each recorded Top-1 accuracy $0.190$ and mean normalized regret
$0.316758$, but their mean Kendall $\tau_b$ values were respectively $0.5493$, $0.5427$ and
$0.5113$. Top-1 accuracy and normalized regret depend on the selected alternative, whereas
Kendall $\tau_b$ depends on the complete ranking. The $\tau_b$ differences therefore establish
different lower-order ranking behaviour on average but are logically uninformative about
argmax equality. That ambiguity motivated the separately frozen diagnostic below.

\subsection{Frozen Post-Closure Fixed-Weight Selection-Map Diagnostic (DQ1)}
\label{supp:sec:fixedselection}

\noindent\textbf{Evidential status: post-closure, descriptive, non-confirmatory.} This
subsection answers DQ1 only. It cannot modify, reinterpret, rescue or reopen the Pilot-B gate,
it defines no pass/fail criterion, and none of its content is confirmatory evidence for RQ1 or
RQ2.

The diagnostic was specified only after the Pilot-B result had been independently audited,
frozen and permanently closed. Its protocol was committed before any context-level
fixed-weight result was reconstructed. It used worlds 21001--21005 at the same reference
condition and the 200 frozen external TEST identities per world. The six structured vectors
were read from the immutable Pilot-B result without re-estimation; Equal was exactly
$(0.1,\ldots,0.1)$. D29 TEST decision matrices and oracle utilities were reconstructed only to
validate the archived aggregates. No target $Y$, model fit/prediction, SHAP, PI, RandomWeights,
bootstrap or protected seed was executed.

For each world, method and context the result stores the selected alternative; all six MOORA
scores and ranks; oracle utilities, ranks and selected alternative; Top-1 match; normalized
regret contribution; and agreement with the frozen MajorityWinner identity. Exactly
$5\times7\times200=7000$ records were independently audited. The audit reconstructed every
selection, anchor-based rank and regret directly from the stored full-precision vectors,
matched all 35 archived method--world aggregates within $10^{-12}$, and rebuilt all 105
pairwise selection-agreement summaries without executing the generator.

Table~\ref{supp:tab:fixedselectionmaps} reports each method's modal selected alternative,
modal share and number of distinct selected alternatives. Entries use the format
``alternative/\allowbreak share/\allowbreak $k$'', where $k=1$ denotes
context-invariant selection in that world.
MajorityWinner (MW) is a single frozen alternative by construction.

\begin{table}[t]
\centering
\caption{Audited selection maps for the seven observed fixed vectors. Each method entry is
modal alternative/modal share/number of distinct selected alternatives over 200 TEST contexts.}
\label{supp:tab:fixedselectionmaps}
\scriptsize
\setlength{\tabcolsep}{2.1pt}
\begin{tabular}{@{}rclllllll@{}}
\toprule
World & MW & Oracle & SHAP & PI & Ridge+ & CRITIC & Entropy & Equal\\
\midrule
21001 & A6 & A6/1.000/1 & A6/1.000/1 & A6/.980/2 & A6/1.000/1 & A6/1.000/1 & A6/1.000/1 & A6/1.000/1\\
21002 & A1 & A1/.700/3 & A1/.950/2 & A1/.630/2 & A1/.875/3 & A3/.800/3 & A1/1.000/1 & A1/1.000/1\\
21003 & A2 & A6/.720/4 & A2/.690/3 & A2/.530/3 & A6/.495/3 & A6/1.000/1 & A6/1.000/1 & A6/1.000/1\\
21004 & A2 & A2/.990/2 & A2/1.000/1 & A2/.775/4 & A2/.990/2 & A3/.750/3 & A3/.985/3 & A2/1.000/1\\
21005 & A2 & A6/.350/4 & A2/.790/3 & A3/.675/5 & A6/.375/4 & A6/.915/3 & A6/.800/3 & A2/1.000/1\\
\bottomrule
\end{tabular}
\end{table}

Equal was context invariant in all five worlds and exactly matched MajorityWinner across all
200 contexts in worlds 21001, 21002, 21004 and 21005. In world 21003, Equal instead selected A6
throughout while MajorityWinner was A2. The numbers of constant-selection worlds were
Oracle 1, SHAP 2, PI 0, Ridge+ 1, CRITIC 2, Entropy 3 and Equal 5; exact all-context
MajorityWinner agreement occurred in 1, 2, 0, 1, 1, 2 and 4 worlds, respectively.

The world-21003 ambiguity is now directly resolved: Equal, CRITIC and Entropy all selected A6
in every TEST context. Their different mean $\tau_b$ values therefore reflect differences below
the argmax, not different selected alternatives. More broadly, the heterogeneous constancy
counts show that fixed-vector MOORA was neither universally context invariant nor uniformly
context responsive over the seven observed vectors. The result supports coarse
observed-vector decision regions as a descriptive hypothesis, but cannot establish the causal
mechanism of the Pilot-B stop, characterize unobserved regions of the weight simplex, rank
methods or generalize beyond these five development worlds.

The new artefact makes paired context resampling technically possible, but the prospectively
frozen diagnostic explicitly excluded bootstrap and defined no inferential or pass/fail role.
No resampling was executed. Any future paired sensitivity analysis requires its own protocol,
frozen before inspecting bootstrap output, and cannot revise the closed Pilot-B decision.

\section{Detailed Weighting and Decision Operators}

The primary information partitions were fixed: Oracle, SHAP, PI, CRITIC and Entropy used WEIGHT;
Ridge+ used FIT; no method used external TEST for fitting or weight estimation. Oracle and SHAP
normalized mean absolute attributions. SHAP was declared undefined if total importance was at
most $10^{-12}$ and was never replaced by Equal weights.

CRITIC applied criterion-wise min--max scaling, population SD and Pearson conflict over
nonconstant columns. Entropy applied the same min--max scaling followed by Shannon
diversification with $0\log0=0$. Ridge+ standardized predictors and target with fold-local
population moments, solved non-negative L2-augmented least squares without an intercept, and
used the frozen grid
\begin{equation}
\tau\in\{10^{-4},10^{-3},10^{-2},10^{-1},1,10,100\}.
\label{supp:eq:ridgegrid}
\end{equation}
Five-fold context-grouped original-scale RMSE selected $\tau$, with the smallest value as tie
break.

Permutation importance kept the fitted XGBoost model fixed. For each criterion, it copied the
complete six-alternative block from source context $\pi(d)$ to destination context $d$ under 20
unique derangements. Their frozen generator was
\begin{equation}
\texttt{SeedSequence([replication, 84001, N])}.
\label{supp:eq:piseedsequence}
\end{equation}
The same
derangements were reused across criteria and factor levels for a given seed and $N$. Importance
was permuted minus baseline WEIGHT MSE; negative means were truncated before normalization and
repeat dispersion used sample SD. CRITIC, Entropy, Ridge+ and PI used Equal weights only under
their frozen zero-information rules, with fallback recorded.

Equal was $(0.1,\ldots,0.1)$. RandomWeights used 200
$\operatorname{Dirichlet}(1,\ldots,1)$ vectors per world from namespace 81001 and reused them
across conditions. MajorityWinner estimated the oracle modal winner on WEIGHT and applied that
constant choice to TEST. MajorityWinner was intentionally oracle-informed: it was the
context-free oracle-modal alternative estimated on WEIGHT and evaluated on TEST, and it was
optimal only for empirical oracle Top-1 frequency among constant choices. It is not claimed to
minimize regret or to constitute an operational transportation policy. It was neither a
deployable competitor nor a universally optimal choice, and it was not stakeholder-derived or
normative. The regret arm of the frozen gate uses this same
prospectively frozen Top-1-derived constant reference rather than a separately optimized
regret-minimizing constant policy; a future benchmark could prospectively define metric-matched
constant references. The frozen gate itself is unchanged by this observation. Direct-XGBoost was a parallel
predictive-decision reference rather than a weighting
method. MOORA used within-context vector normalization and a weighted sum; TOPSIS used the same
normalized benefit matrix as a conditional aggregation robustness operator.

\subsection{Frozen Tie Semantics and Archived Decision Margin}

Within each context, alternatives were sorted by descending raw score and ascending
\texttt{alternative\_id}. The highest unassigned score anchored a tie group; every remaining
alternative within absolute distance $10^{-12}$ of that anchor joined the group and received
the average of the occupied one-based rank positions. Ties were not chained through adjacent
near-equalities. Deterministic Top-1 used
\begin{equation}
\mathcal T_s=\{a:S_{\max,s}-S_{as}\le10^{-12}\}
\label{supp:eq:top1tieset}
\end{equation}
and selected the smallest alternative ID in $\mathcal T_s$. Relative tolerance was zero. These
semantics applied to oracle, MOORA, TOPSIS and Direct-XGBoost scores.

The conditional primary plan also defined the guarded oracle decision margin
\begin{equation}
M_s=\frac{U_{(1),s}^\star-U_{(2),s}^\star}
{\max_aU_{as}^\star-\min_aU_{as}^\star+10^{-12}}.
\label{supp:eq:margin}
\end{equation}
Margin-conditioned analyses were not executed because Pilot B stopped the primary study.

\section{Frozen Pilot-B Gate and Complete Decision}

Pilot B reused development worlds 21001--21005 at the frozen reference
\begin{equation}
(N,c,\rho,\lambda)=(250,0.30,0.4,0.5).
\label{supp:eq:pilotbreference}
\end{equation}
Each world used 200 external TEST contexts and 200 common
RandomWeights draws per world. It was a protocol-level adequacy screen, not independent
population validation. A structured method was eligible only if it was defined without Equal
fallback in all five worlds and Kendall's $\tau_b$ was defined in at least 190 TEST contexts per
world; at least 190 RandomWeights draws also had to be eligible. Cross-world advantages used
medians and required strict positivity in at least four worlds.

The five independently sufficient hard stops were: incomplete coverage; no eligible method
clearing median $\Delta\tau_b=0.05$ or $\Delta$mean-regret $=0.01$ versus RandomWeights; no
eligible method clearing median $\Delta$Top-1 $=0.02$ or $\Delta$mean-regret $=0.01$ versus
MajorityWinner; pooled oracle modal share at least $0.90$; or the frozen random-weight
insensitivity criterion, including its $0.90$ invariant fraction and joint width limits
$0.05$ ($\tau_b$), $0.02$ (Top-1) and $0.01$ (mean regret) in at least four worlds. Passage
required every stop to be false.

\begin{table}[t]
\centering
\caption{Complete frozen Pilot-B gate decision.}
\label{supp:tab:pilotbgates}
\small
\begin{tabularx}{\linewidth}{@{}lXc@{}}
\toprule
Gate & Audited result & Stop\\
\midrule
Coverage & Six structured methods eligible without fallback & No\\
Random equivalence & Four of six methods separated from RandomWeights & No\\
Modal-winner effective tie & Zero of six methods escaped MajorityWinner & \textbf{Yes}\\
Oracle-winner dominance & Pooled modal share $0.526<0.90$ & No\\
Weight insensitivity & Invariant fraction $0.004<0.90$; width-qualified worlds $0<4$ & No\\
\midrule
Overall OR rule & At least one frozen stop fired & \textbf{Yes}\\
\bottomrule
\end{tabularx}
\end{table}

The best strictly positive consistency count was two of five worlds for either Top-1 or mean
regret, and the best cross-world medians were zero. The independently reconstructed decision
therefore barred the protected factorial. This is a negative result about the adequacy of the
frozen comparison at one condition, not evidence that SHAP, weights or MCDM fail universally.

\section{Archived Stage-Wise Fidelity Ladder}

The conditional primary analysis contained nine representations:
\begin{enumerate}
\item \texttt{OracleUtility}, the full interacting oracle;
\item \texttt{OracleMainEffectReference}, with the interaction removed;
\item \texttt{OracleGlobalAttributionNonlinearQ};
\item \texttt{SHAPGlobalAttributionNonlinearQ};
\item \texttt{OracleGlobalAttributionLinearG};
\item \texttt{SHAPGlobalAttributionLinearG};
\item \texttt{OracleWeightMOORA};
\item \texttt{SHAPWeightMOORA};
\item \texttt{DirectXGBoost}, a parallel reference rather than a sequential stage.
\end{enumerate}
Eight paired contrasts isolated interaction removal, oracle global compression, SHAP estimation
on the nonlinear scale, nonlinear-to-linear transformation under each weight vector, MOORA
transformation under each vector and post-MOORA Oracle-versus-SHAP attribution. Kendall's
$\tau_b$ and normalized oracle regret were specified for every representation. The ladder was
diagnostic, not an additive error decomposition, and was not executed because Pilot B stopped
the primary study.

\section{Conditional Primary Factorial}

The four planned primary factors were
\begin{align}
N&\in\{25,50,100,250,1000\}, &
c&\in\{0.10,0.30,0.60\},\label{supp:eq:factors1}\\
\rho&\in\{0,0.4,0.8\}, &
\lambda&\in\{0,0.5,1\}.
\label{supp:eq:factors2}
\end{align}
The resulting $5\times3\times3\times3=135$ conditions were to be repeated over 30
prespecified primary benchmark-instance seeds, giving 4050 seed--condition replications. The
reference condition was
\begin{equation}
(N,c,\rho,\lambda)=(250,0.30,0.4,0.5).
\label{supp:eq:reference}
\end{equation}

\begin{table}[t]
\centering
\caption{Archived nested estimation partitions and fixed TEST size.}
\label{supp:tab:partitions}
\small
\begin{tabular}{@{}rrrr@{}}
\toprule
$N$ & FIT & WEIGHT & external TEST\\
\midrule
25 & 20 & 5 & 200\\
50 & 40 & 10 & 200\\
100 & 80 & 20 & 200\\
250 & 200 & 50 & 200\\
1000 & 800 & 200 & 200\\
\bottomrule
\end{tabular}
\end{table}

Technology-response parameters and common-random-number streams were to be sampled once per
seed and reused across factor levels. For each $(\rho,\lambda)$ cell, the oracle signal scale
$s_U$ was to be calculated on the complete 1000-context FIT+WEIGHT master pool so that changing
$N$ would not change the imposed noise-to-signal ratio. Realized signal-to-noise ratio was to
be logged for every cell. FIT, WEIGHT and external TEST partitions were to remain
context-disjoint, and no TEST row was to be used for model fitting or weight estimation.

\section{Archived Alpha-Dispersion Stress Analysis}

The conditional primary analysis retained the frozen heterogeneous coefficient vector. At the
reference condition only, three additional coefficient structures had been fixed: balanced,
dispersion-aligned and dispersion-anti-aligned. The mapping was obtained from an independent
200-world technology ensemble using
\begin{equation}
D_{j,w}^{attr}
=\frac{1}{M_w}\sum_i\left|q_j(g_{ijw})-\bar q_{j,w}\right|,
\qquad
\bar D_j^{attr}=\E_W[D_{j,W}^{attr}].
\label{supp:eq:qmad}
\end{equation}
The complete ordering agreed between 150 and 200 worlds:
\begin{equation}
C_8>C_{10}>C_5>C_2>C_3>C_4>C_1>C_7>C_6>C_9.
\label{supp:eq:d3order}
\end{equation}

\begin{table}[t]
\centering
\caption{Frozen 200-world mean attribution-scale dispersion.}
\label{supp:tab:d3dispersion}
\small
\begin{tabular}{@{}lc@{}}
\toprule
Criterion & Mean dispersion\\
\midrule
$C_8$ & 0.122873493\\
$C_{10}$ & 0.119403289\\
$C_5$ & 0.110114266\\
$C_2$ & 0.094820283\\
$C_3$ & 0.090482741\\
$C_4$ & 0.088452971\\
$C_1$ & 0.087135074\\
$C_7$ & 0.078706704\\
$C_6$ & 0.049591336\\
$C_9$ & 0.028266887\\
\bottomrule
\end{tabular}
\end{table}

Balanced coefficients were $\alpha_j=0.10$. The aligned and anti-aligned mappings were
\begin{align}
\boldsymbol\alpha^{aligned}:
&(C_8,C_{10},C_5,C_2,C_3,C_4,C_1,C_7,C_6,C_9)
\nonumber\\
&\mapsto(0.16,0.14,0.13,0.12,0.11,0.10,0.08,0.07,0.05,0.04),
\label{supp:eq:alphaaligned}\\
\boldsymbol\alpha^{anti}:
&(C_9,C_6,C_7,C_1,C_4,C_3,C_2,C_5,C_{10},C_8)
\nonumber\\
&\mapsto(0.16,0.14,0.13,0.12,0.11,0.10,0.08,0.07,0.05,0.04).
\label{supp:eq:alphaanti}
\end{align}
The frozen primary heterogeneous assignment had weak rank association with the dispersion
ordering ($\rho_S=0.1273$, $\tau_b=0.0222$), so it was not changed. Although the D3 mapping
itself was a completed development artefact, the planned primary stress comparisons were never
executed.

\section{Archived Bootstrap and Robustness Plan}

At the reference condition, $B=200$ context-level bootstrap replications were to resample FIT
and WEIGHT contexts separately, refit the model, recompute attributions and weights, and
evaluate decisions on the fixed original external TEST contexts. Confidence intervals were to
be reported for criterion weights, total-variation weight error, Kendall's $\tau_b$, Top-1
accuracy and normalized oracle regret. No cross-context probability that a particular ITS was
ranked first was to be interpreted as a global deployment probability.

A 50-oracle ensemble was to leave the frozen D29 technology-response kernel unchanged while
varying oracle coefficient vectors, nonlinearities and semantically admissible interaction
pairs. Relative outcome noise was to be recalibrated to each oracle variant's signal standard
deviation. This ensemble would not replace the response generator and therefore would not
reintroduce the historical v2.1 separability defect through oracle-only perturbations.

Two distribution-shift analyses were specified: holding out the upper 20\% of demand pressure
and holding out the lower 20\% of infrastructure readiness. Secondary robustness analyses were
to include criterion removal, $\pm10/20/30\%$ SHAP-weight perturbations,
leave-one-alternative-out analysis and MOORA--TOPSIS comparison.

\section{Archived Statistical Analysis Plan}

Factor effects were to be summarized with repeated-seed regression or mixed-effects models,
emphasizing effect sizes and uncertainty intervals rather than binary significance alone.
Method comparisons were paired because every method would share the same benchmark instance.
The 30 primary seeds were to be treated as benchmark-instance clusters; alternative--context
rows were not to be interpreted as independent superpopulation replications.

Method availability was to be reported explicitly. For each factorial cell, undefined
Oracle/SHAP weight rates and recorded Equal-fallback rates were to accompany conditional
decision-fidelity summaries. Pairwise contrasts were to use only benchmark instances on which
both methods were defined, with coverage and fallback rates reported alongside them. Undefined
Oracle or SHAP vectors were never to be silently replaced by Equal weights.

Planned outputs included predictive-performance and Direct-XGBoost references, weight-recovery
tables, decision-fidelity summaries, factorial effect summaries, alpha-alignment stress results,
bootstrap intervals, validity-boundary figures and the prespecified stage-wise fidelity ladder.
None of these primary outputs exists under v1.

\section{Reproducibility and Content-Integrity Record}
\label{supp:sec:integrity}

Table~\ref{supp:tab:digests} gives non-identifying truncated content digests. The selected B100
background exactly matched the frozen TreeSHAP identity set.

\begin{table}[t]
\centering
\caption{Truncated, non-identifying content digests of the frozen artefacts underlying the
reported numbers.}
\label{supp:tab:digests}
\small
\begin{tabular}{@{}ll@{}}
\toprule
Artifact & Truncated SHA-256\\
\midrule
Oracle B100 result & \texttt{7f9864bd...439c9fb0}\\
TreeSHAP background identities & \texttt{3184c675...1646eda1}\\
Pilot-B result & \texttt{05072857...9f787a88}\\
Fixed-weight selection-map result & \texttt{a1f293b6...d7e20f4e}\\
\bottomrule
\end{tabular}
\end{table}

The full anonymous provenance package records, in order, the
threshold freeze, evaluator/executor freeze, scientifically invariant direct-file loader repair,
audited Pilot-B result freeze, immutable closure reference, post-closure diagnostic protocol
freeze, implementation freeze and independently audited result freeze. The worktree was clean
after each closure point, and the final complete test run reported 836 passes with three
pre-existing external plotting warnings.
\ifanon
Exact repository identifiers are intentionally omitted here to protect anonymous review.
\else
Exact repository identifiers accompany the public reproducibility archive.
\fi
These facts document software and artefact integrity; they are not statistical evidence.

\section{Prospective Stop and Permanent Closure}

Pilot B was governed by five independently sufficient hard stops: coverage failure, failure to
separate RandomWeights, failure to escape MajorityWinner, pooled oracle-winner dominance and
weight insensitivity. Passage required every flag to be false. In the frozen result, only the
modal-winner effective-tie flag was true, but the overall OR rule therefore barred every
analysis in this supplement.

The immutable version-control, digest and test record appears in
Section~\ref{supp:sec:integrity}. Pilot B v1 is closed and must not be rerun. Any later study
requires a new specification and version reference, explicit disclosure of the v1 failure and
development-only design work; it cannot alter the v1 thresholds, result or closure reference.
The fixed-weight selection-map diagnostic in Section~\ref{supp:sec:fixedselection} is such a
separately frozen descriptive layer; it neither reopened Pilot B nor authorized any protected
analysis.

\end{document}
